%%%%%%%%%%%%%%%%%%%%%%%%%%%%%%%%%%%%%%%%%%%%%%%%%%%%%%%%%%%%%%%%%
\chapter{MAT1073 --- College Algebra for Scientists and Engineers}
\label{chap:unit-f995}
%%%%%%%%%%%%%%%%%%%%%%%%%%%%%%%%%%%%%%%%%%%%%%%%%%%%%%%%%%%%%%%%%

\ChapterWarning

%===============================================================================
\section{Algebraic Properties}
\label{sec:alge-a0fa}
%===============================================================================

Algebraic Properties

%-------------------------------------------------------------------------------
\subsection{Learning Objectives}
\label{ssec:lear-775b}
%-------------------------------------------------------------------------------

\begin{enumerate}
    \item Determine factors and prime factorizations
    \item Determine least common multiples
    \item Use variables and algebraic symbols
    \item Simplify expressions using the order of operations
    \item Evaluate an expression
    \item Identify and combine like terms
    \item Translate an English phrase to an algebraic expression
    \item Correctly identify the Additive \& Multiplicative identity
    \item Correctly identify the Additive \& Multiplicative inverse
    \item Correctly identify the Commutative property
    \item Correctly identify the Associative property
    \item Correctly identify the Distributive property
    \item Correctly explain the algebraic properties of numbers
    \item Correctly apply algebraic properties in procedural explanations of solving mathematical equations
\end{enumerate}

%-------------------------------------------------------------------------------
\subsection{Assessment Instrument}
\label{ssec:asse-370c}
%-------------------------------------------------------------------------------

Assess through short written exercises that ask students to name and apply each algebraic
property in a concrete simplification or equation-solving context. Problems should require
students to justify each step by citing the property used, reinforcing both procedural
fluency and conceptual understanding of identity, inverse, commutative, associative, and
distributive properties.

\textbf{Example problems.}

\begin{enumerate}
    \item Simplify the expression $3(x + 4) - 2x + 5$ and name every algebraic property used at each step.
    \item Given $a \cdot 0 = 0$ and $a + 0 = a$, identify which property each equation illustrates and state the general rule.
    \item Rewrite $7 + (3 + x)$ in two equivalent forms, one using the commutative property and one using the associative property, and explain the difference between them.
\end{enumerate}

%===============================================================================
\section{Fractions}
\label{sec:frac-8099}
%===============================================================================

Fractions

%-------------------------------------------------------------------------------
\subsection{Learning Objectives}
\label{ssec:lear-10a5}
%-------------------------------------------------------------------------------

\begin{enumerate}
    \item Determine equivalent fractions
    \item Determine common denominators to add and subtract fractions
    \item Add fractions
    \item Subtract fractions
    \item Multiply fractions
    \item Dividing fractions
    \item Work with proper and improper fractions
    \item Simplify fractions to lowest terms
\end{enumerate}

%-------------------------------------------------------------------------------
\subsection{Assessment Instrument}
\label{ssec:asse-06cc}
%-------------------------------------------------------------------------------

Assess through computational problems that progress from finding equivalent fractions and
common denominators to performing all four arithmetic operations on proper and improper
fractions. Include at least one problem requiring simplification to lowest terms to
confirm that students can identify the GCF of the numerator and denominator.

\textbf{Example problems.}

\begin{enumerate}
    \item Compute $\dfrac{3}{8} + \dfrac{5}{12}$ and express the result in lowest terms.
    \item Evaluate $\dfrac{7}{4} \div \dfrac{3}{2}$ and convert the answer to a mixed number.
    \item Simplify $\dfrac{18}{24}$ to lowest terms, showing the GCF used.
\end{enumerate}

%===============================================================================
\section{Factoring Polynomials}
\label{sec:fact-754c}
%===============================================================================
%-------------------------------------------------------------------------------
\subsection{Useful Formulas}
\label{sss:usef-1334}
%-------------------------------------------------------------------------------

Computing factors of polynomials requires knowledge of different formulas and some
experience to find out which formula to apply. Below, we provide some important formulas:

\begin{itemize}
    \item $x^2 - y^2 = (x + y)(x - y)$
    \item $x^2 + 2xy + y^2 = (x + y)^2$
    \item $x^2 - 2xy + y^2 = (x - y)^2$
    \item $x^3 - y^3 = (x - y)(x^2 + xy + y^2)$
    \item $x^3 + y^3 = (x + y)(x^2 - xy + y^2)$
    \item $a^3 + b^3 + c^3 - 3abc = (a + b + c)(a^2 + b^2 + c^2 - ab - ac - bc)$
\end{itemize}

%-------------------------------------------------------------------------------
\subsection{Methods}
\label{sss:meth-a9ac}
%-------------------------------------------------------------------------------

Given the expression $x^2 + 3x + 2$, one may ask ``What are the values of $x$ that make this expression 0?'' If we factor, we obtain:
\[
x^2 + 3x + 2 = (x + 2)(x + 1)
\]
If $x = -1$ or $x = -2$, then one of the factors on the right becomes zero. Therefore, the whole expression must be zero. By factoring, we have discovered the values of $x$ that render the expression zero. These values are termed ``roots.''

In general, given a quadratic polynomial $px^2 + qx + r$ that factors as:
\[
px^2 + qx + r = (ax + c)(bx + d)
\]
Then, the roots of the polynomial are:
\[
x = -\frac{c}{a} \text{ and } x = -\frac{d}{b}
\]
A special case to be aware of is the difference of two squares, $a^2 - b^2$. In this case, we can always factor as:
\[
a^2 - b^2 = (a + b)(a - b)
\]
For example, consider $4x^2 - 9$. On initial inspection, we see that both $4x^2$ and $9$ are squares of $2x$ and $3$, respectively. Applying the previous rule, we have:
\[
4x^2 - 9 = (2x + 3)(2x - 3)
\]

%-------------------------------------------------------------------------------
\subsection{The AC Method}
\label{sss:thea-afe9}
%-------------------------------------------------------------------------------

The AC method simplifies the process of factoring. Suppose a quadratic polynomial has the
form:
\[
ax^2 + bx + c
\]
If there are numbers $r$ and $q$ that satisfy:
\[
r \cdot q = ac
\]

\[
r + q = b
\]
Then, we can rewrite the polynomial as:
\[
ax^2 + bx + c = ax^2 + rx + qx + c
\]
And factor out a common term from $(ax^2 + rx)$ and $(qx + c)$ to factor the polynomial.

For example, take the polynomial $2x^2 + 7x - 15$:
\[
ac = 2 \cdot (-15) = -30
\]

\[
10 \cdot (-3) = -30 \text{ and } 10 + (-3) = 7
\]
So we can rewrite the polynomial as:
\[
2x^2 + 10x - 3x - 15
\]
From here, we can factor $2x$ from the first two terms and $-3$ from the last two, which gives us:
\[
2x(x + 5) - 3(x + 5) = (2x - 3)(x + 5)
\]
Thus, the factored form of $2x^2 + 7x - 15$ is:
\[
(2x - 3)(x + 5)
\]

%-------------------------------------------------------------------------------
\subsection{The Quadratic Formula}
\label{sss:theq-f874}
%-------------------------------------------------------------------------------

Given any quadratic equation $ax^2 + bx + c = 0$, where $a \neq 0$, all solutions are given by the quadratic formula:
\[
x = \frac{-b \pm \sqrt{b^2 - 4ac}}{2a}
\]
Note that the value of $b^2 - 4ac$ will affect the number of real solutions:

\begin{itemize}
    \item If $b^2 - 4ac > 0$, there are two real solutions.
    \item If $b^2 - 4ac = 0$, there is one real solution.
    \item If $b^2 - 4ac < 0$, there are no real solutions.
\end{itemize}

%-------------------------------------------------------------------------------
\subsection{Remainder and Factor Theorem}
\label{sss:rema-cd36}
%-------------------------------------------------------------------------------

The polynomial division algorithm is as follows: Suppose $d(x)$ and $p(x)$ are nonzero polynomials where the degree of $p(x)$ is greater than or equal to the degree of $d(x)$. Then there exist two unique polynomials, $q(x)$ and $r(x)$, such that:
\[
p(x) = d(x)q(x) + r(x)
\]
Where either $r(x) = 0$ or the degree of $r(x)$ is strictly less than the degree of $d(x)$.

%-------------------------------------------------------------------------------
\subsection{Remainder Theorem}
\label{sss:rema-0f44}
%-------------------------------------------------------------------------------

Suppose $p(x)$ is a polynomial of degree at least 1 and $c$ is a real number. When $p(x)$ is divided by $(x - c)$, the remainder is $p(c)$.

\textbf{Proof:} By the division algorithm:
\[
p(x) = (x - c)q(x) + r
\]
Where $r$ must be a constant since $d(x) = x - c$ has a degree of 1. Setting $x = c$, we get:
\[
p(c) = (c - c)q(x) + r = r
\]
Thus, the remainder $r = p(c)$.

%-------------------------------------------------------------------------------
\subsection{Factor Theorem}
\label{sss:fact-8420}
%-------------------------------------------------------------------------------

Suppose $p(x)$ is a nonzero polynomial. The real number $c$ is a zero of $p(x)$ if and only if $(x - c)$ is a factor of $p(x)$.

By the division algorithm, $x - c$ is a factor of $p(x)$ if and only if $r = 0$. Since $p(c) = r$ when $p(x)$ is divided by $x - c$, $x - c$ is a factor of $p(x)$ if and only if $p(c) = 0$; that is, if $c$ is a zero of $p(x)$.

%-------------------------------------------------------------------------------
\subsection{Factor Theorem Example Problem}
\label{sss:fact-0b0b}
%-------------------------------------------------------------------------------

Determine if $x + 2$ is a factor of $2x^2 + 3x - 2$.

Since $c$ is positive instead of negative, use this basic identity:
\[
x + 2 = x - (-2)
\]
Now use the factor theorem:
\[
2(-2)^2 + 3(-2) - 2 = 8 - 6 - 2 = 0
\]
Since the result is 0, $(x + 2)$ is a factor of $2x^2 + 3x - 2$.

Thus, it can be re-stated as:
\[
2x^2 + 3x - 2 = (x + 2)(ax + b)
\]
Expanding the right-hand side:
\[
2x^2 + 3x - 2 = ax^2 + x(2a + b) + 2b
\]
Equating like terms, we get:
\[
2 = a
\]

\[
2a + b = 3
\]

\[
2b = -2
\]
Solving these, $a = 2$ and $b = -1$. Thus:
\[
2x^2 + 3x - 2 = (x + 2)(2x - 1)
\]

%-------------------------------------------------------------------------------
\subsection{Example Problems}
\label{sss:exam-b742}
%-------------------------------------------------------------------------------

\textbf{EXAMPLE 1:} Find all the roots of $4x^2 + 7x - 2$.

Finding the roots is equivalent to solving $4x^2 + 7x - 2 = 0$. Applying the quadratic formula with $a = 4$, $b = 7$, $c = -2$, we have:
\[
x = \frac{-7 \pm \sqrt{7^2 - 4 \cdot 4 \cdot (-2)}}{2 \cdot 4}
\]

\[
x = \frac{-7 \pm \sqrt{49 + 32}}{8}
\]

\[
x = \frac{-7 \pm \sqrt{81}}{8}
\]

\[
x = \frac{-7 \pm 9}{8}
\]

\[
x = \frac{2}{8} \text{ or } x = \frac{-16}{8}
\]

\[
x = \frac{1}{4} \text{ or } x = -2
\]
\textbf{EXAMPLE 2:} Factor the polynomial $4x^2 + 7x - 2$.

We know the roots are $x = \frac{1}{4}$ and $x = -2$. Therefore, the polynomial factors as:
\[
4x^2 + 7x - 2 = 4(x - \frac{1}{4})(x + 2)
\]
Or:
\[
4x^2 + 7x - 2 = (4x - 1)(x + 2)
\]
\textbf{EXAMPLE 3:} Factor $4x^2 + 12x + 9$ using the perfect square trinomial method.

We observe that $4x^2 + 12x + 9$ can be rewritten as:
\[
(2x + 3)^2
\]
To verify, expand $(2x + 3)^2$:
\[
(2x + 3)^2 = 4x^2 + 12x + 9
\]

%-------------------------------------------------------------------------------
\subsection{Resources}
\label{sss:reso-55b5}
%-------------------------------------------------------------------------------

\begin{itemize}
    \item \href{https://tutorial.math.lamar.edu/classes/alg/factoring.aspx}{Factoring Polynomials with Example Problems, Paul''s Online Notes}
    \item \href{https://courses.lumenlearning.com/suny-beginalgebra/chapter/5-2-1-factor-trinomials/}{Factoring Polynomials, Lumen Learning}
    \item \href{https://www.youtube.com/watch?v=shHS-ERFVrc}{Difference of Squares. Produced by TA Catherine Sporer, UTSA}
    \item \href{http://www.algebralab.org/lessons/lesson.aspx?file=algebra_factoring.xml}{6 Methods of Factoring with Practice Problems, AlgebraLAB}
    \item \href{https://txwes.edu/media/twu/content-assets/images/academics/academic-success-center/A-C-Method.pdf}{AC Factoring Method, Texas Wesleyan University}
\end{itemize}

%-------------------------------------------------------------------------------
\subsection{Licensing}
\label{sss:lice-7352}
%-------------------------------------------------------------------------------

Content obtained and/or adapted from:

\begin{itemize}
    \item \href{https://en.wikibooks.org/wiki/Calculus/Algebra}{Algebra, Wikibooks: Calculus} under a CC BY-SA license
    \item \href{https://en.wikibooks.org/wiki/Algebra/Factoring_Polynomials}{Factoring Polynomials, Wikibooks: Algebra} under a CC BY-SA license
\end{itemize}

%-------------------------------------------------------------------------------
\subsection{Learning Objectives}
\label{ssec:lear-30ae}
%-------------------------------------------------------------------------------

\begin{enumerate}
    \item Identify factored vs non-factored forms of a polynomial expression
    \item Differentiate between factors and terms of a polynomial expression
    \item Factor out greatest common factors from a polynomial expression
    \item Factor difference of squares into two binomial terms
    \item Factor trinomials into two binomial terms
    \item Multiply and/or distribute to check the factors of a polynomial expression
\end{enumerate}

%-------------------------------------------------------------------------------
\subsection{Assessment Instrument}
\label{ssec:asse-d43e}
%-------------------------------------------------------------------------------

Assess through multi-step problems requiring students to choose and apply an appropriate
factoring method: GCF extraction, difference of squares, perfect square trinomials, the AC
method, or the quadratic formula. Problems should also require use of the Remainder and
Factor Theorems to verify results.

\textbf{Example problems.}

\begin{enumerate}
    \item Factor $6x^2 + 11x - 10$ completely using the AC method and verify the result by expanding.
    \item Use the Factor Theorem to determine whether $(x - 3)$ is a factor of $2x^3 - 5x^2 - 4x + 3$, showing all work.
    \item Find all real roots of $3x^2 - 7x + 2 = 0$ using the quadratic formula and interpret the value of the discriminant.
\end{enumerate}

%===============================================================================
\section{Linear Equations}
\label{sec:line-d9b6}
%===============================================================================

%-------------------------------------------------------------------------------
\subsection{What are Linear Equations?}
\label{sss:what-22ca}
%-------------------------------------------------------------------------------

In the functions section we talk about how a function is like a box that takes an
independent input value and uses a rule defined mathematically to create a unique output
value. The value for the output is dependent on the value that is put in the box. We call
the values that are going into the box the independent values or the domain. We call the
values coming out of the box the dependent values or the range.

Unless we specify differently, on Cartesian graphs the domain is the real numbers. In the
Cartesian Plane section we saw how running different values through a function to identify
the points on the Cartesian plane by picking the first (x) value of the point the domain
and the second (y) value from the range. To restate this: by convention the two variables
for a function on the Cartesian Plane are x for the domain, the independent variable, and
y for the range, the dependent variable. The variable y is the same as writing f(x).
Mathematicians recognize this equivalence but generally prefer to write y because it is
shorter. Because a function definition has an input and an output it must also contain an
equal sign. The section in this book solving equations showed the various operations we
can perform on both sides of an equal sign and still maintain the notion of equivalence.
In this section we plug different values into the independent variable and solve to find
the associated dependent variable. For instance if we start with the equation:

\[
y = x
\]

\[
\text{We can add a } -x \text{ to both sides to get the equation}
\]

\[
y - x = x - x
\]

\[
\text{which we then simplify to}
\]

\[
y - x = 0
\]

\[
\text{Or we can add a } -y \text{ to both sides to get the equation}
\]

\[
y - y = x - y
\]

\[
\text{which we then simplify to}
\]

\[
0 = x - y
\]

\[
\text{Since } y - x = 0  \equiv 0 = x - y \text{ then:}
\]

\[
y - x = 0 = x - y
\]

\[
\text{Using the transitive property}
\]

\[
y - x  = x - y
\]

\[
\text{adding } x + y \text{ to both sides gives us}
\]

\[
(y - x )+( y + x)  = (x - y) + (y + x)
\]

\[
\text{using the associative property we change this to}
\]

\[
(y - y )+( x + x)  = (x - x) + (y + y)
\]

\[
\text{which simplifies to}
\]

\[
2x  = 2y
\]

\[
\text{And divide both sides by } 2
\]

\[
\frac{2x}{2}  = \frac{2y}{2}
\]

\[
\text{To simply reverse the order and show}
\]

\[
x = y
\]

We have not really proved anything mathematically above, but these operations allow us to
manipulate equations to get the dependent variable by itself on one side of the equals
sign. Then we can plug numbers into the independent variable to discover the function
values for those numbers. Then we can draw these values as points on the Cartesian plane
and get a feel for what the function would look like if we could see all the points
defined by the function at once.

%-------------------------------------------------------------------------------
\subsection{Horizontal Linear Equations}
\label{sss:hori-b67b}
%-------------------------------------------------------------------------------

Equations of the form $y = C_2$ are linear functions of the general form $y = mx + b$ where slope $m = 0$ and the constant $C_2$ is the y-intercept $b$ (in the general form). The graph of this zero-slope function is a straight horizontal line, intercepting the y-axis at $C_2$, including zero and extends infinitely in the positive and negative directions for all $R$ values of x.

% [Image: graph1]

The domain for such functions is $R$ covering all real numbers (unless otherwise specified), but the range is just the set $\{ c \}$. The equation $y = 0$ is the x-axis.

%-------------------------------------------------------------------------------
\subsection{Vertical Linear Equations}
\label{sss:vert-19e2}
%-------------------------------------------------------------------------------

Equation $x = C_1$, x is one single value $C_1$ and y, being unrestricted, is every $R$ number. The graph of $x = C_1$ is a straight vertical line where $x = C_1$, covering all positive, negative and zero values of y.

% [Image: x = c, where c = 3]

Its domain is set $\{ C_1 \}$ and range is set $R$ (unless otherwise specified). $x = C_1$ is technically not a function (there is more than one value of y for each value of x), but it''s a relation. Vertical lines have no slope ($m =$ divide by zero, undefined, plus and minus infinity). These are the only types of linear equations of the general form shown previously which are not linear functions. The equation $x = 0$ is exactly the y-axis. Lines with steepness approaching vertical have very large-magnitude slopes but are still functions.

%-------------------------------------------------------------------------------
\subsection{General Linear Equations}
\label{sss:gene-4191}
%-------------------------------------------------------------------------------

We are going to start by looking at simple functions called linear equations. When none of
the instances of x and y in the algebraic expression defining the function rule have
exponents then all the instances of x and y can be combined into just two occurrences. The
graph of the expression can be represented as a straight line. The equation that expresses
the function is considered a linear equation with two variables. The following equation is
a simple example of such a linear equation:

\[
y - x = 2
\]

Since y is the dependent variable it is standing in for the function. We can re-write the expression as $f(x) - x = 2$. If we add an x to both sides the equality property holds and we get the expression $f(x) - x + x = x + 2$. Simplifying we get $f(x) = x + 2$. In the following table we''ll pick 3 values for x, and then calculate the dependent (y) values from f(x).
\[
\begin{array}{|c|c|c|}
\hline
\text{x value} & \text{y value (abscissa)} & \text{Coordinates (x,y)} \\ 
\hline
-1 & 1 & (-1, 1) \\ 
\hline
0 & 2 & (0, 2) \\ 
\hline
1 & 3 & (1, 3) \\ 
\hline
\end{array}
\]
where x and y are variables to be plotted in a two-dimensional Cartesian coordinate graph
as shown here:

% [Image: 3ptsyeqxpl2.png]

% [Image: Y equals x plus 2.PNG]

This function is equivalent to the previous example of a linear equation, $y - x = 2$. The arrows at each end of the line indicate that the line extends infinitely in both directions. All linear functions of a single input variable have or can be algebraically arranged to have the general form:

\[
y = f(x) = mx + b
\]

where x and y are variables, f(x) is the function of x, m is a constant called the slope of the line, and b is a constant which is the ordinate of the y-intercept (i.e. the value of y where the function line crosses the y-axis). The slope indicates the steepness of the line. In the previous example where $y = x + 2$, the slope $m = 1$ and the y-intercept ordinate $b = 2$. The $y = mx + b$ form of a linear function is called the slope-intercept form.

Unless a domain for x is otherwise stated, the domain for linear functions will be assumed to be all real numbers and so the lines in graphs of all linear functions extend infinitely in both directions. Also in linear functions with all real number domains, the range of a linear function will cover the entire set of real numbers for y, unless the slope $m = 0$ and the function equals a constant. In such cases, the range is simply the constant.

Conversely, if a function has the general form $y = mx + b$ or if it can be arranged to have that form, the function is linear. A linear equation with two variables has or can be algebraically rearranged to have the general form:

\[
A x + B y = C
\]

where x and y represent the linear variables, and the letters A, B, and C can represent
any real constants, either positive or negative. Conversely, an equation with two
variables x and y having that general form, or being able to be arranged in that form,
would be linear as long as A and B are not both equal to 0. In the preceding equation,
capital letters are to avoid confusion with other constants in this chapter and for
consistency with Reference 1.

If one divides the preceding equation by B (when B is not 0) and solves for y, the
following form can be obtained:

\[
y = \frac{-A}{B} x + \frac{C}{B}
\]

If one equates $-A/B$ to the slope m and $C/B$ to the y-intercept ordinate b, it can be seen that the general form for a linear equation and the slope-intercept form for a linear function are practically interconvertible except for the fact that, in a linear function, the B constant in the linear equation form cannot equal 0.

%-------------------------------------------------------------------------------
\subsection{X and Y axis intercepts}
\label{sss:xand-5717}
%-------------------------------------------------------------------------------

An axis intercept point is a point where the graph of a function, relation, or equation
intersects the X or Y axes. This section is about finding out how a particular set of
functions: linear functions cross the axes.

We know that the domains of most lines are infinite because they are defined at every value of $X$. The exception is lines that are defined as $X = c$ where $c$ is a number we choose when we write the function. By definition, this line is only defined for one value of $X$. Since the domain maps onto more than one value for the range, this is actually a relationship and not a function. We''ve seen the graph of this relationship is a vertical line that passes through the point $(c,Y)$.

Lines with the equation $X = c$ intersect the X axis once and the Y axis 0 or infinite times.

We can also restrict the range of a function by simply writing $Y = c$, where $c$ is again any number we choose. The graph of this line is a horizontal line that passes through the point $(X, c)$. When $Y = 0$, this line is the same as the X axis. When $c \neq 0$, then the line can never intercept the X axis.

Lines with the equation $Y = c$ intersect the Y axis once and the X axis 0 or infinite times.

When looking at Cartesian graphs and linear equations, we run into a mathematical axiom: ``Two points determine a line.'' We will see how this axiom affects the slope-intercept definition of a line $y = f(x) = mx + b$ in the next section. When two lines intersect, they intersect at a point. If a line is not horizontal or perpendicular, it will have to intersect the X and Y axes once, but only once.

In this page, we are going to accept the statement ``At most one line can be drawn through
any point not on a given line parallel to the given line in a plane.''

We''ve seen that for the equation $Y = mX + b$, the Y intercept will always be at $b$ because that is where $X = 0$.

Using Algebra, we can subtract $b$ from both sides: $Y - b = mX$

and multiply by $\frac{1}{m}$

$\frac{Y}{m} - \frac{b}{m} = X$

we can see that the X intercept is going to be $-\frac{b}{m}$.

An axis intercept may simply refer to the number value on the axis where the intersection occurs. For brevity, we may say the line has an X intercept of 1 and a Y intercept of 2. After graphing just a few lines, you will be able to tell this line points down and runs through quadrants II, I, and IV. With a little more practice, you will be able to know that the equation for the line is $Y = -2x + 2$. We will see that by specifying the two points, we are actually implying the slope of the line. There is an exception to this rule. If we say a line crosses the axes at 0, we know that the line will pass through 2 quadrants instead of 3, but we won''t know which quadrants or how steep the line is. When we look at slope in the next section, we will see why the equations above specify a point and a slope.

When you are trying to graph a linear equation, finding the axes intercepts is often the easiest way to go about doing it. To find the x-intercept, set $y = 0$ and solve for $x$. To find the y-intercept, set $x = 0$ and solve for $y$. For most examples, the intercepts are different points, and a line can be drawn through the two intercepts. If both intercepts are $(0, 0)$, then another point must be determined to graph the line. If the equation is in the form $x = c$ or $y = c$, the horizontal or vertical lines are very simple to plot.

%-------------------------------------------------------------------------------
\subsection{Slope}
\label{sss:slop-0bbe}
%-------------------------------------------------------------------------------

%-------------------------------------------------------------------------------
\subsection{Algebra/Slope}
\label{sss:alge-547c}
%-------------------------------------------------------------------------------

Slope is the change in the vertical distance of a line on a coordinate plane over the change in horizontal difference. In other words, it is the “rise” over the “run” or the steepness of a line. Slope is usually represented by the symbol $m$ like in the equation $y = mx + b$, where $m$, the coefficient of $x$, represents the slope of the line.

Slope is computed by measuring the change in vertical distance divided by the change in
horizontal difference, i.e.:
\[
m = \frac{\Delta y}{\Delta x} = \frac{\text{rise}}{\text{run}}
\]
The Greek uppercase letter $\Delta$ represents change, in this case, change in the $y$-coordinates divided by the change in $x$-coordinates.

%-------------------------------------------------------------------------------
\subsection{Positive Slope/ Negative Slope}
\label{sss:posi-9cc1}
%-------------------------------------------------------------------------------

If a line goes up from left to right, then the slope has to be positive. For example, a slope of $\frac{3}{4}$ would have a “rise” of 3, or go up 3; and a “run” of 4, or go right 4. Both numbers in the slope are either negative or positive in order to have a positive slope.

If a line goes down from left to right, then the slope has to be negative. For example, a slope of $-\frac{3}{4}$ would have a “rise” of -3, or go down 3; and a “run” of 4, or go right 4. Only one number in the slope can be negative for a line to have a negative slope.

%-------------------------------------------------------------------------------
\subsection{Other Types of Slope}
\label{sss:othe-c807}
%-------------------------------------------------------------------------------

There are two special circumstances: no slope and slope of zero. A horizontal line has a
slope of 0, and a vertical line has an undefined slope.

Horizontal lines have the form: $y = a$ ; where $a$ is a constant, i.e. $a \in \mathbb{R}$.

Vertical lines have the form: $x = a$ ; where $a$ is a constant, i.e. $a \in \mathbb{R}$.

%-------------------------------------------------------------------------------
\subsection{Determining Slope}
\label{sss:dete-8960}
%-------------------------------------------------------------------------------

To determine the slope, you need some information. This can include two (or more)
coordinates, a parallel slope and a coordinate, a perpendicular slope and a coordinate, or
the y-intercept and slope.

For completely horizontal lines, the difference in $y$ coordinates between any two points is 0, so the slope $m = 0$, indicating no steepness in the line at all. If the line extends between right-upper $(+,+)$ and left-lower $(-,-)$ directions, then the slope is positive. As the slope increases, the line becomes steeper until the line is almost vertical when the slope is very large. When the slope $m = 1$, the line is diagonal with an angle halfway between the $x$ and $y$ axes. If the line extends between left-upper $(-,+)$ and right-lower $(+,-)$ directions, then the slope is negative. As the slope changes from 0 to very negative numbers, the steepness in the opposite direction increases. Compare the slope $(m)$ values in the following graph of functions $y = 1$ (where $m = 0$), $y = \frac{1}{2}x + 1$, $y = x + 1$, $y = 2x$, $y = -\frac{1}{2}x + 1$, $y = -x + 1$, and $y = -2x + 1$. For all two-variable linear equations that can be converted to linear functions, the same calculation applies to slopes for those lines.

\textbf{Various Linear Slopes}

% [Image: slopes]

For the most part, finding the slope when given information is a simple matter. Simply take the slope equation $y = mx + b$ and replace the variable with whatever information you know, and solve.

%-------------------------------------------------------------------------------
\subsection{Two Coordinates}
\label{sss:twoc-24d3}
%-------------------------------------------------------------------------------

To find the slope with two coordinates, you must first find the slope. Use the standard equation $\frac{y_2 - y_1}{x_2 - x_1}$. Put that into the equation as $m$, and replace $x$ and $y$ with $x$ and $y$ from one of the coordinates. Solve for $b$. Put that into the equation, and you''re done.

Example: $(1, 4)$ $(4, 8)$
\[
m = \frac{y_2 - y_1}{x_2 - x_1} = \frac{4 - 1}{4 - 8} = \frac{3}{-4}
\]
Plug that right in.
\[
y = \frac{3}{-4}x + b
\]

\[
4 = \frac{3}{-4}(1) + b
\]

\[
4 - \frac{3}{-4} = b
\]

\[
\frac{-19}{4} = b
\]
Put that in the equation, and you''re done.
\[
y = \frac{3}{-4}x + \frac{-19}{4}
\]

%-------------------------------------------------------------------------------
\subsection{Parallel Lines}
\label{sss:para-7283}
%-------------------------------------------------------------------------------

If you have to find the slope of a line (Let''s say $AB$) which is parallel to line (Let''s say $XY$), then using the coordinates of line $XY$, you can find the coordinates of the slope of line $AB$. As, the slope of a line parallel to another line is equal, i.e. Slope of $AB$ = Slope of $XY$.

\textbf{Example:}

$AB$ and $XY$ are PARALLEL Lines. Find the slope of $AB$.

Let line $XY$ have the coordinates:

$X(2, 4) = (x_1, y_1)$ and $Y(3, 6) = (x_2, y_2)$

\textbf{Slope of $XY$:}
\[
m = \frac{y_2 - y_1}{x_2 - x_1} = \frac{6 - 4}{3 - 2} = \frac{2}{1} \Rightarrow m = 2
\]
Slope of $AB = $ Slope of $XY$, so Slope of $AB = 2$: 
%-------------------------------------------------------------------------------
\subsection{Resources:}
\label{sss:reso-55b5}
%-------------------------------------------------------------------------------

\begin{itemize}
    \item \href{https://courses.lumenlearning.com/beginalgebra/chapter/introduction-to-solving-linear-equations/}{Solving Linear Equations}, Lumen Learning
    \item \href{https://openstax.org/books/intermediate-algebra-2e/pages/2-1-use-a-general-strategy-to-solve-linear-equations}{General Strategy to Solving Linear Equations}, OpenStax
    \item \href{https://openstax.org/books/intermediate-algebra-2e/pages/2-3-solve-a-formula-for-a-specific-variable}{Solving Linear Equations for a Specific Variable}, OpenStax
    \item \href{https://www.youtube.com/watch?v=bsGyJ0GeRy8}{Solving Linear Equations: Example 1.} Produced by TA Catherine Sporer, UTSA
    \item \href{https://www.youtube.com/watch?v=nGZcfHQEdqo}{Solving Linear Equations: Example 2.} Produced by TA Catherine Sporer, UTSA
\end{itemize}

%-------------------------------------------------------------------------------
\subsection{Licensing}
\label{sss:lice-7352}
%-------------------------------------------------------------------------------

Content obtained and/or adapted from:

\begin{itemize}
    \item \href{https://en.wikibooks.org/wiki/Algebra/Linear_Equations_and_Functions}{Linear Equations and Functions, Wikibooks: Algebra} under a CC BY-SA license
    \item \href{https://en.wikibooks.org/wiki/Algebra/Intercepts}{Intercepts, Wikibooks: Algebra} under a CC BY-SA license
    \item \href{https://en.wikibooks.org/wiki/Algebra/Slope}{Slope, Wikibooks: Algebra} under a CC BY-SA license
\end{itemize}

%-------------------------------------------------------------------------------
\subsection{Learning Objectives}
\label{ssec:lear-fc06}
%-------------------------------------------------------------------------------

\begin{enumerate}
    \item \textbf{1. Understanding and Application of Linear Equations:}
    \item Students will be able to define and explain linear equations, including their components such as slope and y-intercept.
    \item Students will apply their knowledge to solve linear equations and interpret the results in real-world contexts. \textbf{2. Graphing Linear Equations:}
    \item Students will be able to graph linear equations on the Cartesian plane, identifying key features such as the slope, y-intercept, and x-intercept.
    \item Students will interpret and analyze graphs to understand the relationship between the variables. \textbf{3. Manipulating and Transforming Linear Equations:}
    \item Students will be able to perform algebraic manipulations on linear equations, including addition, subtraction, multiplication, and division of both sides of the equation.
    \item Students will demonstrate the ability to transform equations to isolate variables and solve for specific values.
\end{enumerate}

%-------------------------------------------------------------------------------
\subsection{Assessment Instrument}
\label{ssec:asse-9733}
%-------------------------------------------------------------------------------

Assess through problems that require students to identify, graph, and solve linear
equations in one and two variables, including horizontal, vertical, and general-form
equations. Problems should connect the algebraic solution to the geometric interpretation
of slope and intercepts on the Cartesian plane.

\textbf{Example problems.}

\begin{enumerate}
    \item Write the slope-intercept form of the line passing through $(2, -1)$ with slope $m = 3$, then graph the line and label the $y$-intercept.
    \item Solve $4(x - 2) + 3 = 2x + 7$, identifying the algebraic property used at each step.
    \item Determine whether $y = 5$ and $x = -3$ represent functions and state the domain and range of each.
\end{enumerate}

%===============================================================================
\section{Systems of Equations in Two Variables}
\label{sec:syst-660c}
%===============================================================================

%-------------------------------------------------------------------------------
\subsection{Introduction}
\label{sss:intr-8800}
%-------------------------------------------------------------------------------

A system of equations consists of two or more equations made up of two or more variables such that all equations in the system are considered simultaneously. To find the unique solution to a system of equations, we must find a numerical value for each variable in the system that will satisfy all equations in the system at the same time. Some systems may not have a solution (for example, two distinct lines that are parallel to one another) and others may have an infinite number of solutions (for example, $y = 0$ and $y = \sin(x)$, since $\sin(x) = 0$ for an infinite number of $x$-values). In order for a linear system to have a unique solution, there must be at least as many equations as there are variables. Even so, this does not guarantee a unique solution.

Systems of equations can be solved in a number of ways. We can use elimination and/or
substitution to solve for points of intersection shared by all equations in the system. We
can also sometimes graph the equations to find the points shared by all equations, given
that our system is graphable.

Here are some examples of systems of equations with two variables:

\begin{itemize}
    \item \textbf{One solution}: $y = 2x$ and $y = x + 5$. By use of substitution, elimination, or graphing, we can find that the solution to this system is $(5, 10)$, as this is the only point shared by both equations in the system.
\end{itemize}

\begin{itemize}
    \item \textbf{Multiple solutions}: $y = x^2$ and $y = 4$. Since $x^2 = 4$ when $x = 2$ or $-2$, we know there are two solutions to this system: $(2, 4)$ and $(-2, 4)$. This is also clear when we graph these two equations, as $y = 4$ intersects the parabola $y = x^2$ when $x = -2$, and when $x = 2$.
\end{itemize}

\begin{itemize}
    \item \textbf{No solutions}: $y = x^2$ and $y = -3$. Since $x^2$ is positive for all values of $x$, it cannot equal $-3$. When graphed, we can see that these two equations never intersect. Thus there are no solutions to this system of equations.
\end{itemize}

\begin{itemize}
    \item \textbf{Infinite solutions}: $y = \sin(x)$ and $y = 0$. $\sin(x) = 0$ for all $x = \pi k$, where $k$ is some integer. So, this system has an infinite number of solutions of the form $(\pi k, 0)$; that is, the solution set for this system of equations is {..... $( -2\pi , 0 )$, $( -\pi , 0 )$, $( 0 , 0 )$, $( \pi , 0 )$, $( 2\pi , 0 )$ .....}.
\end{itemize}

%-------------------------------------------------------------------------------
\subsection{Linear Equations with Two Variables}
\label{sss:line-f2a6}
%-------------------------------------------------------------------------------

In the previous module, linear equations with two variables were discussed. A single
linear equation having two unknown variables is practically insufficient to solve or even
narrow down the solutions for the two variables, although it does establish a relationship
between them. The relationship is shown graphically as a line. Another linear equation
with the same two variables may be enough to narrow down the solution to the two equations
to one value for the first variable and one value for the second variable, i.e., to solve
the system of two simultaneous linear equations. Let''s see how two linear equations with
the same two unknowns might be related to each other. Since we said it was given that both
equations were linear, the graphs of both equations would be lines in the same
two-dimensional coordinate plane (for a system with two variables). The lines could be
related to each other in the following three ways:

1. The graphs of both equations could coincide giving the same line. This means that the
two equations are providing the same information about how the variables are related to
each other. The two equations are basically the same, perhaps just different versions or
forms of each other. Either one could be mathematically manipulated to produce the other
one. Both lines would have the same slope and the same y-intercept. Such equations are
considered dependent on each other. Since no new information is provided, the addition of
the second equation does not solve the problem by narrowing the solution set down to one
solution.

   \textbf{Example: Dependent linear equations}

   $6x - 3y = 12$

   $y = 2x - 4$

The above two equations provide the same information and result in the same graph, i.e.,
lines which coincide as shown in the following image.

% [Image: graph]

Let''s see how these equations can be mathematically manipulated to show they are
basically the same.

Divide both sides of the first equation $6x - 3y = 12$ by 3 to give
\[
2x - y = 4
\]
Now add $y$ to both sides
\[
2x = 4 + y
\]
Now subtract 4 from both sides
\[
y = 2x - 4
\]
This is the same as the second equation in the example. This is the slope-intercept form of the equation, from which a slope and a $y$-intercept unique to the line can be compared with any other equations in the slope-intercept form.

2. The graphs of two lines could be parallel although not the same. The two lines do not intersect each other at any point. This means there is no solution which satisfies both equations simultaneously, i.e., at the same time. The solution set for this system of simultaneous linear equations is the empty set. Such equations are considered inconsistent with each other and actually give contradictory information if it is claimed they are both true at the same time in the same problem. The parallel lines have equal slopes but different $y$-intercepts.

Sets of equations which have at least one common point which might provide a solution set
are consistent with each other. For example, the dependent equations mentioned previously
are consistent with each other.

\textbf{Example: Inconsistent linear equations}
\[
3x - 2y = -2
\]

\[
3x - 2y = 2
\]
To compare slopes and $y$-intercepts for these two linear equations, we place them in the slope-intercept forms. Subtract $3x$ from both sides of both equations.
\[
3x - 2y = -2
\]

\[
3x - 2y = 2
\]

\[
-2y = -3x - 2
\]

\[
-2y = -3x + 2
\]
Divide both sides of both equations by $-2$ and simplify to get slope-intercept forms for comparison.
\[
\frac{-2y}{-2} = \frac{-3x - 2}{-2}
\]

\[
\frac{-2y}{-2} = \frac{-3x + 2}{-2}
\]

\[
y = \frac{3}{2}x + 1
\]

\[
y = \frac{3}{2}x - 1
\]
Now, both slopes are equal at $\frac{3}{2}$, but the $y$-intercepts at 1 and -1 are different. The lines are parallel. The graphs are shown here:

% [Image: graph]

3. If the two lines are not the same and are not parallel, then they would intersect at one point because they are graphed in the same two-dimensional coordinate plane. The one point of intersection is the ordered pair of numbers which is the solution to the system of two linear equations and two unknowns. The two equations provide enough information to solve the problem and further equations are not needed. Such equations intersecting at a point providing a solution to the problem are considered independent of each other. The lines have different slopes but may or may not have the same $y$-intercept. Because such equations provide at least one solution point, they are consistent with each other.

\textbf{Example: Consistent independent linear equations}
\[
y = 3x - 5
\]

\[
y = -x - 1
\]
Both of these equations are given in the slope-intercept form, so it is easy to compare slopes and $y$-intercepts. For these two linear functions, both slopes are different and both $y$-intercepts are different. This means the lines are neither dependent nor inconsistent, so on a two-dimensional graph they must intersect at some point. In fact, the graph shows the lines intersecting at $(1, -2)$, which is the ordered pair solution to this system of independent simultaneous equations. Visual inspection of a graph cannot be relied on to give perfectly accurate coordinates every time, so either the point is tested with both equations or one of the following two methods is used to determine accurate coordinates for the intersection point.

% [Image: graph]

%-------------------------------------------------------------------------------
\subsection{Resources}
\label{sss:reso-55b5}
%-------------------------------------------------------------------------------

\begin{itemize}
    \item \href{https://tutorial.math.lamar.edu/classes/alg/systemstwovrble.aspx}{Linear Systems with Two Variables, Paul''s Online Notes}
    \item \href{https://openstax.org/books/college-algebra/pages/7-1-systems-of-linear-equations-two-variables}{Systems of Linear Equations with Two Variables, OpenStax}
    \item \href{https://openstax.org/books/college-algebra/pages/7-3-systems-of-nonlinear-equations-and-inequalities-two-variables}{Systems of Nonlinear Equations with Two Variables, OpenStax}
    \item \href{https://www.youtube.com/watch?v=ej25myhYcSg}{Solving Systems of Equations with Elimination, patrickJMT}
    \item \href{https://www.youtube.com/watch?v=cwHR_B9zK7k}{Solving Systems of Equations with Substitution, patrickJMT}
    \item \href{https://www.youtube.com/watch?v=WNkPKv0OTGI}{Solving Systems of Equations with Graphing, patrickJMT}
\end{itemize}

%-------------------------------------------------------------------------------
\subsection{Licensing}
\label{sss:lice-7352}
%-------------------------------------------------------------------------------

Content obtained and/or adapted from:

\begin{itemize}
    \item \href{https://openstax.org/books/college-algebra/pages/7-1-systems-of-linear-equations-two-variables}{Systems of Linear Equations: Two Variables, OpenStax: College Algebra under a CC BY license}
    \item \href{https://en.wikibooks.org/wiki/Algebra/Systems_of_Equations}{Systems of Equations, Wikibooks: Algebra under a CC BY-SA license}
\end{itemize}

%-------------------------------------------------------------------------------
\subsection{Learning Objectives}
\label{ssec:lear-d4a2}
%-------------------------------------------------------------------------------

\begin{enumerate}
    \item \textbf{Identify and Solve Systems with Unique Solutions}: Solve a system of two linear equations to find the unique point of intersection. Example: Given $y = 2x$ and $y = x + 5$, find the solution $(5, 10)$.
    \item \textbf{Determine Consistency and Parallelism}: Analyze systems of linear equations to determine if they are consistent, dependent, or inconsistent. For example, identify that $3x - 2y = -2$ and $3x - 2y = 2$ are inconsistent due to parallel lines with the same slope but different $y$-intercepts.
    \item \textbf{Graph and Interpret Systems with Multiple or Infinite Solutions}: Graph and interpret systems where equations intersect at multiple points or have infinite solutions. For example, recognize that $y = x^2$ and $y = 4$ intersect at $(2, 4)$ and $(-2, 4)$, showing multiple solutions.
\end{enumerate}

%-------------------------------------------------------------------------------
\subsection{Assessment Instrument}
\label{ssec:asse-d238}
%-------------------------------------------------------------------------------

Assess through problems requiring students to solve $2 \times 2$ linear systems by at
least two distinct methods (substitution and elimination) and to classify each system as
consistent independent, consistent dependent, or inconsistent based on its solution set.
Include at least one applied word problem that requires setting up the system from a
verbal description.

\textbf{Example problems.}

\begin{enumerate}
    \item Solve $\begin{cases} 2x + 3y = 12 \\ x - y = 1 \end{cases}$ by elimination and verify the solution by substitution.
    \item Determine whether $\begin{cases} 4x - 2y = 6 \\ 2x - y = 3 \end{cases}$ is consistent independent, consistent dependent, or inconsistent. Justify your answer.
    \item A store sells pens for \$1.50 and notebooks for \$3.00. Ten items cost \$21.00 in total. Write and solve a system of equations to find the number of each item bought.
\end{enumerate}

%===============================================================================
\section{Functions}
\label{sec:func-0eaa}
%===============================================================================

%-------------------------------------------------------------------------------
\subsection{Introduction}
\label{sss:intr-8800}
%-------------------------------------------------------------------------------

Formally, a function $f$ from a set $X$ to a set $Y$ is defined by a set $G$ of ordered pairs $(x, y)$ such that $x \in X$, $y \in Y$, and every element of $X$ is the first component of exactly one ordered pair in $G$. In other words, for every $x$ in $X$, there is exactly one element $y$ such that the ordered pair $(x, y)$ belongs to the set of pairs defining the function $f$. The set $G$ is called the graph of the function. Formally speaking, it may be identified with the function, but this hides the usual interpretation of a function as a process. Therefore, in common usage, the function is generally distinguished from its graph. Whenever one quantity uniquely determines the value of another quantity, we have a function. That is, the set $X$ uniquely determines the set $Y$. You can think of a function as a kind of machine. You feed the machine raw materials, and the machine changes the raw materials into a finished product.

%-------------------------------------------------------------------------------
\subsection{Functions in Everyday Life}
\label{sss:func-5465}
%-------------------------------------------------------------------------------

Think about dropping a ball from a bridge. At each moment in time, the ball is a height
above the ground. The height of the ball is a function of time. It was the job of
physicists to come up with a formula for this function. This type of function is called
real-valued since the ``finished product'' is a number (or, more specifically, a real
number).

Think about a wind storm. At different places, the wind can be blowing in different directions with different intensities. The direction and intensity of the wind can be thought of as a function of position. This is a function of two real variables (a location is described by two values - an $x$ and a $y$) which results in a vector (which is something that can be used to hold a direction and an intensity). These functions are studied in multivariable calculus (which is usually studied after a one-year college-level calculus course). This is a vector-valued function of two real variables.

We will be looking at real-valued functions until studying multivariable calculus. Think
of a real-valued function as an input-output machine; you give the function an input, and
it gives you an output which is a number (more specifically, a real number). For example,
the squaring function takes the input 4 and gives the output value 16. The same squaring
function takes the input -1 and gives the output value 1.

This is an intuitive way to understand functions: a machine that makes the input $x$ go through a transformation $f$ into the output $f(x)$.

% [Image: image]

%-------------------------------------------------------------------------------
\subsection{Notation}
\label{sss:nota-71f9}
%-------------------------------------------------------------------------------

Functions are used so much that there is a special notation for them. The notation is
somewhat ambiguous, so familiarity with it is important in order to understand the
intention of an equation or formula.

Though there are no strict rules for naming a function, it is standard practice to use the letters $f$, $g$, and $h$ to denote functions, and the variable $x$ to denote an independent variable. $y$ is used for both dependent and independent variables.

When discussing or working with a function $f$, it''s important to know not only the function but also its independent variable $x$. Thus, when referring to a function $f$, you usually do not write $f$, but instead $f(x)$. The function is now referred to as `` $f$ of $x$ ''. The name of the function is adjacent to the independent variable (in parentheses). This is useful for indicating the value of the function at a particular value of the independent variable. For instance, if 
\[
f(x) = 7x + 1
\]
and if we want to use the value of $f$ for $x$ equal to 2, then we would substitute 2 for $x$ on both sides of the definition above and write
\[
f(2) = 7(2) + 1 = 14 + 1 = 15
\]
This notation is more informative than leaving off the independent variable and writing simply '' $f$ '', but can be ambiguous since the parentheses next to $f$ can be misinterpreted as multiplication, $2f$. To make sure nobody is too confused, follow this procedure:

1. Define the function $f$ by equating it to some expression.

2. Give a sentence like the following: ``At $x = c$, the function $f$ is...''

3. Evaluate the function.

%-------------------------------------------------------------------------------
\subsection{Description}
\label{sss:desc-67da}
%-------------------------------------------------------------------------------

There are many ways people describe functions. In the examples above, a verbal description
is given (the height of the ball above the earth as a function of time). Here is a list of
ways to describe functions. The top three listed approaches to describing functions are
the most popular.

1. A function is given a name (such as $f$) and a formula for the function is also given. For example, $f(x) = 3x + 2$ describes a function. We refer to the input as the argument of the function (or the independent variable), and to the output as the value of the function at the given argument.

2. A function is described using an equation and two variables. One variable is for the input of the function and one is for the output of the function. The variable for the input is called the independent variable. The variable for the output is called the dependent variable. For example, $y = 3x + 2$ describes a function. The dependent variable appears by itself on the left-hand side of the equal sign.

3. A verbal description of the function.

When a function is given a name (like in number 1 above), the name of the function is usually a single letter of the alphabet (such as $f$ or $g$). Some functions whose names are multiple letters (like the sine function $y = \sin(x)$).

%-------------------------------------------------------------------------------
\subsection{Plugging a Value into a Function}
\label{sss:plug-50c6}
%-------------------------------------------------------------------------------

If we write $f(x) = 3x + 2$, then we know that:

\begin{itemize}
    \item The function $f$ is a function of $x$.
    \item To evaluate the function at a certain number, replace the $x$ with that number.
    \item Replacing $x$ with that number in the right side of the function will produce the function''s output for that certain input.
\end{itemize}

In English, the definition of $f$ is interpreted, ``Given a number, $f$ will return two more than the triple of that number.''

How would we know the value of the function $f$ at 3? We would have the following three thoughts:
\[
f(3) = 3(3) + 2
\]

\[
3(3) + 2 = 9 + 2
\]

\[
9 + 2 = 11
\]
and we would write
\[
f(3) = 3(3) + 2 = 9 + 2 = 11
\]
The value of $f$ at 3 is 11.

Note that $f(3)$ means the value of the dependent variable when $x$ takes on the value of 3. So we see that the number 11 is the output of the function when we give the number 3 as the input. People often summarize the work above by writing ``the value of $f$ at three is eleven'', or simply ``$f$ of three equals eleven''.

%-------------------------------------------------------------------------------
\subsection{Basic Concepts of Functions}
\label{sss:basi-1380}
%-------------------------------------------------------------------------------

The formal definition of a function states that a function is actually a mapping that associates the elements of one set called the domain of the function, $A$, with the elements of another set called the range of the function, $B$. For each value we select from the domain of the function, there exists exactly one corresponding element in the range of the function. The definition of the function tells us which element in the range corresponds to the element we picked from the domain. An example is given below.

%-------------------------------------------------------------------------------
\subsection{Example}
\label{sss:exam-1a79}
%-------------------------------------------------------------------------------

Let function
\[
f(x) = 5x + \frac{5x}{2}
\]
for all $x$. For what value of $x$ gives 
\[
f(x) = 225 \frac{1}{2}
\]
In mathematics, it is important to notice any repetition. If something seems to repeat,
keep a note of that in the back of your mind somewhere.

Here, notice that
\[
f(x) = 5x + \frac{5x}{2}
\]
and
\[
f(x) = 225 \frac{1}{2}
\]
Because $f(x)$ is equal to two different things, it must be the case that the other side of the equal sign to $f(x)$ is equal to the other. This property is known as the transitive property and could thus make the following equation true:
\[
225 \frac{1}{2} = 5x + \frac{5x}{2}
\]
Next, simplify :  make your life easier rather than harder. In this instance, since $5x + \frac{5x}{2}$ has $x$ as a like-term, and the two terms are fractions added to the other, put it over a common denominator and simplify. Similarly, since $225 \frac{1}{2}$ is a mixed fraction, $225 \frac{1}{2} = 225 + \frac{1}{2}$.
\[
225 \frac{1}{2} = 5x + \frac{5x}{2}
\]

\[
\Rightarrow \frac{225}{1} + \frac{1}{2} = \frac{5x}{1} + \frac{5x}{2}
\]

\[
\Rightarrow \frac{450}{2} + \frac{1}{2} = \frac{10x}{2} + \frac{5x}{2}
\]

\[
\Rightarrow \frac{451}{2} = \frac{15x}{2}
\]
Multiply both sides by the reciprocal of the coefficient of $x$, $\frac{2}{15}$ from both sides by multiplying by it.
\[
\frac{2}{15} \cdot \frac{451}{2} = x
\]

\[
\frac{451}{15} = x
\]
or
\[
x = \frac{451}{15}
\]
The value of $x$ that makes 
\[
f(x) = 225 \frac{1}{2}
\]
is
\[
x = \frac{451}{15}
\]
Classically, the element picked from the domain is pictured as something that is fed into
the function and the corresponding element in the range is pictured as the output. Since
we ``pick'' the element in the domain whose corresponding element in the range we want to
find, we have control over what element we pick and hence this element is also known as
the ``independent variable''. The element mapped in the range is beyond our control and is
``mapped to'' by the function. This element is hence also known as the ``dependent
variable'', for it depends on which independent variable we pick. Since the elementary
idea of functions is better understood from the classical viewpoint, we shall use it
hereafter. However, it is still important to remember the correct definition of functions
at all times.

%-------------------------------------------------------------------------------
\subsection{Basic Types of Transformation}
\label{sss:basi-cd64}
%-------------------------------------------------------------------------------

To make it simple, for the function $f(x)$, all of the possible $x$ values constitute the domain, and all of the values $f(x)$ ($y$ on the $x$-$y$ plane) constitute the range. To put it in more formal terms, a function $f$ is a mapping of some element $a \in A$, called the domain, to exactly one element $b \in B$, called the range, such that $f:A \to B$. The image below should help explain the modern definition of a function:

The image demonstrates a mapping of some element $a$ (the circle) in $A$, the domain, to exactly one element $b$ in $B$, the range. $A$ is the domain of the function while $B$ is the range. This transformation from set $A$ to $B$ is an example of a one-to-one function.

% [Image: image]

A function is considered one-to-one if an element $a \in A$ from domain $A$ of function $f$ leads to exactly one element $b \in B$ from range $B$ of the function. By definition, since only one element $b$ is mapped by function $f$ from some element $a$, $f:A \to B$ implies that there exists only one element $b$ from the mapping. Therefore, there exists a one-to-one function because it complies with the definition of a function. 

A function is considered many-to-one if some elements $a_1, a_2, \ldots a_n \in A$ from domain $A$ of function $f$ maps to exactly one element $b \in B$ from range $B$ of the function. Since some elements $(a_1, a_2, \ldots a_n) \in A$ must map onto exactly one element $b \in B$, $f:A \to B$ must be compliant with the definition of a function. Therefore, there exists a many-to-one function.

A function is considered one-to-many if exactly one element $a \in A$ from domain $A$ of function $f$ maps to some elements $(b_1, b_2, \ldots b_n) \in B$ from range $B$ of the function. If some element $a$ maps onto many distinct elements $b$, then $f:A \to B$ is non-functional since there exist many distinct elements $b$. Given many-to-one is non-compliant with the definition of a function, there exists no function that is one-to-many.

The modern definition describes a function sufficiently such that it helps us determine
whether some new type of ``function'' is indeed one. Now that each case is defined above,
you can now prove whether functions are one of these function cases. Here is an example
problem:

%-------------------------------------------------------------------------------
\subsection{One-to-One Example}
\label{sss:onet-5bc6}
%-------------------------------------------------------------------------------

Given: $f(x) = ax + b$, where $a$ and $b$ are constant and $a \neq 0$.

Prove: function $f$ is one-to-one.

Notice that the only changing element in the function $f$ is $x$. To prove a function is one-to-one by applying the definition may be impossible because although two random elements of domain set $A$ may not be many-to-one, there may be some elements $(a_1, a_2, \ldots a_n) \in A$ that may make the function many-to-one. To account for this possibility, we must prove that it is impossible to have some result like that.

Assume there exist two distinct elements $x_1 \neq x_2$ that will result in $f(x_1) = f(x_2)$. This would make the function many-to-one. In consequence,
\[
ax_1 + b = ax_2 + b
\]
Subtract $b$ from both sides of the equation.
\[
\Rightarrow ax_1 = ax_2
\]
Subtract $ax_2$ from both sides of the equation.
\[
\Rightarrow ax_1 - ax_2 = 0
\]
Factor $a$ from both terms to the left of the equation.
\[
\Rightarrow a(x_1 - x_2) = 0
\]
Multiply $\frac{1}{a}$ to both sides of the equation.
\[
\Rightarrow (x_1 - x_2) = 0
\]
Add $x_2$ to both sides of the equation.
\[
\Rightarrow x_1 = x_2
\]
Notice that $x_1 = x_2$. However, this is impossible because $x_1$ and $x_2$ are distinct. Ergo, $f(x_1) \neq f(x_2)$. No two distinct inputs can give the same output. Therefore, the function must be one-to-one.

%-------------------------------------------------------------------------------
\subsection{Basic Concepts}
\label{sss:basi-3f1b}
%-------------------------------------------------------------------------------

The domain is all the elements in set $A$ that can be mapped to the elements in set $B$. The range is those elements in set $B$ which are mapped with the domain. The codomain is all the elements in set $B$.

% [Image: image]

There are a few more important ideas to learn from the modern definition of the function, and it comes from knowing the difference between domain, range, and codomain. We already discussed some of these, yet knowing about sets adds a new definition for each of the following ideas. Therefore, let us discuss these based on these new ideas. Let $A$ and $B$ be a set. If we were to combine these two sets, it would be defined as the cartesian cross product $A \times B$. The subset of this product is the function. The below definitions are a little confusing. The best way to interpret these is to draw an image. To the right of these definitions is the image that corresponds to it.

%-------------------------------------------------------------------------------
\subsection{Definition of Domain, Range, and Codomain of a Function}
\label{sss:defi-62c6}
%-------------------------------------------------------------------------------

The domain is defined to be a set $A$ with all elements $a \in A$ mapping to at least one unique $b \in B$.

The set of elements in $B$ is the range of the function mapping $f:A \to B$ in the cartesian cross product, whereby the set of all elements $a \in A$ maps to some element $b \in B$.

The codomain is the set $B$, where it is not necessarily the case that all elements $b \in B$ were mapped by some $a \in A$.

Note that the codomain is not as important as the other two concepts.

Take $f(x) = \sqrt{1 - x^2}$ for example:

Because of the square root, the content in the square root should be strictly not smaller
than 0.

  \[
1 - x^2 \geq 0
\]

  \[
\Leftrightarrow -1 \leq x \leq 1
\]

Thus the domain is

  \[
x \in [-1, 1]
\]

\begin{itemize}
    \item \textbf{The domain of the function is the interval from -1 to 1}:
\end{itemize}
% [Image: domain]

Correspondingly, the range will be

  \[
f(x) \in [0, 1]
\]

\begin{itemize}
    \item \textbf{The range of the function is the interval from 0 to 1}:
\end{itemize}
% [Image: range]

%-------------------------------------------------------------------------------
\subsection{Other Types of Transformation}
\label{sss:othe-065a}
%-------------------------------------------------------------------------------

There is one more final aspect that needs to be defined. We already have a good idea of
what makes a mapping a function (e.g. it cannot be one-to-many). However, three more
definitions that you will often hear will be necessary to talk about: injective,
surjective, bijective.

% [Image: mapping]

\begin{itemize}
    \item The function mapping on the left is an example of an injective function because it is one-to-one. However, it is not surjective because the range and the codomain are not the same.
    \item A function is \textbf{injective} if it is one-to-one.
    \item A function is \textbf{surjective} if it is ``onto.'' That is, the mapping $f:A \to B$ has $B$ as the range of the function, where the codomain and the range of the function are the same.
    \item A function is \textbf{bijective} if it is both surjective and injective.
\end{itemize}

Again, the above definitions are often very confusing. Another image is provided to the right of the bullet points. Along with that, another example is also provided. Let us analyze the function $g(x) = ax^2$.

Given: $g(x) = ax^2$, where $a$ is constant and $a \neq 0$.

Prove: function $g$ is non-surjective and non-injective.

Notice that the only changing element in the function $g$ is $x$. Let us check to see the conditions of this function.

Is it injective? Let $g(x) = y$, and solve for $x$. First, divide by $a$.
\[
y = ax^2
\]

\[
\Rightarrow \frac{y}{a} = x^2
\]
Then take the square root of $x^2$. By definition, 
\[
\sqrt{x^2} = \pm x
\]
so
\[
\frac{y}{a} = x^2 \Rightarrow \pm \sqrt{\frac{y}{a}} = \pm x
\]
Notice, however, what we learned from the above manipulation. As a result of solving for $x$, we found that there are two solutions for $x$. However, this resulted in two different values from $\frac{y}{a}$. This means that for some individual $x$ that gives $y$, there are two different inputs that result in the same value. Because $f(x_1) = f(x_2)$ when $x_1 \neq x_2$, this is by definition non-injective.

Is it surjective? As a natural consequence of what we learned about As a result of solving for $x$, we found that there are two solutions for $x$. However, this resulted in two different values from $\frac{y}{a}$. This means that for some individual $x$ that gives $y$, there are two different inputs that result in the same value. Because $f(x_1) = f(x_2)$ when $x_1 \neq x_2$, this is by definition non-injective.

Is it surjective? As a natural consequence of what we learned about inputs, let us
determine the range of the function. After all, the only way to determine if something is
surjective is to see if the range applies to all real numbers. A good way to determine
this is by finding a pattern involving our domains. Let us say we input a negative number
for the domain:
\[
f(-x) = a(-x)^2 = a \cdot (-x) \cdot (-x) = ax^2
\]
This seemingly trivial exercise tells us that negative numbers give us non-negative numbers for our range. This is huge information! For $x \in \mathbb{R}$, the function always results $y > 0$ for our range. For the set $B$, we have elements in that set that have no mappings from the set $A$. As such, $y \leq 0$ is the codomain of set $B$. Therefore, this function is non-surjective!

%-------------------------------------------------------------------------------
\subsection{Tests for Equations}
\label{sss:test-28db}
%-------------------------------------------------------------------------------

%-------------------------------------------------------------------------------
\subsection{The Vertical Line Test}
\label{sss:thev-18cb}
%-------------------------------------------------------------------------------

The vertical line test is a systematic test to find out if an equation involving $x$ and $y$ can serve as a function (with $x$ the independent variable and $y$ the dependent variable). Simply graph the equation and draw a vertical line through each point of the $x$-axis. If any vertical line ever touches the graph at more than one point, then the equation is not a function; if the line always touches at most one point of the graph, then the equation is a function.

The circle, below, is not a function because the vertical line intercepts two points on the graph when $x = 1$.

% [Image: vertical line test]

%-------------------------------------------------------------------------------
\subsection{The Horizontal Line and the Algebraic 1-1 Test}
\label{sss:theh-25b4}
%-------------------------------------------------------------------------------

Similarly, the horizontal line test, though it does not test if an equation is a function,
tests if a function is injective (one-to-one). If any horizontal line ever touches the
graph at more than one point, then the function is not one-to-one; if the line always
touches at most one point on the graph, then the function is one-to-one.

The algebraic 1-1 test is the non-geometric way to see if a function is one-to-one. The
basic concept is that:

Assume there is a function $f$. If:

\[
f(a) = f(b) \text{ and } a = b,
\]

then

function $f$ is one-to-one.

Here is an example: prove that $f(x) = \frac{1 - 2x}{1 + x}$ is injective.

Since the notation is the notation for a function, the equation is a function. So we only need to prove that it is injective. Let $a$ and $b$ be the inputs of the function and that $f(a) = f(b)$. Thus,
\[
\begin{aligned}
\frac{1 - 2a}{1 + a} = \frac{1 - 2b}{1 + b}
&\;\Longleftrightarrow\; (1 + b)(1 - 2a) = (1 + a)(1 - 2b) \\
&\;\Longleftrightarrow\; 1 - 2a + b - 2ab = 1 - 2b + a - 2ab \\
&\;\Longleftrightarrow\; 1 - 2a + b = 1 - 2b + a \\
&\;\Longleftrightarrow\; -2a + b = -2b + a \\
&\;\Longleftrightarrow\; -3a = -3b \\
&\;\Longleftrightarrow\; a = b.
\end{aligned}
\]
So, the result is $a = b$, proving that the function is injective.

Another example is proving that $g(x) = x^2$ is not injective.

Using the same method, one can find that $a = \pm b$, which is not $a = b$. So, the function is not injective.

%-------------------------------------------------------------------------------
\subsection{Remarks}
\label{sss:rema-46b8}
%-------------------------------------------------------------------------------

The following arise as a direct consequence of the definition of functions:

\begin{itemize}
    \item By definition, for each ``input'' a function returns only one ``output'', corresponding to that input. While the same output may correspond to more than one input, one input cannot correspond to more than one output. This is expressed graphically as the vertical line test: a line drawn parallel to the axis of the dependent variable (normally vertical) will intersect the graph of a function only once. However, a line drawn parallel to the axis of the independent variable (normally horizontal) may intersect the graph of a function as many times as it likes. Equivalently, this has an algebraic (or formula-based) interpretation. We can always say if $a = b$, then $f(a) = f(b)$, but if we only know that $f(a) = f(b)$ then we can''t be sure that $a = b$.
    \item Each function has a set of values, the function''s domain, which it can accept as input. Perhaps this set is all positive real numbers; perhaps it is the set {pork, mutton, beef}. This set must be implicitly/explicitly defined in the definition of the function. You cannot feed the function an element that isn''t in the domain, as the function is not defined for that input element.
    \item Each function has a set of values, the function''s range, which it can output. This may be the set of real numbers. It may be the set of positive integers or even the set {0, 1}. This set, too, must be implicitly/explicitly defined in the definition of the function.
\end{itemize}

Functions are an important foundation of mathematics. This modern interpretation is a result of the hard work of the mathematicians of history. It was not until recently that the definition of the relation was introduced by René Descartes in \textit{Geometry} (1637). Nearly a century later, the standard notation ($f(x) = y$) was first introduced by Leonhard Euler in \textit{Introductio in Analysin Infinitorum} and \textit{Institutiones Calculi Differentialis}. The term function was also a new innovation during Euler''s time as well, being introduced by Gottfried Wilhelm Leibniz in a 1673 letter ``to describe a quantity related to points of a curve, such as a coordinate or curve''s slope.'' Finally, the modern definition of the function being the relation among sets was first introduced in 1908 by Godfrey Harold Hardy where there is a relation between two variables $x$ and $y$ such that ``to some values of $x$ at any rate correspond values of $y$.''

%-------------------------------------------------------------------------------
\subsection{Important Functions}
\label{sss:impo-1cf2}
%-------------------------------------------------------------------------------

%-------------------------------------------------------------------------------
\subsection{Polynomials}
\label{sss:poly-f9e4}
%-------------------------------------------------------------------------------

Polynomial functions are the most common and most convenient functions in the world of calculus. Their frequent appearances and their applications on computer calculations have made them important. A polynomial in a single indeterminate $x$ can always be written (or rewritten) in the form:
\[
f(x) = a_n x^n + a_{n-1} x^{n-1} + a_{n-2} x^{n-2} + \cdots + a_2 x^2 + a_1 x + a_0
\]
To be more concise, it can also be written in the summation form:
\[
\sum_{k=0}^n a_k x^k
\]

%-------------------------------------------------------------------------------
\subsection{Constant}
\label{sss:cons-617a}
%-------------------------------------------------------------------------------

When $n = 0$, the polynomial can be rewritten into the following function:
\[
f(x) = C,
\]
where $C$ is a constant.

The graph of this function is a horizontal line passing through the point $(0, C)$.

%-------------------------------------------------------------------------------
\subsection{Linear}
\label{sss:line-9a93}
%-------------------------------------------------------------------------------

Two linear functions are shown in the image. One is $f(x) = \frac{1}{2}x + 2$ and the other is $g(x) = -2x + 5$. 

% [Image: constant]

When $n = 1$, the polynomial can be rewritten into
\[
f(x) = mx + b,
\]
where $m$ and $b$ are constants.

The graph of this function is a straight line passing through the points $(0, b)$ and $\left(-\frac{b}{m}, 0\right)$, and the slope of this function is $m$.

%-------------------------------------------------------------------------------
\subsection{Quadratic}
\label{sss:quad-d2ac}
%-------------------------------------------------------------------------------

When $n = 2$, the polynomial can be rewritten into
\[
f(x) = ax^2 + bx + c,
\]
where $a$, $b$, $c$ are constants.

The graph of this function is a parabola, like the trajectory of a basketball thrown into
the hoop.

% [Image: graph]

If there are questions about the quadratic formula and other basic polynomial concepts,
please review the respective chapters in Algebra.

%-------------------------------------------------------------------------------
\subsection{Trigonometric}
\label{sss:trig-4a83}
%-------------------------------------------------------------------------------

Trigonometric functions are also very important because they can connect algebra and
geometry. Trigonometric functions are not explained in detail here due to their importance
and difficulty.

%-------------------------------------------------------------------------------
\subsection{Exponential and Logarithmic}
\label{sss:expo-d74f}
%-------------------------------------------------------------------------------

Exponential and logarithmic functions have a unique identity when calculating the
derivatives, so this is a great time to review those functions. The exponential function
is defined as:
\[
f(x) = a^x,
\]
where $a$ is a constant.

The logarithmic function is defined as:
\[
f(x) = \log_b x,
\]
where $b$ is a constant.

A special number will be frequently seen in these functions: Euler''s constant, also known as the base of the natural logarithm. Notated as $e$, it is defined as
\[
e = \sum_{n=0}^{\infty} \frac{1}{n!}.
\]
The curve on the left is an exponential function while the curve on the right is a
logarithmic one:

% [Image: log]

%-------------------------------------------------------------------------------
\subsection{Signum}
\label{sss:sign-327b}
%-------------------------------------------------------------------------------

The Signum (sign) function is defined as:

\begin{itemize}
    \item \(sgn(x) = -1\) if \(x < 0\)
    \item \(sgn(x) = 0\) if \(x = 0\)
    \item \(sgn(x) = 1\) if \(x > 0\)
\end{itemize}

%-------------------------------------------------------------------------------
\subsection{Properties of Functions}
\label{sss:prop-3b03}
%-------------------------------------------------------------------------------

Sometimes, a lot of function manipulations cannot be achieved without some help from basic
properties of functions.

%-------------------------------------------------------------------------------
\subsection{Domain and Range}
\label{sss:doma-1606}
%-------------------------------------------------------------------------------

This concept is discussed above.

%-------------------------------------------------------------------------------
\subsection{Growth}
\label{sss:grow-5ff5}
%-------------------------------------------------------------------------------

Although it seems obvious to spot a function increasing or decreasing, without the help of
graphing software, we need a mathematical way to spot whether the function is increasing
or decreasing, or else our human minds cannot immediately comprehend the huge amount of
information.

Assume a function $f(x)$ with inputs $x_1$, $x_2$, and that $x_1 \in [a,b]$, $x_2 \in [a,b]$, and $x_2 > x_1$ at all times.

If for all $x_1$ and $x_2$, $f(x_2) - f(x_1) > 0$, then

$f(x)$ is increasing in $[a,b]$

If for all $x_1$ and $x_2$, $f(x_2) - f(x_1) < 0$, then

$f(x)$ is decreasing in $[a,b]$

%-------------------------------------------------------------------------------
\subsection{Example}
\label{sss:exam-1a79}
%-------------------------------------------------------------------------------

In which intervals is $f(x) = \frac{1}{1 - x^2}$ increasing?

Firstly, the domain is important. Because the denominator cannot be 0, the domain for this
function is

$x \in (-\infty, -1) \cup (-1, 1) \cup (1, \infty)$

In $(-\infty, -1)$, the growth of the function is:

Let $x_1, x_2 \in [-\infty, -1]$ and $x_2 > x_1$. Thus,
\[
f(x_2) - f(x_1) = \frac{1}{1 - x_2^2} - \frac{1}{1 - x_1^2} = \frac{x_2^2 - x_1^2}{(1 - x_2^2)(1 - x_1^2)}
\]
$\because$ both $x_1, x_2 \in [-\infty, -1]$

$\therefore (1 - x_2^2)(1 - x_1^2) > 0$

$\because x_2 > x_1$ and $x_1 < x_2 < 0$

$\therefore x_2^2 - x_1^2 < 0$

So,
\[
f(x_2) - f(x_1) = \frac{x_2^2 - x_1^2}{(1 - x_2^2)(1 - x_1^2)} < 0
\]
$f(x)$ is decreasing in $(-\infty, -1)$

In $(-1, 1)$

Let $x_1, x_2 \in [-1, 1]$ and $x_2 > x_1$. Thus,
\[
f(x_2) - f(x_1) = \frac{1}{1 - x_2^2} - \frac{1}{1 - x_1^2} = \frac{x_2^2 - x_1^2}{(1 - x_2^2)(1 - x_1^2)}
\]
$\because$ both $x_1, x_2 \in [-1, 1]$

$\therefore (1 - x_2^2)(1 - x_1^2) > 0$

However, the sign of $x_2^2 - x_1^2$ in $(-1, 1)$ cannot be determined. It can only be determined in $(-1, 0)$ and $(0, 1)$.

In $(-1, 0)$

$\because x_2 > x_1$ and $x_1 < x_2 < 0$

$\therefore x_2^2 - x_1^2 < 0$

In $(0, 1)$

$\because x_2 > x_1$ and $0 < x_1 < x_2$

$\therefore x_2^2 - x_1^2 > 0$

As a result, $f(x)$ is decreasing in $(-1, 0)$ and increasing in $(0, 1)$.

In $(1, \infty)$

Let $x_1, x_2 \in [1, \infty]$ and $x_2 > x_1$. Thus,
\[
f(x_2) - f(x_1) = \frac{1}{1 - x_2^2} - \frac{1}{1 - x_1^2} = \frac{x_2^2 - x_1^2}{(1 - x_2^2)(1 - x_1^2)}
\]
$\because$ both $x_1, x_2 \in [1, \infty]$

$\therefore (1 - x_2^2)(1 - x_1^2) > 0$

$\because x_2 > x_1$ and $0 < x_1 < x_2$

$\therefore x_2^2 - x_1^2 > 0$

So,
\[
f(x_2) - f(x_1) = \frac{x_2^2 - x_1^2}{(1 - x_2^2)(1 - x_1^2)} > 0
\]
$f(x)$ is increasing in $(1, \infty)$.

Therefore, the intervals in which the function is increasing are $(0, 1) \cup (1, \infty)$.

$\blacksquare$

After learning derivatives, there will be more ways to determine the growth of a function.

%-------------------------------------------------------------------------------
\subsection{Parity}
\label{sss:pari-aaba}
%-------------------------------------------------------------------------------

The properties odd and even are associated with symmetry. While even functions have a symmetry about the $y$-axis, odd functions are symmetric about the origin. In mathematical terms:

A function is even when $f(-x) = f(x)$ 

A function is odd when $f(-x) = -f(x)$

%-------------------------------------------------------------------------------
\subsection{Example}
\label{sss:exam-1a79}
%-------------------------------------------------------------------------------

Prove that $f(x) = \frac{1}{1 - x^2}$ is an even function.

$\because f(-x) = \frac{1}{1 - (-x)^2} = \frac{1}{1 - x^2} = f(x)$

$\therefore f(x)$ is an even function

%-------------------------------------------------------------------------------
\subsection{Manipulating Functions}
\label{sss:mani-12b7}
%-------------------------------------------------------------------------------

%-------------------------------------------------------------------------------
\subsection{Addition, Subtraction, Multiplication, and Division of Functions}
\label{sss:addi-f096}
%-------------------------------------------------------------------------------

For two real-valued functions, we can add the functions, multiply the functions, raised to
a power, etc. For example:

\begin{itemize}
    \item If we add the functions $y = 3x + 2$ and $y = x^2$, we obtain $y = x^2 + 3x + 2$.
    \item If we subtract $y = 3x + 2$ from $y = x^2$, we obtain $y = x^2 - (3x + 2)$. We can also write this as $y = x^2 - 3x - 2$.
    \item If we multiply the function $y = 3x + 2$ and the function $y = x^2$, we obtain $y = (3x + 2)x^2$. We can also write this as $y = 3x^3 + 2x^2$.
    \item If we divide the function $y = 3x + 2$ by the function $y = x^2$, we obtain $y = \frac{3x + 2}{x^2}$.
\end{itemize}

If a math problem wants you to add two functions $f$ and $g$, there are two ways that the problem will likely be worded:

\begin{itemize}
    \item If you are told that $f(x) = 3x + 2$, that $g(x) = x^2$, that $h(x) = f(x) + g(x)$ and asked about $h$, then you are being asked to add two functions. Your answer would be $h(x) = x^2 + 3x + 2$.
    \item If you are told that $f(x) = 3x + 2$, that $g(x) = x^2$ and you are asked about $f + g$, then you are being asked to add two functions. The addition of $f$ and $g$ is called $f + g$. Your answer would be $(f + g)(x) = x^2 + 3x + 2$.
\end{itemize}

Similar statements can be made for subtraction, multiplication, and division.

Let $f(x) = 3x + 2$ and $g(x) = x^2$. Let''s add, subtract, multiply, and divide.

\[
(f + g)(x) = f(x) + g(x) = (3x + 2) + (x^2) = x^2 + 3x + 2
\]

\[
(f - g)(x) = f(x) - g(x) = (3x + 2) - (x^2) = -x^2 + 3x + 2
\]

\[
(f \times g)(x) = f(x) \times g(x) = (3x + 2) \times (x^2) = 3x^3 + 2x^2
\]

\[
\left(\frac{f}{g}\right)(x) = \frac{f(x)}{g(x)} = \frac{3x + 2}{x^2}
\]

%-------------------------------------------------------------------------------
\subsection{Composition of Functions}
\label{sss:comp-37f3}
%-------------------------------------------------------------------------------

We begin with a fun (and not too complicated) application of composition of functions
before we talk about what composition of functions is.

%-------------------------------------------------------------------------------
\subsection{Example: Dropping a Ball}
\label{sss:exam-c784}
%-------------------------------------------------------------------------------

If we drop a ball from a bridge which is 20 meters above the ground, then the height of
our ball above the earth is a function of time. The physicists tell us that if we measure
time in seconds and distance in meters, then the formula for height in terms of time is
\[
h = -4.9t^2 + 20
\]
Suppose we are tracking the ball with a camera and always want the ball to be in the
center of our picture. Suppose we have
\[
\theta = f(h)
\]
The angle will depend upon the height of the ball above the ground and the height above
the ground depends upon time. So the angle will depend upon time. This can be written as
\[
\theta = f(-4.9t^2 + 20)
\]
We replace $h$ with what it is equal to. This is the essence of composition.

%-------------------------------------------------------------------------------
\subsection{Composition of Functions}
\label{sss:comp-37f3}
%-------------------------------------------------------------------------------

Composition of functions is another way to combine functions which is different from
addition, subtraction, multiplication, or division.

The value of a function $f$ depends upon the value of another variable $x$; however, that variable could be equal to another function $g$, so its value depends on the value of a third variable. If this is the case, then the first variable is a function $h$ of the third variable; this function ($h$) is called the composition of the other two functions ($f$ and $g$).

%-------------------------------------------------------------------------------
\subsection{Example: Composing Two Functions}
\label{sss:exam-d4fc}
%-------------------------------------------------------------------------------

Let
\[
f(x) = 3x + 2
\]
and
\[
g(x) = x^2
\]
The composition of $f$ with $g$ is read as either ``f composed with g'' or ``f of g of x.''

Let
\[
h(x) = f(g(x))
\]
Then
\[
h(x) = f(g(x)) = f(x^2) = 3(x^2) + 2 = 3x^2 + 2
\]
Sometimes a math problem asks you to compute $(f \circ g)(x)$ when they want you to compute $f(g(x))$.

Here, $h$ is the composition of $f$ and $g$ and we write $h = f \circ g$. Note that composition is not commutative:
\[
f(g(x)) = 3x^2 + 2
\]
and
\[
g(f(x)) = g(3x + 2) = (3x + 2)^2 = 9x^2 + 12x + 4
\]
so
\[
f(g(x)) \neq g(f(x))
\]
Composition of functions is very common, mainly because functions themselves are common.
For instance, squaring and sine are both functions:
\[
\text{square}(x) = x^2
\]

\[
\text{sine}(x) = \sin(x)
\]
Thus, the expression
\[
\sin^2(x)
\] 
is a composition of functions:
\[
\sin^2(x) = \text{square}(\sin(x)) = \text{square}(\text{sine}(x))
\]
(Note that this is not the same as
\[
\text{sine}(\text{square}(x)) = \sin(x^2)
\]
Since the function sine equals $\frac{1}{2}$ if $x = \frac{\pi}{6}$,
\[
\text{square}(\sin(\frac{\pi}{6})) = \text{square}(\frac{1}{2})
\]
Since the function square equals $\frac{1}{4}$ if $x = \frac{\pi}{6}$,
\[
\sin^2(\frac{\pi}{6}) = \text{square}(\sin(\frac{\pi}{6})) = \text{square}(\frac{1}{2}) = \frac{1}{4}
\]

%-------------------------------------------------------------------------------
\subsection{Transformations}
\label{sss:tran-dfc3}
%-------------------------------------------------------------------------------

Transformations are a type of function manipulation that are very common. They consist of
multiplying, dividing, adding, or subtracting constants to either the input or the output.
Multiplying by a constant is called dilation and adding a constant is called translation.
Here are a few examples:

\begin{itemize}
    \item $ f(2 \times x) $ :  Dilation
    \item $ f(x + 2) $ :  Translation
    \item $ 2 \times f(x) $ :  Dilation
    \item $ 2 + f(x) $ :  Translation
\end{itemize}

%-------------------------------------------------------------------------------
\subsection{Examples of Horizontal and Vertical Translations}
\label{sss:exam-30cf}
%-------------------------------------------------------------------------------

Examples of horizontal and vertical translations can be seen below.

% [Image: translation]

The red graphs represent functions in their ''original'' state, the solid blue graphs have
been translated (shifted) horizontally, and the dashed graphs have been translated
vertically.

%-------------------------------------------------------------------------------
\subsection{Examples of Horizontal and Vertical Dilations}
\label{sss:exam-e155}
%-------------------------------------------------------------------------------

% [Image: dilations]

Dilations are demonstrated in a similar fashion. The function
\[
f(2 \times x)
\]
has had its input doubled. One way to think about this is that now any change in the input will be doubled. If I add one to $x$, I add two to the input of $f$, so it will now change twice as quickly. Thus, this is a horizontal dilation by $\frac{1}{2}$ because the distance to the $y$-axis has been halved. A vertical dilation, such as
\[
2 \times f(x)
\]
is slightly more straightforward. In this case, you double the output of the function. The output represents the distance from the $x$-axis, so in effect, you have made the graph of the function ''taller''. Here are a few basic examples where $a$ is any positive constant:

%-------------------------------------------------------------------------------
\subsection{Transformations}
\label{sss:tran-dfc3}
%-------------------------------------------------------------------------------

\begin{itemize}
    \item \textbf{Original graph}: $f(x)$
    \item \textbf{Rotation about origin}: $-f(-x)$
    \item \textbf{Horizontal translation by $a$ units left}: $f(x + a)$
    \item \textbf{Horizontal translation by $a$ units right}: $f(x - a)$
    \item \textbf{Horizontal dilation by a factor of $a$}: $f\left(x \times \frac{1}{a}\right)$
    \item \textbf{Vertical dilation by a factor of $a$}: $a \times f(x)$
    \item \textbf{Vertical translation by $a$ units down}: $f(x) - a$
    \item \textbf{Vertical translation by $a$ units up}: $f(x) + a$
    \item \textbf{Reflection about the $x$-axis}: $-f(x)$
    \item \textbf{Reflection about the $y$-axis}: $f(-x)$
\end{itemize}

%-------------------------------------------------------------------------------
\subsection{Inverse Functions}
\label{sss:inve-f81e}
%-------------------------------------------------------------------------------

We call $g(x)$ the inverse function of $f(x)$ if, for all $x$:
\[
g(f(x)) = f(g(x)) = x
\]
A function $f(x)$ has an inverse function if and only if $f(x)$ is one-to-one. For example, the inverse of $f(x) = x + 2$ is $g(x) = x - 2$. The function $f(x) = \sqrt{1 - x^2}$ has no inverse because it is not injective. Similarly, the inverse functions of trigonometric functions have to undergo transformations and limitations to be considered valid functions.

%-------------------------------------------------------------------------------
\subsection{Notation}
\label{sss:nota-71f9}
%-------------------------------------------------------------------------------

The inverse function of $f$ is denoted as $f^{-1}(x)$. Thus, $f^{-1}(x)$ is defined as the function that follows this rule:
\[
f(f^{-1}(x)) = f^{-1}(f(x)) = x
\]
To determine $f^{-1}(x)$ when given a function $f$, substitute $f^{-1}(x)$ for $x$ and substitute $x$ for $f(x)$. Then solve for $f^{-1}(x)$, provided that it is also a function.

%-------------------------------------------------------------------------------
\subsection{Example}
\label{sss:exam-1a79}
%-------------------------------------------------------------------------------

Given $f(x) = 2x - 7$, find $f^{-1}(x)$.

Substitute $f^{-1}(x)$ for $x$ and $x$ for $f(x)$. Then solve for $f^{-1}(x)$:
\[
f(x) = 2x - 7 \Leftrightarrow x = 2[f^{-1}(x)] - 7 \Leftrightarrow x + 7 = 2[f^{-1}(x)] \Leftrightarrow \frac{x + 7}{2} = f^{-1}(x)
\]
To check your work, confirm that $f^{-1}(f(x)) = x$:
\[
f^{-1}(f(x)) = f^{-1}(2x - 7) = \frac{(2x - 7) + 7}{2} = \frac{2x}{2} = x
\]
If $f$ isn''t one-to-one, then, as we said before, it doesn''t have an inverse. Then this method will fail.

%-------------------------------------------------------------------------------
\subsection{Example}
\label{sss:exam-1a79}
%-------------------------------------------------------------------------------

Given $f(x) = x^2$, find $f^{-1}(x)$.

Substitute $f^{-1}(x)$ for $x$ and $x$ for $f(x)$. Then solve for $f^{-1}(x)$:
\[
f(x) = x^2 \Leftrightarrow x = (f^{-1}(x))^2 \Leftrightarrow f^{-1}(x) = \pm \sqrt{x}
\]
Since there are two possibilities for $f^{-1}(x)$, it''s not a function. Thus $f(x) = x^2$ doesn''t have an inverse. Of course, we could also have found this out from the graph by applying the Horizontal Line Test. It''s useful to have lots of ways to solve a problem, since in a specific case some of them might be very difficult while others might be easy. For example, we might only know an algebraic expression for $f(x)$ but not a graph.

%-------------------------------------------------------------------------------
\subsection{Resources}
\label{sss:reso-55b5}
%-------------------------------------------------------------------------------

\begin{itemize}
    \item \href{https://www.khanacademy.org/math/algebra/x2f8bb11595b61c86:functions/x2f8bb11595b61c86:evaluating-functions/v/what-is-a-function}{What is a Function?, Khan Academy}
    \item \href{https://www.khanacademy.org/math/algebra2/functions_and_graphs}{Function and Their Graphs, Khan Academy}
    \item \href{https://www.cliffsnotes.com/study-guides/calculus/precalculus/functions/relations-vs-functions}{Relations vs. Functions, Cliff''s Notes}
    \item \href{https://courses.lumenlearning.com/ivytech-collegealgebra/chapter/determine-whether-a-relation-represents-a-function/}{Determining if a Relation is a Function, Lumen Learning}
    \item \href{https://courses.lumenlearning.com/waymakercollegealgebra/chapter/identify-functions-using-graphs/}{Identifying Functions Using Graphs, Lumen Learning}
\end{itemize}

%-------------------------------------------------------------------------------
\subsection{Licensing}
\label{sss:lice-7352}
%-------------------------------------------------------------------------------

Content obtained and/or adapted from:

\begin{itemize}
    \item \href{https://en.wikibooks.org/wiki/Calculus/Functions}{Functions, Wikibooks: Calculus} under a CC BY-SA license
\end{itemize}

%-------------------------------------------------------------------------------
\subsection{Learning Objectives}
\label{ssec:lear-4795}
%-------------------------------------------------------------------------------

\begin{enumerate}
    \item Define and differentiate between domain, range, and codomain of a function.
    \item Apply vertical and horizontal line tests to determine if a relation represents a function or if it is injective.
    \item Evaluate functions using their formulas and solve problems involving function mappings, including real-valued and vector-valued functions. Quadratic Functions 1.\textbf{Identify and Describe Quadratic Functions:}
    \item Students will be able to identify the general form of a quadratic function $f(x) = ax^2 + bx + c$, describe its properties (such as vertex, axis of symmetry, and direction of opening), and explain how changes in $a$, $b$, and $c$ affect the graph of the function. 2.\textbf{Solve Quadratic Equations Using the Quadratic Formula:}
    \item Students will be able to apply the quadratic formula $ x = \frac{-b \pm \sqrt{b^2 - 4ac}}{2a} $ to solve quadratic equations and interpret the solutions in the context of real-world problems. 3.\textbf{Graph Quadratic Functions and Analyze Their Properties:}
    \item Students will be able to graph quadratic functions, identify key features such as the vertex, axis of symmetry, and x-intercepts, and analyze how transformations (translations, dilations) affect the graph. Exponential and Logarithmic Functions 1.\textbf{Understand and Apply Exponential Functions:}
    \item Students will be able to define exponential functions $ f(x) = a^x $, understand the impact of different bases and coefficients on the function's growth, and solve problems involving exponential growth and decay. 2.\textbf{Understand and Apply Logarithmic Functions:}
    \item Students will be able to define logarithmic functions $ f(x) = \log_b x $, apply properties of logarithms to simplify expressions, and solve logarithmic equations. 3.\textbf{Utilize Euler’s Constant $ e $ in Exponential and Logarithmic Functions:}
    \item Students will be able to use Euler's constant $ e $ to solve problems involving natural exponential functions $ f(x) = e^x $ and natural logarithms $ f(x) = \ln(x) $, and understand its significance in continuous growth and decay processes. Signum Function 1.\textbf{Evaluate the Signum Function:}
    \item Students will be able to evaluate the signum function $ \text{sgn}(x) $ for given values of $ x $ and interpret its meaning in different contexts. 2.\textbf{Apply the Signum Function in Real-World Scenarios:}
    \item Students will be able to apply the signum function to model and solve problems involving different types of inputs and outputs, including scenarios involving directional or categorical data. Properties of Functions 1.\textbf{Determine the Domain and Range of Functions:}
    \item Students will be able to determine the domain and range of various functions, including rational, polynomial, and piecewise functions, and explain how these properties affect the behavior of the functions. 2.\textbf{Analyze and Describe Function Growth:}
    \item Students will be able to analyze whether a function is increasing or decreasing over specified intervals and use this information to describe the function's behavior in real-world contexts. 3.\textbf{Identify and Apply Function Symmetry:}
    \item Students will be able to determine if a function is even, odd, or neither by analyzing its symmetry about the y-axis and the origin, and apply these properties to solve related problems.
\end{enumerate}

%-------------------------------------------------------------------------------
\subsection{Assessment Instrument}
\label{ssec:asse-1481}
%-------------------------------------------------------------------------------

Assess through problems that require students to determine whether a given relation is a
function, evaluate a function at specific inputs, and interpret function behavior from
graphs and tables. Include a problem that distinguishes domain from range and one that
applies the vertical line test.

\textbf{Example problems.}

\begin{enumerate}
    \item Determine whether $\{(1,2),\,(2,3),\,(2,4),\,(3,5)\}$ is a function. Explain your reasoning.
    \item Given $f(x) = 3x^2 - x + 4$, evaluate $f(-2)$ and $f(a + 1)$.
    \item Apply the vertical line test to the circle $x^2 + y^2 = 9$ and state whether it defines $y$ as a function of $x$.
\end{enumerate}

%===============================================================================
\section{Function Notation}
\label{sec:func-4102}
%===============================================================================


\textbf{See Functions for more information.}

Functions are used so much that there is a special notation for them. The notation is
somewhat ambiguous, so familiarity with it is important in order to understand the
intention of an equation or formula.

Though there are no strict rules for naming a function, it is standard practice to use the letters $f$, $g$, and $h$ to denote functions, and the variable $x$ to denote an independent variable. $y$ is used for both dependent and independent variables.

When discussing or working with a function $f$, it's important to know not only the function, but also its independent variable $x$. Thus, when referring to a function $f$, you usually do not write $f$, but instead $f(x)$. The function is now referred to as ``$f$ of $x$''. The name of the function is adjacent to the independent variable (in parentheses). This is useful for indicating the value of the function at a particular value of the independent variable. For instance, if
\[
f(x) = 7x + 1
\]
and if we want to use the value of $f$ for $x$ equal to $2$, then we would substitute $2$ for $x$ on both sides of the definition above and write
\[
f(2) = 7(2) + 1 = 14 + 1 = 15
\]
This notation is more informative than leaving off the independent variable and writing simply ''$f$'', but can be ambiguous since the parentheses next to $f$ can be misinterpreted as multiplication, $2f$. To make sure nobody is too confused, follow this procedure:

1. Define the function $f$ by equating it to some expression.

2. Give a sentence like the following: ``At $x = c$, the function $f$ is...''

3. Evaluate the function.

%-------------------------------------------------------------------------------
\subsection{Function Description}
\label{sss:func-987a}
%-------------------------------------------------------------------------------

There are many ways which people describe functions. In the examples above, a verbal
description is given (the height of the ball above the earth as a function of time). Here
is a list of ways to describe functions. The top three listed approaches to describing
functions are the most popular.

1. A function is given a name (such as $f$) and a formula for the function is also given. For example, $f(x) = 3x + 2$ describes a function. We refer to the input as the argument of the function (or the independent variable), and to the output as the value of the function at the given argument.

2. A function is described using an equation and two variables. One variable is for the input of the function and one is for the output of the function. The variable for the input is called the independent variable. The variable for the output is called the dependent variable. For example, $y = 3x + 2$ describes a function. The dependent variable appears by itself on the left hand side of the equal sign.

3. A verbal description of the function.

When a function is given a name (like in number 1 above), the name of the function is usually a single letter of the alphabet (such as $f$ or $g$). Some functions whose names are multiple letters (like the sine function $y = \sin(x)$).

%-------------------------------------------------------------------------------
\subsection{Plugging a value into a function}
\label{sss:plug-50c6}
%-------------------------------------------------------------------------------

If we write $f(x) = 3x + 2$, then we know that

\begin{itemize}
    \item The function $f$ is a function of $x$.
    \item To evaluate the function at a certain number, replace the $x$ with that number.
    \item Replacing $x$ with that number in the right side of the function will produce the function's output for that certain input.
    \item In English, the definition of $f$ is interpreted, ``Given a number, $f$ will return two more than the triple of that number.''
\end{itemize}

How would we know the value of the function $f$ at $3$? We would have the following three thoughts:
\[
f(3) = 3(3) + 2
\]

\[
3(3) + 2 = 9 + 2
\]

\[
9 + 2 = 11
\]
and we would write
\[
f(3) = 3(3) + 2 = 9 + 2 = 11
\]
The value of $f$ at $3$ is $11$.

Note that $f(3)$ means the value of the dependent variable when $x$ takes on the value of $3$. So we see that the number $11$ is the output of the function when we give the number $3$ as the input. People often summarize the work above by writing ``the value of $f$ at three is eleven'', or simply ``$f$ of three equals eleven''.

%-------------------------------------------------------------------------------
\subsection{Resources}
\label{sss:reso-55b5}
%-------------------------------------------------------------------------------

\begin{itemize}
    \item \href{https://mathresearch.utsa.edu/wikiFiles/MAT1053/Function_and_Function_Notation/MAT1053_M1.1Functions_and_Function_Notation.pdf}{Function and Function Notation}, Book Chapter
    \item \href{https://mathresearch.utsa.edu/wikiFiles/MAT1053/Function_and_Function_Notation/MAT1053_M1.1_Functions_and_Function_NotationGN.pdf}{Function and Function Notation} Guided Notes
\end{itemize}

%-------------------------------------------------------------------------------
\subsection{Licensing}
\label{sss:lice-7352}
%-------------------------------------------------------------------------------

Content obtained and/or adapted from:

\begin{itemize}
    \item \href{https://en.wikibooks.org/wiki/Calculus/Functions}{Functions, Wikibooks: Calculus} under a CC BY-SA license
\end{itemize}

%-------------------------------------------------------------------------------
\subsection{Learning Objectives}
\label{ssec:lear-2f05}
%-------------------------------------------------------------------------------

\begin{enumerate}
    \item \textbf{1. Understand Function Notation} Students will be able to identify and explain the notation used for functions, including the use of letters such as $f$, $g$, and $h$ and the variables $x$ and $y$. \textbf{2. Evaluate Functions} Students will be able to evaluate a function $f(x)$ at a given value of $x$ by substituting the value into the function''s expression and simplifying to find the output. \textbf{3. Describe Functions} Students will be able to describe a function using different methods, including verbal descriptions, equations with independent and dependent variables, and by evaluating functions at specific points.
\end{enumerate}

%-------------------------------------------------------------------------------
\subsection{Assessment Instrument}
\label{ssec:asse-80b4}
%-------------------------------------------------------------------------------

Assess through exercises requiring students to translate between verbal descriptions,
algebraic rules, and function notation, and to evaluate composed or nested function
expressions. Problems should require correct substitution of non-numeric inputs such as $x
+ h$.

\textbf{Example problems.}

\begin{enumerate}
    \item Let $g(x) = \dfrac{x}{x-1}$. Compute $g(3)$, $g(0)$, and $g(x+1)$, simplifying each result.
    \item Rewrite ``the output is three more than twice the input'' in function notation, then evaluate the function at $x = 5$.
    \item Given $h(x) = x^2 + 2x$, compute and simplify $\dfrac{h(x+k) - h(x)}{k}$.
\end{enumerate}

%===============================================================================
\section{Domain of a Function}
\label{sec:doma-23ec}
%===============================================================================
%-------------------------------------------------------------------------------
\subsection{Definition}
\label{sss:defi-3061}
%-------------------------------------------------------------------------------

In mathematics, the domain or set of departure of a function is the set into which all of the input of the function is constrained to fall. It is the set $X$ in the notation $f: X \rightarrow Y$, and is alternatively denoted as $\text{dom}(f)$. Since a (total) function is defined on its entire domain, its domain coincides with its domain of definition. However, this coincidence is no longer true for a partial function, since the domain of definition of a partial function can be a proper subset of the domain.

A domain is part of a function $f$ if $f$ is defined as a triple $(X, Y, G)$, where $X$ is called the domain of $f$, $Y$ its codomain, and $G$ its graph.

A domain is not part of a function $f$ if $f$ is defined as just a graph. For example, it is sometimes convenient in set theory to permit the domain of a function to be a proper class $X$, in which case there is formally no such thing as a triple $(X, Y, G)$. With such a definition, functions do not have a domain, although some authors still use it informally after introducing a function in the form $f: X \rightarrow Y$.

For instance, the domain of cosine is the set of all real numbers, while the domain of the square root consists only of numbers greater than or equal to $0$ (ignoring complex numbers in both cases).

If the domain of a function is a subset of the real numbers and the function is represented in a Cartesian coordinate system, then the domain is represented on the $x$-axis. The domain of a function $f$ can be thought of as the set of all $x$ values that can be plugged into $f(x)$ that return a valid output. For example, if we have a function $g(x)$ in the Cartesian plane, the domain is all of the $x$ values such that $g(x)$ is a real number.

%-------------------------------------------------------------------------------
\subsection{Examples:}
\label{sss:exam-bfeb}
%-------------------------------------------------------------------------------

\begin{itemize}
    \item Let $S$ be a set of ordered pairs such that $S = \\{(1, 2), (2, 3), (4, 7), (13, 9), (-20, 0)\\}$.
\end{itemize}

\begin{itemize}
    \item The domain is the set of all $x$ values of $S$, so the domain is $\\{-20, 1, 2, 4, 13\\}$.
\end{itemize}

\begin{itemize}
    \item The domain of $g(x) = \frac{1}{x}$ is all real numbers \textbf{EXCEPT} $0$, since $\frac{1}{0}$ is not defined.
\end{itemize}

\begin{itemize}
    \item The domain of $h(x) = \sqrt{x}$ is $[0, \infty)$, since $\sqrt{x}$ is only defined when $x$ is nonnegative (that is, when $x$ is greater than or equal to $0$).
\end{itemize}

%-------------------------------------------------------------------------------
\subsection{Resources}
\label{sss:reso-55b5}
%-------------------------------------------------------------------------------

\begin{itemize}
    \item \href{https://www.youtube.com/watch?v=Q3NWljhiSJg}{Domain and Range: Basic Idea, patrickJMT}
    \item \href{https://courses.lumenlearning.com/ivytech-collegealgebra/chapter/find-the-domain-of-a-function-defined-by-an-equation/}{Finding Domain with an Equation, Lumen Learning}
    \item \href{https://courses.lumenlearning.com/ivytech-collegealgebra/chapter/find-domain-and-range-from-graphs/}{Finding Domain and Range with Graphs, Lumen Learning}
    \item \href{https://www.youtube.com/watch?v=w81y25anEOM}{Finding the Domain Algebraically, patrickJMT}
    \item \href{https://www.youtube.com/watch?v=BxaYyS6lsQ4}{Finding Domain and Range of a Piecewise Function, patrickJMT}
    \item \href{https://mathculus.com/how-to-find-the-domain-of-a-function-algebraically/}{How to Find Domain + Example Problems, Math Culus}
\end{itemize}

%-------------------------------------------------------------------------------
\subsection{Licensing}
\label{sss:lice-7352}
%-------------------------------------------------------------------------------

Content obtained and/or adapted from:

\begin{itemize}
    \item \href{https://en.wikipedia.org/wiki/Domain_of_a_function}{Domain of a Function, Wikipedia} under a CC BY-SA license
\end{itemize}

%-------------------------------------------------------------------------------
\subsection{Learning Objectives}
\label{ssec:lear-ccc8}
%-------------------------------------------------------------------------------

\begin{enumerate}
    \item \textbf{Identify the domain of a function:} Students will be able to define the domain of a function and explain its significance in terms of the input values that a function can accept.
    \item \textbf{Determine the domain of a function from its equation:} Students will be able to determine the domain of various types of functions, including polynomial, rational, and square root functions, based on their equations.
    \item \textbf{Analyze the domain of a function represented graphically:} Students will be able to identify the domain of a function when it is represented in a Cartesian coordinate system.
\end{enumerate}

%-------------------------------------------------------------------------------
\subsection{Assessment Instrument}
\label{ssec:asse-a1fc}
%-------------------------------------------------------------------------------

Assess through problems requiring students to determine the domain of rational and radical
functions algebraically (denominator $\neq 0$ and radicand $\geq 0$). At least one problem
should ask students to express the domain in both interval notation and set-builder
notation.

\textbf{Example problems.}

\begin{enumerate}
    \item Find the domain of $f(x) = \dfrac{3x + 1}{x^2 - 4}$ and write it in interval notation.
    \item Determine the domain of $g(x) = \sqrt{5 - 2x}$ and write it in set-builder notation.
    \item Identify a real-world scenario whose natural domain is restricted to positive integers and write a corresponding function.
\end{enumerate}

%===============================================================================
\section{Range of a Function}
\label{sec:rang-baf4}
%===============================================================================
%-------------------------------------------------------------------------------
\subsection{Definition}
\label{sss:defi-3061}
%-------------------------------------------------------------------------------

In mathematics, the range of a function may refer to either of two closely related
concepts:

\begin{itemize}
    \item The codomain of the function
    \item The image of the function
\end{itemize}

Given two sets $X$ and $Y$, a binary relation $f$ between $X$ and $Y$ is a (total) function (from $X$ to $Y$) if for every $x$ in $X$ there is exactly one $y$ in $Y$ such that $f$ relates $x$ to $y$. The sets $X$ and $Y$ are called domain and codomain of $f$, respectively. The image of $f$ is then the subset of $Y$ consisting of only those elements $y$ of $Y$ such that there is at least one $x$ in $X$ with $f(x) = y$.

In algebra, the range (or codomain) of a function is all of the possible outputs of the function. That is, if $x$ is any element of the domain of some function $f$, then $f(x)$ is in the range of the function $f$.

%-------------------------------------------------------------------------------
\subsection{Examples:}
\label{sss:exam-bfeb}
%-------------------------------------------------------------------------------

\begin{itemize}
    \item Let $S$ be a set of ordered pairs such that $S = \\{(1, 2), (2, 3), (4, 7), (13, 9), (-20, 0)\\}$. The range is the set of all $y$ values of $S$, so the range is $\{0, 2, 3, 7, 9\}$.
\end{itemize}

\begin{itemize}
    \item The range of $g(x) = \frac{1}{x}$ is all real numbers \textbf{EXCEPT} for $0$. We know this because for all nonzero real numbers $M$, $\frac{1}{M}$ is a nonzero number and is in the domain of $g(x)$ (since the domain of this function is all nonzero numbers). So, we know that $\frac{1}{\left(\frac{1}{M}\right)} = M$ is in the range, where $M$ is all nonzero numbers. There is no real number $M$ such that $\frac{1}{M} = 0$ though, which is why $0$ is not in the range of $g(x)$.
\end{itemize}

\begin{itemize}
    \item The range of $h(x) = x^2 + 2$ is $[2, \infty)$. We can see this on the graph of $h(x)$ easily: the lowest point, or vertex, of the parabola is at $(0, 2)$, so $2$ is in the range. The parabola extends up to infinity on either side of the vertex, so we know that the range must be all numbers from $2$ to infinity.
\end{itemize}

%-------------------------------------------------------------------------------
\subsection{Resources}
\label{sss:reso-55b5}
%-------------------------------------------------------------------------------

\begin{itemize}
    \item \href{https://www.intmath.com/functions-and-graphs/2a-domain-and-range.php}{Domain and Range, Interactive Mathematics}
    \item \href{https://www.youtube.com/watch?v=Q3NWljhiSJg}{Domain and Range: Basic Idea, patrickJMT}
    \item \href{https://courses.lumenlearning.com/ivytech-collegealgebra/chapter/find-domain-and-range-from-graphs/}{Finding Domain and Range with Graphs, Lumen Learning}
    \item \href{https://www.youtube.com/watch?v=BxaYyS6lsQ4}{Finding Domain and Range of a Piecewise Function, patrickJMT}
    \item \href{https://mathculus.com/how-to-find-the-domain-and-range-of-a-function-algebraically/}{How to Find Range + Example Problems, Math Culus}
\end{itemize}

%-------------------------------------------------------------------------------
\subsection{Licensing}
\label{sss:lice-7352}
%-------------------------------------------------------------------------------

Content obtained and/or adapted from:

\begin{itemize}
    \item \href{https://en.wikipedia.org/wiki/Range_(mathematics}{Range of a function, Wikipedia}) under a CC BY-SA license
\end{itemize}

%-------------------------------------------------------------------------------
\subsection{Learning Objectives}
\label{ssec:lear-0b2d}
%-------------------------------------------------------------------------------

\begin{enumerate}
    \item \textbf{Understand the Range of a Function}: Students will be able to identify and explain the range of a function, distinguishing between the codomain and the image.
    \item \textbf{Determine the Range for Given Functions}: Students will be able to determine the range of a function given its equation, focusing on various types of functions (e.g., rational, quadratic).
    \item \textbf{Analyze Graphs to Find the Range}: Students will be able to analyze and interpret graphs of functions to accurately identify the range, recognizing key features such as vertices and asymptotes.
\end{enumerate}

%-------------------------------------------------------------------------------
\subsection{Assessment Instrument}
\label{ssec:asse-7e68}
%-------------------------------------------------------------------------------

Assess through problems that require students to determine the range of absolute-value,
quadratic, and constant functions by analyzing behavior, and to explain why the range of a
function may differ from its codomain.

\textbf{Example problems.}

\begin{enumerate}
    \item State the range of $f(x) = |x - 3| + 2$ and explain why no output can be below $2$.
    \item Determine the range of $g(x) = -x^2 + 4$ by identifying the vertex.
    \item A function maps every real number to its square. State the domain, range, and codomain if the codomain is $\mathbb{R}$.
\end{enumerate}

%===============================================================================
\section{Toolkit Functions}
\label{sec:tool-4c3a}
%===============================================================================
%-------------------------------------------------------------------------------
\subsection{Constant Function}
\label{sss:cons-6354}
%-------------------------------------------------------------------------------

For the constant function $f(x) = c$, the domain consists of all real numbers; there are no restrictions on the input. The only output value is the constant $c$, so the range is the set $\{c\}$ that contains this single element. In interval notation, this is written as $[c, c]$, the interval that both begins and ends with $c$.

% [Image: constant]

\begin{itemize}
    \item \textbf{Domain:} $(-\infty ,\infty)$
    \item \textbf{Range:} $[c, c]$
\end{itemize}

%-------------------------------------------------------------------------------
\subsection{Identity Function}
\label{sss:iden-48e9}
%-------------------------------------------------------------------------------

For the identity function $f(x) = x$, there is no restriction on $x$. Both the domain and range are the set of all real numbers.

% [Image: identity]

\begin{itemize}
    \item \textbf{Domain:} $(-\infty ,\infty)$
    \item \textbf{Range:} $(-\infty ,\infty)$
\end{itemize}

%-------------------------------------------------------------------------------
\subsection{Absolute Value Function}
\label{sss:abso-ffeb}
%-------------------------------------------------------------------------------

For the absolute value function $f(x) = |x|$, there is no restriction on $x$. The outputs are $x$ for $x \geq 0$ and $-x$ for $x < 0$, so the range is all numbers greater than or equal to 0.

% [Image: absolute]

\begin{itemize}
    \item \textbf{Domain:} $(-\infty ,\infty)$
    \item \textbf{Range:} $[0 ,\infty)$
\end{itemize}

%-------------------------------------------------------------------------------
\subsection{Quadratic Function}
\label{sss:quad-a21f}
%-------------------------------------------------------------------------------

For the quadratic function $f(x) = x^2$, the domain is all real numbers since the horizontal extent of the graph is the whole real number line. Because the graph does not include any negative values for the range, the range is only nonnegative real numbers.

% [Image: quad]

\begin{itemize}
    \item \textbf{Domain:} $(-\infty ,\infty)$
    \item \textbf{Range:} $[0 ,\infty)$
\end{itemize}

%-------------------------------------------------------------------------------
\subsection{Cubic Function}
\label{sss:cubi-6dd3}
%-------------------------------------------------------------------------------

For the cubic function $f(x) = x^3$, the domain is all real numbers because the horizontal extent of the graph is the whole real number line. The same applies to the vertical extent of the graph, so the domain and range include all real numbers.

% [Image: cubic]

\begin{itemize}
    \item \textbf{Domain:} $(-\infty ,\infty)$
    \item \textbf{Range:} $(-\infty ,\infty)$
\end{itemize}

%-------------------------------------------------------------------------------
\subsection{Rational Function}
\label{sss:rati-42db}
%-------------------------------------------------------------------------------

For the rational function (also known as the reciprocal function) $f(x) = \frac{1}{x}$, we cannot divide by 0, so we must exclude 0 from the domain. Further, 1 divided by any value can never be 0, so the range also will not include 0.

% [Image: rational]

\begin{itemize}
    \item \textbf{Domain:} $(-\infty ,0) \cup (0 ,\infty)$
    \item \textbf{Range:} $(-\infty ,0) \cup (0 ,\infty)$
\end{itemize}

%-------------------------------------------------------------------------------
\subsection{Squared Rational Function}
\label{sss:squa-de87}
%-------------------------------------------------------------------------------

For the squared rational function $f(x) = \frac{1}{x^2}$, we cannot divide by 0, so we must exclude 0 from the domain. Further, 1 divided by any value can never be 0, so the range also will not include 0. Also, since $x^2 > 0$ for all $x \neq 0$, the range will only consist of positive numbers.

% [Image: sqaured rat]

\begin{itemize}
    \item \textbf{Domain:} $(-\infty ,0) \cup (0 ,\infty)$
    \item \textbf{Range:} $(0 ,\infty)$
\end{itemize}

%-------------------------------------------------------------------------------
\subsection{Square Root Function}
\label{sss:squa-a510}
%-------------------------------------------------------------------------------

For the square root function $f(x) = \sqrt{x}$, we cannot take the square root of a negative real number, so the domain must be 0 or greater. The range also excludes negative numbers because the square root of a positive number $x$ is defined to be positive, even though the square of the negative number $-\sqrt{x}$ also gives us $x$.

% [Image: sq root]

\begin{itemize}
    \item \textbf{Domain:} $[0 ,\infty)$
    \item \textbf{Range:} $[0 ,\infty)$
\end{itemize}

%-------------------------------------------------------------------------------
\subsection{Cube Root Function}
\label{sss:cube-28dd}
%-------------------------------------------------------------------------------

For the cube root function $f(x) = \sqrt[3]{x}$, the domain and range include all real numbers. Note that there is no problem taking a cube root, or any odd-integer root, of a negative number, and the resulting output is negative (it is an odd function).

% [Image: cubic root]

\begin{itemize}
    \item \textbf{Domain:} $(-\infty ,\infty)$
    \item \textbf{Range:} $(-\infty ,\infty)$
\end{itemize}

%-------------------------------------------------------------------------------
\subsection{Resources}
\label{sss:reso-55b5}
%-------------------------------------------------------------------------------

\begin{itemize}
    \item \href{https://mathresearch.utsa.edu/wikiFiles/MAT1053/Domain_Range_and_Toolkit_Functions/MAT1053_M1.2Domain_Range_and_Toolkit_Functions.pdf}{Domain Range and Toolkit Functions, Book Chapter}
    \item \href{https://mathresearch.utsa.edu/wikiFiles/MAT1053/Domain_Range_and_Toolkit_Functions/MAT1053_M1.2Domain_Range_and_Toolkit_FunctionsGN.pdf}{Guided Notes}
    \item \href{https://courses.lumenlearning.com/cuny-hunter-collegealgebra/chapter/toolkit-functions/}{Toolkit Functions, Lumen Learning}
\end{itemize}

%-------------------------------------------------------------------------------
\subsection{Licensing}
\label{sss:lice-7352}
%-------------------------------------------------------------------------------

Content obtained and/or adapted from:

\begin{itemize}
    \item \href{https://courses.lumenlearning.com/cuny-hunter-collegealgebra/chapter/toolkit-functions/}{Toolkit Functions, Lumen Learning: Hunter College MATH101} under a CC0 (public domain) license.
\end{itemize}

%-------------------------------------------------------------------------------
\subsection{Learning Objectives}
\label{ssec:lear-68e5}
%-------------------------------------------------------------------------------

\begin{enumerate}
    \item \textbf{1. Identify and describe the domain and range of constant, identity, and absolute value functions.}
    \item Students will be able to explain the characteristics of the domain and range for the constant function $f(x) = c$, identity function $f(x) = x$, and absolute value function $f(x) = |x|$. \textbf{2. Graph constant, identity, and absolute value functions on the Cartesian plane.}
    \item Students will accurately graph the constant function $f(x) = c$, identity function $f(x) = x$, and absolute value function $f(x) = |x|$ on the Cartesian plane, identifying key points and intervals. \textbf{3. Compare the behavior of quadratic and cubic functions.}
    \item Students will compare and contrast the domain, range, and graphical behavior of the quadratic function $f(x) = x^2$ and cubic function $f(x) = x^3$.
\end{enumerate}

%-------------------------------------------------------------------------------
\subsection{Assessment Instrument}
\label{ssec:asse-3d47}
%-------------------------------------------------------------------------------

Assess through problems requiring students to identify the name, algebraic rule, domain,
range, and key graph features of each toolkit function. Include a matching problem
(equations to graphs) and a piecewise evaluation problem.

\textbf{Example problems.}

\begin{enumerate}
    \item Sketch and label $f(x) = x$, $f(x) = x^2$, and $f(x) = |x|$ on the same axes, stating the domain and range of each.
    \item Evaluate $f(x) = \begin{cases} x^2 & x < 0 \\ \sqrt{x} & x \geq 0 \end{cases}$ at $x = -3$, $x = 0$, and $x = 4$.
    \item Compare the domain and range of $f(x) = x^3$ and $g(x) = \sqrt[3]{x}$ and determine whether they are inverses of each other.
\end{enumerate}

%===============================================================================
\section{Composition of Functions}
\label{sec:comp-37f3}
%===============================================================================
In mathematics, function composition is an operation that takes two functions $f$ and $g$ and produces a function $h$ such that $h(x) = g(f(x))$. In this operation, the function $g$ is applied to the result of applying the function $f$ to $x$. That is, the functions $f : X \to Y$ and $g : Y \to Z$ are composed to yield a function that maps $x$ in $X$ to $g(f(x))$ in $Z$.

Intuitively, if $z$ is a function of $y$, and $y$ is a function of $x$, then $z$ is a function of $x$. The resulting composite function is denoted $g \circ f : X \to Z$, defined by $(g \circ f)(x) = g(f(x))$ for all $x$ in $X$. The notation $g \circ f$ is read as ``g circle f'', ``g round f'', ``g about f'', ``g composed with f'', ``g after f'', ``g following f'', ``g of f'', ``f then g'', or ``g on f'', or ``the composition of $g$ and $f$''. Intuitively, composing functions is a chaining process in which the output of function $f$ feeds the input of function $g$.

The composition of functions is a special case of the composition of relations, sometimes also denoted by $\circ$. As a result, all properties of composition of relations are true of composition of functions, though the composition of functions has some additional properties.

Composition of functions is different from multiplication of functions and has quite
different properties; in particular, composition of functions is not commutative.

%-------------------------------------------------------------------------------
\subsection{Examples}
\label{sss:exam-bfeb}
%-------------------------------------------------------------------------------

1. \textbf{Concrete example for the composition of two functions.}

% [Image: example]

\begin{itemize}
    \item Composition of functions on a finite set: If $f = \{(1, 1), (2, 3), (3, 1), (4, 2)\}$, and $g = \{(1, 2), (2, 3), (3, 1), (4, 2)\}$, then $g \circ f = \{(1, 2), (2, 1), (3, 2), (4, 3)\}$, as shown in the figure.
\end{itemize}

2. \textbf{Composition of functions on an infinite set:}

\begin{itemize}
    \item If $f : \mathbb{R} \to \mathbb{R}$ (where $\mathbb{R}$ is the set of all real numbers) is given by $f(x) = 2x + 4$ and $g : \mathbb{R} \to \mathbb{R}$ is given by $g(x) = x^3$, then:
    \item $(f \circ g)(x) = f(g(x)) = f(x^3) = 2x^3 + 4$, and
    \item $(g \circ f)(x) = g(f(x)) = g(2x + 4) = (2x + 4)^3$.
    \item If an airplane''s altitude at time $t$ is $a(t)$, and the air pressure at altitude $x$ is $p(x)$, then $(p \circ a)(t)$ is the pressure around the plane at time $t$.
\end{itemize}

%-------------------------------------------------------------------------------
\subsection{Properties}
\label{sss:prop-7469}
%-------------------------------------------------------------------------------

\begin{itemize}
    \item The composition of functions is always associative: a property inherited from the composition of relations. That is, if $f$, $g$, and $h$ are composable, then $f \circ (g \circ h) = (f \circ g) \circ h$. Since the parentheses do not change the result, they are generally omitted.
\end{itemize}

\begin{itemize}
    \item In a strict sense, the composition $g \circ f$ is only meaningful if the codomain of $f$ equals the domain of $g$; in a wider sense, it is sufficient that the former be a subset of the latter. Moreover, it is often convenient to tacitly restrict the domain of $f$, such that $f$ produces only values in the domain of $g$. For example, the composition $g \circ f$ of the functions $f : \mathbb{R} \to \mathbb{R}$ defined by $f(x) = 9 - x^2$ and $g : \mathbb{R} \to \mathbb{R}$ defined by $g(x) = \sqrt{x}$ can be defined on the interval $[-3, +3]$.
\end{itemize}

\begin{itemize}
    \item The functions $g$ and $f$ are said to commute with each other if $g \circ f = f \circ g$. Commutativity is a special property, attained only by particular functions, and often in special circumstances. For example, $|x| + 3 = |x + 3|$ only when $x \geq 0$. The picture shows another example.
\end{itemize}

% [Image: graph]

\begin{itemize}
    \item The composition of one-to-one (injective) functions is always one-to-one. Similarly, the composition of onto (surjective) functions is always onto. It follows that the composition of two bijections is also a bijection. The inverse function of a composition (assumed invertible) has the property that $(f \circ g)^{-1} = g^{-1} \circ f^{-1}$.
\end{itemize}

%-------------------------------------------------------------------------------
\subsection{Resources}
\label{sss:reso-55b5}
%-------------------------------------------------------------------------------

\begin{itemize}
    \item \href{https://www.khanacademy.org/math/precalculus/x9e81a4f98389efdf:composite/x9e81a4f98389efdf:composing/v/function-composition}{Intro to Composing Functions, Khan Academy}
    \item \href{https://courses.lumenlearning.com/wmopen-collegealgebra/chapter/introduction-compositions-of-functions/}{Compositions of Functions, Lumen Learning}
    \item \href{https://www.mathsisfun.com/sets/functions-composition.html}{Composition of Functions, Math Is Fun}
\end{itemize}

%-------------------------------------------------------------------------------
\subsection{Licensing}
\label{sss:lice-7352}
%-------------------------------------------------------------------------------

Content obtained and/or adapted from:

\begin{itemize}
    \item \href{https://en.wikipedia.org/wiki/Function_composition}{Function composition, Wikipedia under a CC BY-SA license}
\end{itemize}

%-------------------------------------------------------------------------------
\subsection{Learning Objectives}
\label{ssec:lear-9e9f}
%-------------------------------------------------------------------------------

\begin{enumerate}
    \item \textbf{1. Understand and Apply Function Composition}:
    \item Students will be able to compose two functions $f$ and $g$ to form a new function $h$ such that $h(x) = g(f(x))$. \textbf{2. Identify and Interpret Composite Functions}:
    \item Students will be able to identify composite functions and interpret the chaining process of applying $f$ and $g$ in different orders. \textbf{3. Properties and Characteristics of Composite Functions}:
    \item Students will be able to explain the associative property of function composition, recognize non-commutativity, and understand the conditions under which functions commute.
\end{enumerate}

%-------------------------------------------------------------------------------
\subsection{Assessment Instrument}
\label{ssec:asse-c630}
%-------------------------------------------------------------------------------

Assess through problems requiring students to compute $(f \circ g)(x)$ and $(g \circ
f)(x)$, verify whether they are equal, determine the domain of the composition, and
decompose a given composite expression into inner and outer functions.

\textbf{Example problems.}

\begin{enumerate}
    \item Let $f(x) = 2x - 1$ and $g(x) = x^2 + 3$. Compute $(f \circ g)(x)$ and $(g \circ f)(x)$ and state whether the results are equal.
    \item Find the domain of $(f \circ g)(x)$ where $f(x) = \sqrt{x}$ and $g(x) = x - 5$.
    \item Express $h(x) = (3x + 2)^4$ as a composition $f(g(x))$, identifying $f$ and $g$ explicitly.
\end{enumerate}

%===============================================================================
\section{Inverse Functions}
\label{sec:inve-f81e}
%===============================================================================

In mathematics, an inverse function (or anti-function) is a function that ``reverses'' another function: if the function $f$ applied to an input $x$ gives a result of $y$, then applying its inverse function $g$ to $y$ gives the result $x$, i.e., $g(y) = x$ if and only if $f(x) = y$. The inverse function of $f$ is also denoted as $f^{-1}$.

As an example, consider the real-valued function of a real variable given by $f(x) = 5x - 7$. Thinking of this as a step-by-step procedure (namely, take a number $x$, multiply it by 5, then subtract 7 from the result), to reverse this and get $x$ back from some output value, say $y$, we would undo each step in reverse order. In this case, it means to add 7 to $y$, and then divide the result by 5. In functional notation, this inverse function would be given by,
\[
g(y) = \frac{y + 7}{5}.
\]
With $y = 5x - 7$ we have that $f(x) = y$ and $g(y) = x$.

Not all functions have inverse functions. Those that do are called invertible. For a function $f: X \to Y$ to have an inverse, it must have the property that for every $y$ in $Y$, there is exactly one $x$ in $X$ such that $f(x) = y$. This property ensures that a function $g: Y \to X$ exists with the necessary relationship with $f$.

%-------------------------------------------------------------------------------
\subsection{Definitions}
\label{sss:defi-9fdc}
%-------------------------------------------------------------------------------

If $f$ maps $X$ to $Y$, then $f^{-1}$ maps $Y$ back to $X$.

Let $f$ be a function whose domain is the set $X$, and whose codomain is the set $Y$. Then $f$ is invertible if there exists a function $g$ with domain $Y$ and codomain $X$, with the property:
\[
f(x) = y \Leftrightarrow g(y) = x.
\]
If $f$ is invertible, then the function $g$ is unique, which means that there is exactly one function $g$ satisfying this property. Moreover, it also follows that the ranges of $g$ and $f$ equal their respective codomains. The function $g$ is called the inverse of $f$, and is usually denoted as $f^{-1}$, a notation introduced by John Frederick William Herschel in 1813.

Stated otherwise, a function, considered as a binary relation, has an inverse if and only if the converse relation is a function on the codomain $Y$, in which case the converse relation is the inverse function. Not all functions have an inverse. For a function to have an inverse, each element $y \in Y$ must correspond to no more than one $x \in X$; a function $f$ with this property is called one-to-one or an injection. If $f^{-1}$ is to be a function on $Y$, then each element $y \in Y$ must correspond to some $x \in X$. Functions with this property are called surjections. This property is satisfied by definition if $Y$ is the image of $f$, but may not hold in a more general context. To be invertible, a function must be both an injection and a surjection. Such functions are called bijections. The inverse of an injection $f: X \to Y$ that is not a bijection (that is, not a surjection), is only a partial function on $Y$, which means that for some $y \in Y$, $f^{-1}(y)$ is undefined. If a function $f$ is invertible, then both it and its inverse function $f^{-1}$ are bijections.

% [Image: mapping]

Another convention is used in the definition of functions, referred to as the
``set-theoretic'' or ``graph'' definition using ordered pairs, which makes the codomain
and image of the function the same. Under this convention, all functions are surjective,
so bijectivity and injectivity are the same. Authors using this convention may use the
phrasing that a function is invertible if and only if it is an injection. The two
conventions need not cause confusion, as long as it is remembered that in this alternate
convention, the codomain of a function is always taken to be the image of the function.

%-------------------------------------------------------------------------------
\subsection{Properties}
\label{sss:prop-7469}
%-------------------------------------------------------------------------------

Since a function is a special type of binary relation, many of the properties of an
inverse function correspond to properties of converse relations.

%-------------------------------------------------------------------------------
\subsection{Uniqueness}
\label{sss:uniq-754c}
%-------------------------------------------------------------------------------

If an inverse function exists for a given function $f$, then it is unique. This follows since the inverse function must be the converse relation, which is completely determined by $f$.

%-------------------------------------------------------------------------------
\subsection{Symmetry}
\label{sss:symm-40e7}
%-------------------------------------------------------------------------------

There is a symmetry between a function and its inverse. Specifically, if $f$ is an invertible function with domain $X$ and codomain $Y$, then its inverse $f^{-1}$ has domain $Y$ and image $X$, and the inverse of $f^{-1}$ is the original function $f$. In symbols, for functions $f: X \to Y$ and $f^{-1}: Y \to X$,
\[
f^{-1} \circ f = \operatorname{id}_X
\]
and
\[
f \circ f^{-1} = \operatorname{id}_Y.
\]
This statement is a consequence of the implication that for $f$ to be invertible it must be bijective. The involutory nature of the inverse can be concisely expressed by
\[
\left(f^{-1}\right)^{-1} = f.
\]
The inverse of a composition of functions is given by
\[
(g \circ f)^{-1} = f^{-1} \circ g^{-1}.
\]
% [Image: inverse]

Notice that the order of $g$ and $f$ have been reversed; to undo $f$ followed by $g$, we must first undo $g$, and then undo $f$.

For example, let $f(x) = 3x$ and let $g(x) = x + 5$. Then the composition $g \circ f$ is the function that first multiplies by three and then adds five,
\[
(g \circ f)(x) = 3x + 5.
\]
To reverse this process, we must first subtract five, and then divide by three,
\[
(g \circ f)^{-1}(x) = \frac{1}{3}(x - 5).
\]
This is the composition $(f^{-1} \circ g^{-1})(x)$.

%-------------------------------------------------------------------------------
\subsection{Self-inverses}
\label{sss:self-4f9b}
%-------------------------------------------------------------------------------

If $X$ is a set, then the identity function on $X$ is its own inverse:
\[
\operatorname{id}_X^{-1} = \operatorname{id}_X.
\]
More generally, a function $f : X \to X$ is equal to its own inverse, if and only if the composition $f \circ f$ is equal to $\operatorname{id}_X$. Such a function is called an involution.

%-------------------------------------------------------------------------------
\subsection{Resources}
\label{sss:reso-55b5}
%-------------------------------------------------------------------------------

\begin{itemize}
    \item \href{https://courses.lumenlearning.com/collegealgebra2017/chapter/introduction-inverse-functions/}{Inverse Functions, Lumen Learning}
    \item \href{https://www.khanacademy.org/math/algebra/x2f8bb11595b61c86:functions/x2f8bb11595b61c86:inverse-functions-intro/v/introduction-to-function-inverses}{Introduction to Inverse Function Video, Khan Academy}
\end{itemize}

%-------------------------------------------------------------------------------
\subsection{Licensing}
\label{sss:lice-7352}
%-------------------------------------------------------------------------------

Content obtained and/or adapted from:

\begin{itemize}
    \item \href{https://en.wikipedia.org/wiki/Inverse_function}{Inverse function, Wikipedia under a CC BY-SA license}
\end{itemize}

%-------------------------------------------------------------------------------
\subsection{Learning Objectives}
\label{ssec:lear-eeb6}
%-------------------------------------------------------------------------------

\begin{enumerate}
    \item Identify and denote the inverse function $f^{-1}$ for a given function $f$ and understand the notation and concept of inverse functions.
    \item Calculate the inverse of a linear function by reversing the operations of the original function and express the inverse function in functional notation.
    \item Determine if a function is invertible by analyzing if it is bijective, ensuring it meets the criteria of being both one-to-one and onto.
\end{enumerate}

%-------------------------------------------------------------------------------
\subsection{Assessment Instrument}
\label{ssec:asse-a1b3}
%-------------------------------------------------------------------------------

Assess through problems requiring students to verify inverse function pairs using the
composition criterion $f(f^{-1}(x)) = x$, find the algebraic formula for an inverse, and
use the horizontal line test to determine whether a function is one-to-one.

\textbf{Example problems.}

\begin{enumerate}
    \item Find the inverse of $f(x) = \dfrac{2x - 3}{5}$ and verify using composition.
    \item Determine whether $g(x) = x^2$ on $\mathbb{R}$ has an inverse. If not, restrict the domain so that it does and state the resulting inverse function.
    \item Describe how to obtain the graph of $f^{-1}$ from the graph of $f$, and apply this procedure to $f(x) = 2^x$.
\end{enumerate}

%===============================================================================
\section{Exponential Properties}
\label{sec:expo-7ac7}
%===============================================================================
Exponential properties can be used to manipulate equations involving exponential
expressions and/or functions. Here are some important exponential properties:

\begin{itemize}
    \item \textbf{Negative exponent property}:
\end{itemize}
  \[
a^{-n} = \frac{1}{a^n} \quad \text{and} \quad \frac{1}{a^{-n}} = a^n
\]

\begin{itemize}
    \item \textbf{Product of like bases}:
\end{itemize}
  \[
a^m a^n = a^{m+n}
\]

\begin{itemize}
    \item \textbf{Quotient of like bases}:
\end{itemize}
  \[
\frac{a^m}{a^n} = a^m a^{-n} = a^{m-n}
\]

\begin{itemize}
    \item \textbf{Multiple powers}:
\end{itemize}
  \[
(a^m)^n = a^{mn} = (a^n)^m
\]

\begin{itemize}
    \item \textbf{Product to a power}:
\end{itemize}
  \[
(ab)^n = a^n b^n
\]

\begin{itemize}
    \item \textbf{Quotient to a power}:
\end{itemize}
  \[
\left(\frac{a}{b}\right)^n = \frac{a^n}{b^n}
\]

%-------------------------------------------------------------------------------
\subsection{Special cases involving 0:}
\label{sss:spec-4783}
%-------------------------------------------------------------------------------

\begin{itemize}
    \item For any nonzero number $a$, $a^0 = 1$.
    \item For any positive number $m$, $0^m = 0$.
    \item $0^m$ does not exist if $m$ is negative (since $0^{-n} = \frac{1}{0^n} = \frac{1}{0}$).
    \item $0^0$ is indeterminate or undefined depending on the context.
\end{itemize}

%-------------------------------------------------------------------------------
\subsection{Resources}
\label{sss:reso-55b5}
%-------------------------------------------------------------------------------

\begin{itemize}
    \item \href{https://tutoring.asu.edu/sites/default/files/exponentialandlogrithmicproperties.pdf}{Exponential and Logarithmic Properties, Arizona State University}
    \item \href{https://www.khanacademy.org/math/cc-eighth-grade-math/cc-8th-numbers-operations/cc-8th-exponent-properties/a/exponent-properties-review}{Exponent Properties Review, Khan Academy}
\end{itemize}

%-------------------------------------------------------------------------------
\subsection{Learning Objectives}
\label{ssec:lear-29e5}
%-------------------------------------------------------------------------------

\begin{enumerate}
    \item \textbf{Understanding Exponential Properties}: Students will identify and apply the properties of exponents in various mathematical contexts.
    \item \textbf{Manipulating Exponential Expressions}: Students will simplify and manipulate exponential expressions using the correct properties.
    \item \textbf{Special Cases with Exponents}: Students will recognize and correctly handle special cases involving exponents, including zero and negative exponents.
\end{enumerate}

%-------------------------------------------------------------------------------
\subsection{Assessment Instrument}
\label{ssec:asse-0458}
%-------------------------------------------------------------------------------

Assess through multi-rule simplification problems requiring students to apply product,
quotient, power, negative, and zero exponent properties in combination on both numeric and
variable bases.

\textbf{Example problems.}

\begin{enumerate}
    \item Simplify $\dfrac{(2x^3 y^{-2})^3}{4x^{-1} y^4}$, leaving no negative exponents.
    \item Evaluate $\left(\dfrac{3^{-2} \cdot 3^5}{3^0}\right)^2$ without a calculator.
    \item Show algebraically that $\left(\dfrac{a}{b}\right)^{-n} = \left(\dfrac{b}{a}\right)^n$ for $a, b \neq 0$.
\end{enumerate}

%===============================================================================
\section{Exponential Functions}
\label{sec:expo-6c26}
%===============================================================================

%-------------------------------------------------------------------------------
\subsection{Introduction}
\label{sss:intr-8800}
%-------------------------------------------------------------------------------

In mathematics, you can find the inverse of an exponential function by switching $x$ and $y$ around: 
\[
y = b^x \text{ becomes } x = b^y.
\]
The problem arises in finding the value of $y$. The logarithmic function solves this problem. All conversions of a logarithmic function into an exponential function follow the same pattern:
\[
x = b^y \text{ becomes } y = \log_b x.
\]
If a logarithm is given without a specified base, assume $b = 10$. Also, for logarithmic functions, $b > 0$ and $b \neq 1$. There are 2 cases where the log equals $x$:
\[
\log_b b^X = X \text{ and } b^{\log_b X} = X.
\]
To recap, a logarithm is the inverse function of an exponent. For example, the inverse of
the function
\[
f(x) = 3^x
\]
is
\[
f^{-1}(x) = \log_3 x.
\]
In general,
\[
y = b^x \iff x = \log_b y
\]
given that $b > 0$.

%-------------------------------------------------------------------------------
\subsection{Laws of Logarithmic Functions}
\label{sss:laws-71be}
%-------------------------------------------------------------------------------

When $X$ and $Y$ are positive:

1. \textbf{Product Rule}:

   \[
\log_b (XY) = \log_b X + \log_b Y
\]

2. \textbf{Quotient Rule}:

   \[
\log_b \left(\frac{X}{Y}\right) = \log_b X - \log_b Y
\]

3. \textbf{Power Rule}:

   \[
\log_b (X^k) = k \log_b X
\]

%-------------------------------------------------------------------------------
\subsection{Change of Base}
\label{sss:chan-d752}
%-------------------------------------------------------------------------------

When $x$ and $b$ are positive real numbers and not equal to 1, you can write 
\[
\log_a x
\]
as
\[
\frac{\log_b x}{\log_b a}.
\]
This works for natural logs as well.

\textbf{Example}:

Solving a Logarithmic Equation:

A logarithmic equation is an equation wherein one or more of the terms is a logarithm.

\textbf{Example Problem}:

Solve
\[
\log x + \log(x+2) = 2
\]
\textbf{Solution}:
\[
\log x + \log(x+2) = 2
\]

\[
\log (x(x+2)) = 2
\]

\[
x(x+2) = 100
\]

\[
x^2 + 2x = 100
\]

\[
(x+1)^2 = 101
\]

\[
x+1 = \pm \sqrt{101}
\]

\[
x = -1 \pm \sqrt{101}
\]
\textbf{Another Example}:

Calculate
\[
\log_2 8 = \frac{\log 8}{\log 2} = \frac{0.9}{0.3} = 3.
\]
Verify:
\[
2^3 = 8.
\]

%-------------------------------------------------------------------------------
\subsection{Resources}
\label{sss:reso-55b5}
%-------------------------------------------------------------------------------

\begin{itemize}
    \item \href{https://mathresearch.utsa.edu/wikiFiles/MAT1053/Logarithmic_Functions/MAT1053_M5.2Logarithmic_Functions.pdf}{Logarithmic Functions, Book Chapter}
    \item \href{https://mathresearch.utsa.edu/wikiFiles/MAT1053/Logarithmic_Functions/MAT1053_M5.2Logarithmic_FunctionsGN.pdf}{Guided Notes}
\end{itemize}

%-------------------------------------------------------------------------------
\subsection{Licensing}
\label{sss:lice-7352}
%-------------------------------------------------------------------------------

Content obtained and/or adapted from:

\begin{itemize}
    \item \href{https://en.wikibooks.org/wiki/A-level_Mathematics/OCR/C2/Logarithms_and_Exponentials}{Logarithms and Exponentials, Wikibooks: A-level Mathematics/OCR/C2 under a CC BY-SA license}
\end{itemize}

%-------------------------------------------------------------------------------
\subsection{Learning Objectives}
\label{ssec:lear-6683}
%-------------------------------------------------------------------------------

\begin{enumerate}
    \item 1.\textbf{Understand the Inverse Relationship}: Given an exponential function $y = b^x$, convert it into its logarithmic form $x = \log_b y$ and identify the base $b$ and the argument $y$. 2.\textbf{Apply Logarithmic Laws}: Use the product rule, quotient rule, and power rule for logarithms to simplify expressions involving logarithms, such as applying $ \log_b (XY) = \log_b X + \log_b Y, $ $ \log_b \left(\frac{X}{Y}\right) = \log_b X - \log_b Y, $ and $ \log_b (X^k) = k \log_b X. $ 3.\textbf{Solve Logarithmic Equations}: Solve equations involving logarithms by applying logarithmic properties, such as $ \log_b x + \log_b (x + 2) = 2, $ to find the value of $x$ through conversion and algebraic manipulation.
\end{enumerate}

%-------------------------------------------------------------------------------
\subsection{Assessment Instrument}
\label{ssec:asse-1f50}
%-------------------------------------------------------------------------------

Assess through problems requiring students to identify key features of exponential growth
and decay functions, evaluate and graph them, and solve basic exponential equations by
rewriting both sides with a common base. Include at least one applied problem involving
compound interest or population growth.

\textbf{Example problems.}

\begin{enumerate}
    \item Sketch $f(x) = 2^x$ and $g(x) = \left(\tfrac{1}{2}\right)^x$ on the same axes and state the domain, range, and horizontal asymptote of each.
    \item Solve for $x$: $4^{x-1} = 8^{x+2}$.
    \item A bacterial population doubles every 3 hours from an initial count of 500. Write an exponential model and find the population after 12 hours.
\end{enumerate}

%===============================================================================
\section{Logarithmic Properties}
\label{sec:loga-ebc4}
%===============================================================================
%-------------------------------------------------------------------------------
\subsection{Equality Rule}
\label{sss:equa-b299}
%-------------------------------------------------------------------------------

If
\[
\log_a(x) = \log_a(y)
\]
then
\[
x = y
\]

%-------------------------------------------------------------------------------
\subsection{Product Rule}
\label{sss:prod-1f6a}
%-------------------------------------------------------------------------------

\[
\log_a(xy) = \log_a(x) + \log_a(y)
\]

%-------------------------------------------------------------------------------
\subsection{Quotient Rule}
\label{sss:quot-e55e}
%-------------------------------------------------------------------------------

\[
\log_a \left(\frac{x}{y}\right) = \log_a(x) - \log_a(y)
\]

%-------------------------------------------------------------------------------
\subsection{Power Rule}
\label{sss:powe-f8c4}
%-------------------------------------------------------------------------------

\[
\log_a(x^n) = n \log_a(x)
\]

%-------------------------------------------------------------------------------
\subsection{Change-of-Base Rule}
\label{sss:chan-eefd}
%-------------------------------------------------------------------------------

\[
\log_a(x) = \frac{\log_b(x)}{\log_b(a)}
\]
where $b$ is any positive number not equal to 1.

%-------------------------------------------------------------------------------
\subsection{Resources}
\label{sss:reso-55b5}
%-------------------------------------------------------------------------------

\begin{itemize}
    \item \href{https://mathresearch.utsa.edu/wikiFiles/MAT1053/Logarithmic\%20Properties/MAT1053_M6.1Logarithmic_Properties.pdf}{Logarithmic Properties, Book Chapter}
    \item \href{https://mathresearch.utsa.edu/wikiFiles/MAT1053/Logarithmic\%20Properties/MAT1053_M6.1Logarithmic_PropertiesGN.pdf}{Guided Notes}
\end{itemize}

%-------------------------------------------------------------------------------
\subsection{Learning Objectives}
\label{ssec:lear-f632}
%-------------------------------------------------------------------------------

\begin{enumerate}
    \item Student Learning Objectives (SLOs)
    \item \textbf{Equality Rule}: Given the equation $ \log_a(x) = \log_a(y) $ students will be able to solve for $x$ and $y$ by setting $ x = y. $
    \item \textbf{Product Rule}: Students will apply the product rule to simplify expressions involving logarithms by using $ \log_a(xy) = \log_a(x) + \log_a(y). $
    \item \textbf{Quotient and Power Rules}: Students will demonstrate the ability to use the quotient rule $ \log_a \left(\frac{x}{y}\right) = \log_a(x) - \log_a(y) $ and the power rule $ \log_a(x^n) = n \log_a(x) $ to solve and simplify logarithmic expressions.
\end{enumerate}

%-------------------------------------------------------------------------------
\subsection{Assessment Instrument}
\label{ssec:asse-812c}
%-------------------------------------------------------------------------------

Assess through problems that require students to expand and condense logarithmic
expressions using the product, quotient, and power rules, and to apply the change-of-base
formula for numerical evaluation.

\textbf{Example problems.}

\begin{enumerate}
    \item Expand $\log_3\!\left(\dfrac{x^2 \sqrt{y}}{z^3}\right)$ completely using logarithm properties.
    \item Condense $3\ln x - \tfrac{1}{2}\ln y + 2\ln z$ into a single logarithm.
    \item Use the change-of-base formula to evaluate $\log_7 50$ to four decimal places.
\end{enumerate}

%===============================================================================
\section{Logarithmic Functions}
\label{sec:loga-7d37}
%===============================================================================
%-------------------------------------------------------------------------------
\subsection{Introduction}
\label{sss:intr-8800}
%-------------------------------------------------------------------------------

In mathematics, you can find the inverse of an exponential function by switching $x$ and $y$ around: $y = b^x \text{ becomes } x = b^y$.  The problem arises in finding the value of $y$. The logarithmic function solves this problem. All conversions of a logarithmic function into an exponential function follow the same pattern:  $x = b^y \text{ becomes } y = \log_b x$.  If a logarithm is given without a specified base, assume $b = 10$. Also, for logarithmic functions, $b > 0$ and $b \neq 1$. There are 2 cases where the log equals $x$:  $\log_b b^X = X \text{ and } b^{\log_b X} = X$.

To recap, a logarithm is the inverse function of an exponent. For example, the inverse of
the function
\[
f(x) = 3^x
\]
is
\[
f^{-1}(x) = \log_3 x.
\]
In general,
\[
y = b^x \iff x = \log_b y
\]
given that $b > 0$.

%-------------------------------------------------------------------------------
\subsection{Laws of Logarithmic Functions}
\label{sss:laws-71be}
%-------------------------------------------------------------------------------

When $X$ and $Y$ are positive:

1. \textbf{Product Rule}:

   \[
\log_b (XY) = \log_b X + \log_b Y
\]

2. \textbf{Quotient Rule}:

   \[
\log_b \left(\frac{X}{Y}\right) = \log_b X - \log_b Y
\]

3. \textbf{Power Rule}:

   \[
\log_b (X^k) = k \log_b X
\]

%-------------------------------------------------------------------------------
\subsection{Change of Base}
\label{sss:chan-d752}
%-------------------------------------------------------------------------------

When $x$ and $b$ are positive real numbers and not equal to 1, you can write 
\[
\log_a x
\]
as
\[
\frac{\log_b x}{\log_b a}.
\]
This works for natural logs as well.

\textbf{Example}:

Solving a Logarithmic Equation:

A logarithmic equation is an equation wherein one or more of the terms is a logarithm.

\textbf{Example Problem}:

Solve
\[
\log x + \log(x+2) = 2
\]
\textbf{Solution}:
\[
\log x + \log(x+2) = 2
\]

\[
\log (x(x+2)) = 2
\]

\[
x(x+2) = 100
\]

\[
x^2 + 2x = 100
\]

\[
(x+1)^2 = 101
\]

\[
x+1 = \pm \sqrt{101}
\]

\[
x = -1 \pm \sqrt{101}
\]
\textbf{Another Example}:

Calculate
\[
\log_2 8 = \frac{\log 8}{\log 2} = \frac{0.9}{0.3} = 3.
\]
Verify:
\[
2^3 = 8.
\]

%-------------------------------------------------------------------------------
\subsection{Resources}
\label{sss:reso-55b5}
%-------------------------------------------------------------------------------

\begin{itemize}
    \item \href{https://mathresearch.utsa.edu/wikiFiles/MAT1053/Logarithmic_Functions/MAT1053_M5.2Logarithmic_Functions.pdf}{Logarithmic Functions, Book Chapter}
    \item \href{https://mathresearch.utsa.edu/wikiFiles/MAT1053/Logarithmic_Functions/MAT1053_M5.2Logarithmic_FunctionsGN.pdf}{Guided Notes}
\end{itemize}

%-------------------------------------------------------------------------------
\subsection{Licensing}
\label{sss:lice-7352}
%-------------------------------------------------------------------------------

Content obtained and/or adapted from:

\begin{itemize}
    \item \href{https://en.wikibooks.org/wiki/A-level_Mathematics/OCR/C2/Logarithms_and_Exponentials}{Logarithms and Exponentials, Wikibooks: A-level Mathematics/OCR/C2 under a CC BY-SA license}
\end{itemize}

%-------------------------------------------------------------------------------
\subsection{Learning Objectives}
\label{ssec:lear-1bc3}
%-------------------------------------------------------------------------------

\begin{enumerate}
    \item 1.\textbf{Understand the Inverse Relationship}: Given an exponential function $y = b^x$, convert it into its logarithmic form $x = \log_b y$ and identify the base $b$ and the argument $y$. 2.\textbf{Apply Logarithmic Laws}: Use the product rule, quotient rule, and power rule for logarithms to simplify expressions involving logarithms, such as applying $ \log_b (XY) = \log_b X + \log_b Y, $ $ \log_b \left(\frac{X}{Y}\right) = \log_b X - \log_b Y, $ and $ \log_b (X^k) = k \log_b X. $ 3.\textbf{Solve Logarithmic Equations}: Solve equations involving logarithms by applying logarithmic properties, such as $ \log_b x + \log_b (x + 2) = 2, $ to find the value of $x$ through conversion and algebraic manipulation.
\end{enumerate}

%-------------------------------------------------------------------------------
\subsection{Assessment Instrument}
\label{ssec:asse-ebc6}
%-------------------------------------------------------------------------------

Assess through problems requiring students to graph logarithmic functions, identify domain
and vertical asymptote, convert between exponential and logarithmic form, and solve basic
logarithmic equations. At least one problem should involve the natural logarithm.

\textbf{Example problems.}

\begin{enumerate}
    \item Sketch $f(x) = \log_2(x - 1) + 3$ and state its domain, range, and vertical asymptote.
    \item Solve for $x$: $\log_4(x + 3) + \log_4(x - 1) = 2$.
    \item Rewrite $e^{3x} = 7$ in logarithmic form and solve for $x$ in exact form.
\end{enumerate}

%===============================================================================
\section{Quadratic Equations}
\label{sec:quad-7bb4}
%===============================================================================
%-------------------------------------------------------------------------------
\subsection{Quadratic Formula}
\label{sss:quad-20dd}
%-------------------------------------------------------------------------------

The quadratic formula for the roots of the general quadratic equation is:
\[
x = \frac{-b \pm \sqrt{b^2 - 4ac}}{2a}
\]

%-------------------------------------------------------------------------------
\subsection{Definition}
\label{sss:defi-3061}
%-------------------------------------------------------------------------------

In algebra, a quadratic equation (from the Latin quadratus for ``square'') is any equation
that can be rearranged in standard form as
\[
ax^2 + bx + c = 0
\]
where $x$ represents an unknown, and $a$, $b$, and $c$ represent known numbers, where $a \neq 0$. If $a = 0$, then the equation is linear, not quadratic, as there is no $ax^2$ term. The numbers $a$, $b$, and $c$ are the coefficients of the equation and may be distinguished by calling them, respectively, the quadratic coefficient, the linear coefficient and the constant or free term.

The values of $x$ that satisfy the equation are called solutions of the equation, and roots or zeros of the expression on its left-hand side. A quadratic equation has at most two solutions. If there is only one solution, one says that it is a double root. If all the coefficients are real numbers, there are either two real solutions, or a single real double root, or two complex solutions. A quadratic equation always has two roots, if complex roots are included and a double root is counted for two.

A quadratic equation can be factored into an equivalent equation
\[
ax^2 + bx + c = a(x - r)(x - s) = 0
\]
where $r$ and $s$ are the solutions for $x$. Completing the square on a quadratic equation in standard form results in the quadratic formula, which expresses the solutions in terms of $a$, $b$, and $c$. Solutions to problems that can be expressed in terms of quadratic equations were known as early as 2000 BC.

Because the quadratic equation involves only one unknown, it is called ``univariate''. The quadratic equation contains only powers of $x$ that are non-negative integers, and therefore it is a polynomial equation. In particular, it is a second-degree polynomial equation, since the greatest power is two.

%-------------------------------------------------------------------------------
\subsection{Solving the Quadratic Equation}
\label{sss:solv-8ae7}
%-------------------------------------------------------------------------------

A quadratic equation with real or complex coefficients has two solutions, called roots.
These two solutions may or may not be distinct, and they may or may not be real.

%-------------------------------------------------------------------------------
\subsection{Factoring by Inspection}
\label{sss:fact-1b61}
%-------------------------------------------------------------------------------

It may be possible to express a quadratic equation $ax^2 + bx + c = 0$ as a product $(px + q)(rx + s) = 0$. In some cases, it is possible, by simple inspection, to determine values of $p$, $q$, $r$, and $s$ that make the two forms equivalent to one another. If the quadratic equation is written in the second form, then the ``Zero Factor Property'' states that the quadratic equation is satisfied if $px + q = 0$ or $rx + s = 0$. Solving these two linear equations provides the roots of the quadratic.

For most students, factoring by inspection is the first method of solving quadratic equations to which they are exposed. If one is given a quadratic equation in the form $x^2 + bx + c = 0$, the sought factorization has the form $(x + q)(x + s)$, and one has to find two numbers $q$ and $s$ that add up to $b$ and whose product is $c$ (this is sometimes called ``Vieta''''s rule'' and is related to Vieta''''s formulas). As an example, $x^2 + 5x + 6$ factors as $(x + 3)(x + 2)$. The more general case where $a$ does not equal 1 can require a considerable effort in trial and error guess-and-check, assuming that it can be factored at all by inspection.

Except for special cases such as where $b = 0$ or $c = 0$, factoring by inspection only works for quadratic equations that have rational roots. This means that the great majority of quadratic equations that arise in practical applications cannot be solved by factoring by inspection.

%-------------------------------------------------------------------------------
\subsection{Completing the Square}
\label{sss:comp-a359}
%-------------------------------------------------------------------------------

The process of completing the square makes use of the algebraic identity
\[
x^2 + 2hx + h^2 = (x + h)^2,
\]
which represents a well-defined algorithm that can be used to solve any quadratic equation. Starting with a quadratic equation in standard form, $ax^2 + bx + c = 0$:

1. Divide each side by $a$, the coefficient of the squared term.

2. Subtract the constant term $c/a$ from both sides.

3. Add the square of one-half of $b/a$, the coefficient of $x$, to both sides. This ``completes the square'', converting the left side into a perfect square.

4. Write the left side as a square and simplify the right side if necessary.

5. Produce two linear equations by equating the square root of the left side with the
positive and negative square roots of the right side.

6. Solve each of the two linear equations.

We illustrate use of this algorithm by solving $2x^2 + 4x - 4 = 0$:

1. $x^2 + 2x - 2 = 0$

2. $x^2 + 2x = 2$

3. $x^2 + 2x + 1 = 2 + 1$

4. $(x + 1)^2 = 3$

5. $x + 1 = \pm \sqrt{3}$

6. $x = -1 \pm \sqrt{3}$

The plus–minus symbol ``$\pm$'' indicates that both $x = -1 + \sqrt{3}$ and $x = -1 - \sqrt{3}$ are solutions of the quadratic equation.

%-------------------------------------------------------------------------------
\subsection{Quadratic Formula and Its Derivation}
\label{sss:quad-d91c}
%-------------------------------------------------------------------------------

Completing the square can be used to derive a general formula for solving quadratic
equations, called the quadratic formula. The mathematical proof will now be briefly
summarized. It can easily be seen, by polynomial expansion, that the following equation is
equivalent to the quadratic equation:
\[
\left(x + \frac{b}{2a}\right)^2 = \frac{b^2 - 4ac}{4a^2}.
\]
Taking the square root of both sides, and isolating $x$, gives:
\[
x = \frac{-b \pm \sqrt{b^2 - 4ac}}{2a}.
\]
Some sources, particularly older ones, use alternative parameterizations of the quadratic equation such as $ax^2 + 2bx + c = 0$ or $ax^2 - 2bx + c = 0$, where $b$ has a magnitude one half of the more common one, possibly with opposite sign. These result in slightly different forms for the solution, but are otherwise equivalent.

A number of alternative derivations can be found in the literature. These proofs are
simpler than the standard completing the square method, represent interesting applications
of other frequently used techniques in algebra, or offer insight into other areas of
mathematics.

A lesser known quadratic formula, as used in Muller''''s method, provides the same roots
via the equation:
\[
x = \frac{2c}{-b \pm \sqrt{b^2 - 4ac}}.
\]
This can be deduced from the standard quadratic formula by Vieta''''s formulas, which assert that the product of the roots is $c/a$.

One property of this form is that it yields one valid root when $a = 0$, while the other root contains division by zero, because when $a = 0$, the quadratic equation becomes a linear equation, which has one root. By contrast, in this case, the more common formula has a division by zero for one root and an indeterminate form $0/0$ for the other root. On the other hand, when $c = 0$, the more common formula yields two correct roots whereas this form yields the zero root and an indeterminate form $0/0$.

%-------------------------------------------------------------------------------
\subsection{Reduced Quadratic Equation}
\label{sss:redu-a641}
%-------------------------------------------------------------------------------

It is sometimes convenient to reduce a quadratic equation so that its leading coefficient is one. This is done by dividing both sides by $a$, which is always possible since $a$ is non-zero. This produces the reduced quadratic equation:
\[
x^2 + px + q = 0,
\]
where $p = b/a$ and $q = c/a$. This monic polynomial equation has the same solutions as the original.

The quadratic formula for the solutions of the reduced quadratic equation, written in
terms of its coefficients, is:
\[
x = \frac{1}{2} \left(-p \pm \sqrt{p^2 - 4q}\right),
\]
or equivalently:
\[
x = -\frac{p}{2} \pm \sqrt{\left(\frac{p}{2}\right)^2 - q}.
\]

%-------------------------------------------------------------------------------
\subsection{Discriminant}
\label{sss:disc-db3e}
%-------------------------------------------------------------------------------

This figure plots three quadratic functions on a single Cartesian plane graph to
illustrate the effects of discriminant values. When the discriminant, delta, is positive,
the parabola intersects the x-axis at two points. When delta is zero, the vertex of the
parabola touches the x-axis at a single point. When delta is negative, the parabola does
not intersect the x-axis at all.

% [Image: graph]

In the quadratic formula, the expression underneath the square root sign is called the
discriminant of the quadratic equation, and is often represented using an upper case D or
an upper case Greek delta:

\[
\Delta = b^2 - 4ac.
\]

A quadratic equation with real coefficients can have either one or two distinct real
roots, or two distinct complex roots. In this case the discriminant determines the number
and nature of the roots. There are three cases:

1. If the discriminant is positive, then there are two distinct roots

   \[
\frac{-b + \sqrt{\Delta}}{2a} \text{ and } \frac{-b - \sqrt{\Delta}}{2a},
\]

both of which are real numbers. For quadratic equations with rational coefficients, if the
discriminant is a square number, then the roots are rational: in other cases they may be
quadratic irrationals.

2. If the discriminant is zero, then there is exactly one real root

   \[
-\frac{b}{2a},
\]

sometimes called a repeated or double root.

3. If the discriminant is negative, then there are no real roots. Rather, there are two
distinct (non-real) complex roots

   \[
-\frac{b}{2a} + i \frac{\sqrt{-\Delta}}{2a} \text{ and } -\frac{b}{2a} - i \frac{\sqrt{-\Delta}}{2a},
\]

   which are complex conjugates of each other. In these expressions, $i$ is the imaginary unit.

Thus the roots are distinct if and only if the discriminant is non-zero, and the roots are
real if and only if the discriminant is non-negative.

%-------------------------------------------------------------------------------
\subsection{Geometric Interpretation}
\label{sss:geom-084c}
%-------------------------------------------------------------------------------

Graph of

\[
y = ax^2 + bx + c,
\] 

where $a$ and the discriminant 

\[
b^2 - 4ac
\] 

are positive, with:

\begin{itemize}
    \item Roots and y-intercept in red
    \item Vertex and axis of symmetry in blue
    \item Focus and directrix in pink
\end{itemize}

% [Image: graph]

The function $f(x) = ax^2 + bx + c$ is a quadratic function. The graph of any quadratic function has the same general shape, which is called a parabola. The location and size of the parabola, and how it opens, depend on the values of $a$, $b$, and $c$. As shown in Figure 1, if $a > 0$, the parabola has a minimum point and opens upward. If $a < 0$, the parabola has a maximum point and opens downward. The extreme point of the parabola, whether minimum or maximum, corresponds to its vertex. The x-coordinate of the vertex will be located at 

\[
x = -\frac{b}{2a},
\] 

and the y-coordinate of the vertex may be found by substituting this x-value into the function. The y-intercept is located at the point $(0, c)$.

The solutions of the quadratic equation $ax^2 + bx + c = 0$ correspond to the roots of the function $f(x) = ax^2 + bx + c$, since they are the values of $x$ for which $f(x) = 0$. As shown in Figure 2, if $a$, $b$, and $c$ are real numbers and the domain of $f$ is the set of real numbers, then the roots of $f$ are exactly the x-coordinates of the points where the graph touches the x-axis. As shown in Figure 3, if the discriminant is positive, the graph touches the x-axis at two points; if zero, the graph touches at one point; and if negative, the graph does not touch the x-axis.

%-------------------------------------------------------------------------------
\subsection{Quadratic Factorization}
\label{sss:quad-1f80}
%-------------------------------------------------------------------------------

The term

\[
x - r
\] 

is a factor of the polynomial

\[
ax^2 + bx + c
\] 

if and only if $r$ is a root of the quadratic equation 

\[
ax^2 + bx + c = 0.
\] 

It follows from the quadratic formula that

\[
ax^2 + bx + c = a \left(x - \frac{-b + \sqrt{b^2 - 4ac}}{2a}\right) \left(x - \frac{-b - \sqrt{b^2 - 4ac}}{2a}\right).
\]

In the special case $b^2 = 4ac$ where the quadratic has only one distinct root (i.e., the discriminant is zero), the quadratic polynomial can be factored as 

\[
ax^2 + bx + c = a \left(x + \frac{b}{2a}\right)^2.
\]

%-------------------------------------------------------------------------------
\subsection{Graphical Solution}
\label{sss:grap-2f4d}
%-------------------------------------------------------------------------------

A quadratic function without real root: \[
y = (x - 5)^2 + 9.
\] The ``3'' is the imaginary part of the x-intercept. The real part is the x-coordinate of the vertex. Thus the roots are \[
5 \pm 3i.
\]

% [Image: graph]

The solutions of the quadratic equation

\[
ax^2 + bx + c = 0
\]

may be deduced from the graph of the quadratic function

\[
y = ax^2 + bx + c,
\]

which is a parabola.

\begin{itemize}
    \item If the parabola intersects the x-axis in two points, there are two real roots, which are the x-coordinates of these two points (also called x-intercept).
    \item If the parabola is tangent to the x-axis, there is a double root, which is the x-coordinate of the contact point between the graph and parabola.
    \item If the parabola does not intersect the x-axis, there are two complex conjugate roots. Although these roots cannot be visualized on the graph, their real and imaginary parts can be.
\end{itemize}

Let $h$ and $k$ be respectively the x-coordinate and the y-coordinate of the vertex of the parabola (that is the point with maximal or minimal y-coordinate). The quadratic function may be rewritten

\[
y = a(x - h)^2 + k.
\]

Let $d$ be the distance between the point of y-coordinate $2k$ on the axis of the parabola, and a point on the parabola with the same y-coordinate (see the figure; there are two such points, which give the same distance, because of the symmetry of the parabola). Then the real part of the roots is $h$, and their imaginary parts are $\pm d$. That is, the roots are

\[
h + id \quad \text{and} \quad h - id,
\]

or in the case of the example of the figure

\[
5 + 3i \quad \text{and} \quad 5 - 3i.
\]

%-------------------------------------------------------------------------------
\subsection{Avoiding Loss of Significance}
\label{sss:avoi-6288}
%-------------------------------------------------------------------------------

Although the quadratic formula provides an exact solution, the result is not exact if real
numbers are approximated during the computation, as usual in numerical analysis, where
real numbers are approximated by floating point numbers (called ``reals'' in many
programming languages). In this context, the quadratic formula is not completely stable.

This occurs when the roots have different orders of magnitude, or, equivalently, when $b^2$ and $b^2 - 4ac$ are close in magnitude. In this case, the subtraction of two nearly equal numbers will cause loss of significance or catastrophic cancellation in the smaller root. To avoid this, the root that is smaller in magnitude, $r$, can be computed as 

\[
\frac{c / a}{R}
\] 

where $R$ is the root that is bigger in magnitude.

A second form of cancellation can occur between the terms $b^2$ and $4ac$ of the discriminant, that is when the two roots are very close. This can lead to loss of up to half of correct significant figures in the roots.

%-------------------------------------------------------------------------------
\subsection{Examples and Applications}
\label{sss:exam-0907}
%-------------------------------------------------------------------------------

% [Image: figure]

The trajectory of the cliff jumper is parabolic because horizontal displacement is a
linear function of time

\[
x = v_x t,
\]

while vertical displacement is a quadratic function of time

\[
y = \frac{1}{2} a t^2 + v_y t + h.
\] 

As a result, the path follows the quadratic equation

\[
y = \frac{a}{2v_x^2} x^2 + \frac{v_y}{v_x} x + h,
\]

where $v_x$ and $v_y$ are the horizontal and vertical components of the original velocity, $a$ is gravitational acceleration, and $h$ is the original height. The $a$ value should be considered negative here, as its direction (downwards) is opposite to the height measurement (upwards).

The golden ratio is found as the positive solution of the quadratic equation

\[
x^2 - x - 1 = 0.
\]

The equations of the circle and the other conic sections: ellipses, parabolas, and
hyperbolas: are quadratic equations in two variables.

Given the cosine or sine of an angle, finding the cosine or sine of the angle that is half
as large involves solving a quadratic equation.

The process of simplifying expressions involving the square root of an expression
involving the square root of another expression involves finding the two solutions of a
quadratic equation.

Descartes'''' theorem states that for every four kissing (mutually tangent) circles, their
radii satisfy a particular quadratic equation.

The equation given by Fuss'''' theorem, giving the relation among the radius of a
bicentric quadrilateral''''s inscribed circle, the radius of its circumscribed circle, and
the distance between the centers of those circles, can be expressed as a quadratic
equation for which the distance between the two circles'''' centers in terms of their
radii is one of the solutions. The other solution of the same equation in terms of the
relevant radii gives the distance between the circumscribed circle''''s center and the
center of the excircle of an ex-tangential quadrilateral.

Critical points of a cubic function and inflection points of a quartic function are found
by solving a quadratic equation.

%-------------------------------------------------------------------------------
\subsection{Resources}
\label{sss:reso-55b5}
%-------------------------------------------------------------------------------

\begin{itemize}
    \item \href{https://mathresearch.utsa.edu/wikiFiles/MAT1053/Quadratic_Functions/MAT1053_M2.2Quadratic_Functions.pdf}{Quadratic Functions} Book Chapters
\end{itemize}

\begin{itemize}
    \item \href{https://mathresearch.utsa.edu/wikiFiles/MAT1053/Quadratic_Functions/MAT1053_M2.2Quadratic_FunctionsGN.pdf}{Guided Notes}
\end{itemize}

%-------------------------------------------------------------------------------
\subsection{Licensing}
\label{sss:lice-7352}
%-------------------------------------------------------------------------------

Content obtained and/or adapted from:

\begin{itemize}
    \item \href{https://en.wikipedia.org/wiki/Quadratic_equation}{Quadratic equation, Wikipedia} under a CC BY-SA license
\end{itemize}

%-------------------------------------------------------------------------------
\subsection{Learning Objectives}
\label{ssec:lear-46cc}
%-------------------------------------------------------------------------------

\begin{enumerate}
    \item 1.\textbf{Understanding and Applying the Quadratic Formula:}
    \item Students will be able to derive and apply the quadratic formula $ x = \frac{-b \pm \sqrt{b^2 - 4ac}}{2a} $ to solve any quadratic equation of the form $ax^2 + bx + c = 0$, identifying the nature of the roots based on the discriminant ($\Delta = b^2 - 4ac$). 2.\textbf{Factoring Quadratic Equations:}
    \item Students will be able to factor quadratic equations of the form $ax^2 + bx + c = 0$ into the product of two binomials $(px + q)(rx + s) = 0$, using methods such as factoring by inspection and Vieta's formulas when applicable. 3.\textbf{Solving Quadratic Equations by Completing the Square:}
    \item Students will be able to solve quadratic equations by completing the square, transforming equations into the form $ (x + h)^2 = k $ and subsequently solving for $x$ to find the roots of the equation.
\end{enumerate}

%-------------------------------------------------------------------------------
\subsection{Assessment Instrument}
\label{ssec:asse-de40}
%-------------------------------------------------------------------------------

Assess through problems requiring students to solve quadratic equations by at least three
methods: factoring, completing the square, and the quadratic formula. Problems should also
require students to use the discriminant to predict the number and type of solutions
before solving.

\textbf{Example problems.}

\begin{enumerate}
    \item Solve $x^2 - 5x + 6 = 0$ by factoring, by completing the square, and by the quadratic formula, and verify that all three methods agree.
    \item Use the discriminant to classify the solutions of $3x^2 + x + 2 = 0$, then find all solutions.
    \item A rectangular garden has area $40\,\mathrm{m}^2$ and its length is 3~m more than its width. Write and solve a quadratic equation to find the dimensions.
\end{enumerate}

%===============================================================================
\section{Quadratic Functions}
\label{sec:quad-e70d}
%===============================================================================
In algebra, a quadratic function, a quadratic polynomial, or a polynomial of degree 2, is
a polynomial function where the highest-degree term is of the second degree.

A quadratic polynomial with two real roots (crossings of the x-axis) and no complex roots
can be written as:
\[
f(x) = ax^2 + bx + c, \quad a \neq 0
\]
% [Image: graph]

In this case, the graph of the function is a parabola with an axis of symmetry parallel to
the y-axis. If the quadratic function is set to zero, it forms a quadratic equation, and
the solutions are called the roots of the function.

For a bivariate case involving variables $x$ and $y$, the quadratic function has the form:
\[
f(x, y) = ax^2 + by^2 + cxy + dx + ey + f
\]
where at least one of $a$, $b$, or $c$ is non-zero. Setting this function equal to zero results in a conic section, which could be a circle, ellipse, parabola, or hyperbola.

In three variables $x$, $y$, and $z$, a quadratic function is expressed as:
\[
f(x, y, z) = ax^2 + by^2 + cz^2 + dxy + exz + fyz + gx + hy + iz + j
\]
where at least one of the coefficients $a$, $b$, $c$, $d$, $e$, or $f$ is non-zero. For more than three variables, the resulting surface is called a quadric, with the highest degree term being of degree 2.

%-------------------------------------------------------------------------------
\subsection{Etymology}
\label{sss:etym-a4f6}
%-------------------------------------------------------------------------------

The adjective \textit{quadratic} comes from the Latin word \textit{quadratum} (``square''). In algebra, a term like $x^2$ is called a square because it represents the area of a square with side $x$.

%-------------------------------------------------------------------------------
\subsection{Terminology}
\label{sss:term-40dc}
%-------------------------------------------------------------------------------

%-------------------------------------------------------------------------------
\subsection{Coefficients}
\label{sss:coef-f6ff}
%-------------------------------------------------------------------------------

The coefficients of a polynomial are often real or complex numbers, but a polynomial may
be defined over any ring.

%-------------------------------------------------------------------------------
\subsection{Degree}
\label{sss:degr-fb81}
%-------------------------------------------------------------------------------

When referring to a ``quadratic polynomial,'' authors may mean either ``having degree
exactly 2'' or ``having degree at most 2.'' If the degree is less than 2, this is
sometimes called a ``degenerate case.'' The term ``order'' is also used to denote
``degree,'' e.g., a second-order polynomial.

%-------------------------------------------------------------------------------
\subsection{Variables}
\label{sss:vari-87cd}
%-------------------------------------------------------------------------------

A quadratic polynomial may involve a single variable $x$ (the univariate case) or multiple variables such as $x$, $y$, and $z$ (the multivariate case).

%-------------------------------------------------------------------------------
\subsection{The One-Variable Case}
\label{sss:theo-63b6}
%-------------------------------------------------------------------------------

Any single-variable quadratic polynomial can be written as:
\[
f(x) = ax^2 + bx + c
\]
where $x$ is the variable, and $a$, $b$, and $c$ are the coefficients. In elementary algebra, such polynomials often appear in the form of a quadratic equation:
\[
ax^2 + bx + c = 0
\]
The solutions to this equation are called the roots of the quadratic polynomial and can be
found through factorization, completing the square, graphing, Newton''s method, or using
the quadratic formula. Each quadratic polynomial has an associated quadratic function,
whose graph is a parabola.

%-------------------------------------------------------------------------------
\subsection{Bivariate Case}
\label{sss:biva-dd5b}
%-------------------------------------------------------------------------------

Any quadratic polynomial with two variables can be written as:
\[
f(x, y) = ax^2 + by^2 + cxy + dx + ey + f
\]
where $x$ and $y$ are the variables and $a$, $b$, $c$, $d$, $e$, and $f$ are the coefficients. Such polynomials are fundamental to the study of conic sections, characterized by setting $f(x, y)$ to zero. Quadratic polynomials with three or more variables correspond to quadric surfaces and hypersurfaces. In linear algebra, quadratic polynomials can be generalized to quadratic forms on vector spaces.

%-------------------------------------------------------------------------------
\subsection{Forms of a Univariate Quadratic Function}
\label{sss:form-3213}
%-------------------------------------------------------------------------------

A univariate quadratic function can be expressed in three formats:

1. \textbf{Standard Form}:

   \[
f(x) = ax^2 + bx + c
\]

2. \textbf{Factored Form}:

   \[
f(x) = a(x - r_1)(x - r_2)
\]

   where $r_1$ and $r_2$ are the roots of the quadratic function.

3. \textbf{Vertex Form}:

   \[
f(x) = a(x - h)^2 + k
\]

   where $h$ and $k$ are the $x$ and $y$ coordinates of the vertex, respectively.

The coefficient $a$ is the same in all three forms. To convert the standard form to factored form, one uses the quadratic formula to determine the roots $r_1$ and $r_2$. Converting to vertex form involves completing the square. To convert from factored form (or vertex form) to standard form, one multiplies, expands, and/or distributes the factors.

%-------------------------------------------------------------------------------
\subsection{Graph of the Univariate Function}
\label{sss:grap-3928}
%-------------------------------------------------------------------------------

1. For $a = \{0.1, 0.3, 1, 3\}$:

   \[
f(x) = ax^2
\]

2. For $b = \{1, 2, 3, 4\}$:

   \[
f(x) = x^2 + bx
\]

3. For $b = \{-1, -2, -3, -4\}$:

   \[
f(x) = x^2 + bx
\]

Regardless of the format, the graph of a univariate quadratic function:
\[
f(x) = ax^2 + bx + c
\]
is a parabola. Specifically:

\begin{itemize}
    \item If $a > 0$, the parabola opens upwards.
    \item If $a < 0$, the parabola opens downwards.
    \item The coefficient $a$ affects the curvature of the graph; a larger magnitude of $a$ results in a more sharply curved graph.
\end{itemize}

The coefficients $b$ and $a$ together control the axis of symmetry of the parabola (the $x$-coordinate of the vertex and the $h$ parameter in the vertex form), which is located at:
\[
x = -\frac{b}{2a}
\]
The coefficient $c$ controls the height of the parabola where it intercepts the $y$-axis.

% [Image: graph]

%-------------------------------------------------------------------------------
\subsection{Vertex}
\label{sss:vert-2b5b}
%-------------------------------------------------------------------------------

The vertex of a parabola is the point where it turns, also known as the turning point. If the quadratic function is in vertex form, the vertex is $(h, k)$. By completing the square, the standard form:
\[
f(x) = ax^2 + bx + c
\]
can be converted into:
\[
f(x) = a(x - h)^2 + k = a\left(x - \frac{-b}{2a}\right)^2 + \left(c - \frac{b^2}{4a}\right)
\]
So the vertex $(h, k)$ in standard form is:
\[
\left(-\frac{b}{2a}, c - \frac{b^2}{4a}\right)
\]
If the quadratic function is in factored form:
\[
f(x) = a(x - r_1)(x - r_2)
\]
the average of the two roots, $\frac{r_1 + r_2}{2}$, is the $x$-coordinate of the vertex, and thus the vertex $(h, k)$ is:
\[
\left(\frac{r_1 + r_2}{2}, f\left(\frac{r_1 + r_2}{2}\right)\right)
\]
The vertex is the maximum point if $a < 0$, or the minimum point if $a > 0$.

The vertical line:
\[
x = h = -\frac{b}{2a}
\]
that passes through the vertex is also the axis of symmetry of the parabola.

% [Image: graph]

% [Image: graph]

%-------------------------------------------------------------------------------
\subsection{Maximum and Minimum Points}
\label{sss:maxi-056e}
%-------------------------------------------------------------------------------

Using calculus, the vertex point, being a maximum or minimum of the function, can be found
by solving for the roots of the derivative:
\[
f(x) = ax^2 + bx + c \Rightarrow f''(x) = 2ax + b
\]
Setting $f''(x) = 0$ results in:
\[
x = -\frac{b}{2a}
\]
The corresponding function value is:
\[
f(x) = a\left(-\frac{b}{2a}\right)^2 + b\left(-\frac{b}{2a}\right) + c = c - \frac{b^2}{4a}
\]
Thus, the vertex point coordinates $(h, k)$ can be expressed as:
\[
\left(-\frac{b}{2a}, c - \frac{b^2}{4a}\right)
\]

%-------------------------------------------------------------------------------
\subsection{Roots of the Univariate Function}
\label{sss:root-c66c}
%-------------------------------------------------------------------------------

\textbf{Graph of $y = ax^2 + bx + c$}, where $a$ and the discriminant $b^2 - 4ac$ are positive, with:

\begin{itemize}
    \item Roots and y-intercept in red
    \item Vertex and axis of symmetry in blue
    \item Focus and directrix in pink
\end{itemize}

% [Image: graph]

\textbf{Visualisation of the complex roots of $y = ax^2 + bx + c$}: The parabola is rotated 180° about its vertex (orange). Its x-intercepts are rotated 90° around their midpoint, and the Cartesian plane is interpreted as the complex plane (green).

% [Image: graph]

%-------------------------------------------------------------------------------
\subsection{Exact Roots}
\label{sss:exac-bff6}
%-------------------------------------------------------------------------------

The roots (or zeros), $r_1$ and $r_2$, of the univariate quadratic function
\[
f(x) = ax^2 + bx + c = a(x - r\_1)(x - r\_2),
\]
are the values of $x$ for which $f(x) = 0$.

When the coefficients $a$, $b$, and $c$ are real or complex, the roots are:
\[
r\_1 = \frac{-b - \sqrt{b^2 - 4ac}}{2a},
\]

\[
r\_2 = \frac{-b + \sqrt{b^2 - 4ac}}{2a}.
\]

%-------------------------------------------------------------------------------
\subsection{Upper Bound on the Magnitude of the Roots}
\label{sss:uppe-d41d}
%-------------------------------------------------------------------------------

The modulus of the roots of a quadratic $ax^2 + bx + c$ can be no greater than
\[
\frac{\max(|a|, |b|, |c|)}{|a|} \times \phi,
\]
where $\phi$ is the golden ratio $\frac{1 + \sqrt{5}}{2}$.

%-------------------------------------------------------------------------------
\subsection{The Square Root of a Univariate Quadratic Function}
\label{sss:thes-bc34}
%-------------------------------------------------------------------------------

The square root of a univariate quadratic function gives rise to one of the four conic
sections, almost always either an ellipse or a hyperbola.

If $a > 0$, then the equation
\[
y = \pm \sqrt{ax^2 + bx + c}
\]
describes a hyperbola, as can be seen by squaring both sides. The directions of the axes of the hyperbola are determined by the ordinate of the minimum point of the corresponding parabola $y_p = ax^2 + bx + c$. If the ordinate is negative, then the hyperbola''s major axis (through its vertices) is horizontal, while if the ordinate is positive then the hyperbola''s major axis is vertical.

If $a < 0$, then the equation
\[
y = \pm \sqrt{ax^2 + bx + c}
\]
describes either a circle or another ellipse or nothing at all. If the ordinate of the maximum point of the corresponding parabola $y_p = ax^2 + bx + c$ is positive, then its square root describes an ellipse, but if the ordinate is negative then it describes an empty locus of points.

%-------------------------------------------------------------------------------
\subsection{Iteration}
\label{sss:iter-6f01}
%-------------------------------------------------------------------------------

To iterate a function $f(x) = ax^2 + bx + c$, one applies the function repeatedly, using the output from one iteration as the input to the next.

One cannot always deduce the analytic form of $f^{(n)}(x)$, which means the $n$th iteration of $f(x)$. (The superscript can be extended to negative numbers, referring to the iteration of the inverse of $f(x)$ if the inverse exists.) But there are some analytically tractable cases.

For example, for the iterative equation
\[
f(x) = a(x - c)^2 + c,
\]
one has
\[
f(x) = a(x - c)^2 + c = h^{-1}(g(h(x))),
\]
where
\[
g(x) = ax^2
\]
and
\[
h(x) = x - c.
\]
So by induction,
\[
f^{(n)}(x) = h^{-1}(g^{(n)}(h(x)))
\]
can be obtained, where $g^{(n)}(x)$ can be easily computed as
\[
g^{(n)}(x) = a^{2^n - 1} x^{2^n}.
\]
Finally, we have
\[
f^{(n)}(x) = a^{2^n - 1}(x - c)^{2^n} + c
\]
as the solution.

See Topological conjugacy for more detail about the relationship between $f$ and $g$. And see Complex quadratic polynomial for the chaotic behavior in the general iteration.

%-------------------------------------------------------------------------------
\subsection{The Logistic Map}
\label{sss:thel-fba9}
%-------------------------------------------------------------------------------

\[
x\_{n+1} = rx\_n (1 - x\_n), \quad 0 \leq x\_0 < 1
\]
with parameter $2 < r < 4$ can be solved in certain cases, one of which is chaotic and one of which is not. In the chaotic case $r = 4$, the solution is
\[
x\_n = \sin^2(2^n \theta \pi)
\]
where the initial condition parameter $\theta$ is given by
\[
\theta = \frac{1}{\pi} \sin^{-1}(x\_0^{1/2}).
\]
For rational $\theta$, after a finite number of iterations $x_n$ maps into a periodic sequence. But almost all $\theta$ are irrational, and for irrational $\theta$, $x_n$ never repeats itself – it is non-periodic and exhibits sensitive dependence on initial conditions, so it is said to be chaotic.

The solution of the logistic map when $r = 2$ is
\[
x\_n = \frac{1}{2} - \frac{1}{2} (1 - 2x\_0)^{2^n}
\]
for $x_0 \in [0, 1)$. Since $(1 - 2x_0) \in (-1, 1)$ for any value of $x_0$ other than the unstable fixed point 0, the term $(1 - 2x_0)^{2^n}$ goes to 0 as $n$ goes to infinity, so $x_n$ goes to the stable fixed point $\frac{1}{2}$.

%-------------------------------------------------------------------------------
\subsection{Bivariate (Two-Variable) Quadratic Function}
\label{sss:biva-483b}
%-------------------------------------------------------------------------------

A bivariate quadratic function is a second-degree polynomial of the form
\[
f(x, y) = Ax^2 + By^2 + Cx + Dy + Exy + F
\]
where $A$, $B$, $C$, $D$, and $E$ are fixed coefficients and $F$ is the constant term. Such a function describes a quadratic surface. Setting $f(x, y)$ equal to zero describes the intersection of the surface with the plane $z = 0$, which is a locus of points equivalent to a conic section.

%-------------------------------------------------------------------------------
\subsection{Minimum/Maximum}
\label{sss:mini-36e1}
%-------------------------------------------------------------------------------

If $4AB - E^2 < 0$, the function has no maximum or minimum; its graph forms a hyperbolic paraboloid.

If $4AB - E^2 > 0$, the function has a minimum if $A > 0$, and a maximum if $A < 0$; its graph forms an elliptic paraboloid. In this case, the minimum or maximum occurs at $(x_m, y_m)$ where:
\[
x\_m = -\frac{2BC - DE}{4AB - E^2},
\]

\[
y\_m = -\frac{2AD - CE}{4AB - E^2}.
\]
If $4AB - E^2 = 0$ and $DE - 2CB = 2AD - CE \neq 0$, the function has no maximum or minimum; its graph forms a parabolic cylinder.

If $4AB - E^2 = 0$ and $DE - 2CB = 2AD - CE = 0$, the function achieves the maximum/minimum at a line: a minimum if $A > 0$ and a maximum if $A < 0$; its graph forms a parabolic cylinder.

%-------------------------------------------------------------------------------
\subsection{Resources}
\label{sss:reso-55b5}
%-------------------------------------------------------------------------------

\begin{itemize}
    \item \href{https://mathresearch.utsa.edu/wikiFiles/MAT1053/Quadratic_Functions/MAT1053_M2.2Quadratic_Functions.pdf}{Quadratic Functions, Book Chapters}
    \item \href{https://mathresearch.utsa.edu/wikiFiles/MAT1053/Quadratic_Functions/MAT1053_M2.2Quadratic_FunctionsGN.pdf}{Guided Notes}
\end{itemize}

%-------------------------------------------------------------------------------
\subsection{Licensing}
\label{sss:lice-7352}
%-------------------------------------------------------------------------------

Content obtained and/or adapted from:

\begin{itemize}
    \item \href{https://en.wikipedia.org/wiki/Quadratic_function}{Quadratic function, Wikipedia under a CC BY-SA license}
\end{itemize}

%-------------------------------------------------------------------------------
\subsection{Learning Objectives}
\label{ssec:lear-6e2c}
%-------------------------------------------------------------------------------

\begin{enumerate}
    \item \textbf{Identify and Interpret the Vertex Form of a Quadratic Function} Students will be able to identify and convert between the standard form, factored form, and vertex form of a quadratic function. They will be able to:
    \item Recognize and use the vertex form of a quadratic function, $f(x) = a(x - h)^2 + k$, to determine the vertex $(h, k)$ of the parabola.
    \item Convert between the standard form $f(x) = ax^2 + bx + c$ and vertex form by completing the square.
    \item Understand how the vertex form reveals the minimum or maximum point of the quadratic function depending on the sign of $a$.
    \item \textbf{Calculate and Interpret the Roots of a Quadratic Function} Students will be able to:
    \item Find the roots of a quadratic function using the quadratic formula: $ r\_1 = \frac{-b - \sqrt{b^2 - 4ac}}{2a}, \quad r\_2 = \frac{-b + \sqrt{b^2 - 4ac}}{2a} $
    \item Interpret the nature of the roots based on the discriminant $b^2 - 4ac$, identifying real, repeated, or complex roots.
    \item Graph the quadratic function $f(x) = ax^2 + bx + c$ and accurately plot the roots and the y-intercept.
    \item \textbf{Apply Quadratic Functions to Model Real-World Situations} Students will be able to:
    \item Use quadratic functions to model real-world scenarios such as projectile motion, optimization problems, or financial calculations.
    \item Analyze the quadratic function to determine key features such as the vertex, axis of symmetry, and roots in the context of the application.
    \item Solve practical problems by setting up and solving quadratic equations, and interpret the solutions in terms of the original context.
\end{enumerate}

%-------------------------------------------------------------------------------
\subsection{Assessment Instrument}
\label{ssec:asse-6423}
%-------------------------------------------------------------------------------

Assess through problems requiring students to identify the vertex, axis of symmetry, and
intercepts of a quadratic function from standard, vertex, and intercept forms, and to
convert between forms. Include at least one applied optimization problem involving a
maximum or minimum value.

\textbf{Example problems.}

\begin{enumerate}
    \item Rewrite $f(x) = 2x^2 - 8x + 3$ in vertex form by completing the square, then identify the vertex and axis of symmetry.
    \item Find the $x$-intercepts, $y$-intercept, vertex, and axis of symmetry of $g(x) = -x^2 + 4x + 5$ and sketch the parabola.
    \item The height of a projectile is $h(t) = -4.9t^2 + 20t + 1.5$ metres. Find the maximum height and the time at which it occurs.
\end{enumerate}

%===============================================================================
\section{Dividing Polynomials}
\label{sec:divi-c574}
%===============================================================================

%-------------------------------------------------------------------------------
\subsection{Polynomial Long Division}
\label{sss:poly-7d92}
%-------------------------------------------------------------------------------

Suppose we would like to divide one polynomial by another. The procedure is similar to
long division of numbers and is illustrated in the following example:

%-------------------------------------------------------------------------------
\subsection{Example 1}
\label{sss:exam-acc3}
%-------------------------------------------------------------------------------

Divide $x^2 - 2x - 15$ (the dividend or numerator) by $(x + 3)$ (the divisor or denominator). 
Similar to long division of numbers, we set up our problem as follows:

\[
\begin{array}{r|l}
x+3 & x^2 - 2x - 15 \\
\end{array}
\]

First, we ask: how many times does $x+3$ go into $x^2$? 
Divide the leading term of the dividend by the leading term of the divisor:
\[
\frac{x^2}{x} = x.
\]

So it goes in $x$ times:

\[
\begin{array}{r|l}
x+3 & x^2 - 2x - 15 \\
    & x \\
\end{array}
\]

Next, multiply $x(x+3) = x^2 + 3x$ and subtract:

\[
\begin{array}{r|l}
x+3 & x^2 - 2x - 15 \\
    & x \\
    & -(x^2 + 3x) \\
\end{array}
\]

Now perform the subtraction:

\[
(x^2 - 2x) - (x^2 + 3x) = -5x.
\]

Bring down the $-15$:

\[
\begin{array}{r|l}
x+3 & x^2 - 2x - 15 \\
    & x \\
    & -(x^2 + 3x) \\
    & -5x - 15 \\
\end{array}
\]

Repeat the process:

\[
\frac{-5x}{x} = -5.
\]

Multiply and subtract:

\[
\begin{array}{r|l}
x+3 & x^2 - 2x - 15 \\
    & x - 5 \\
    & -(x^2 + 3x) \\
    & -5x - 15 \\
    & -(-5x - 15) \\
    & 0 \\
\end{array}
\]

Thus, the quotient is $x - 5$.In this case, we have no remainder.

%-------------------------------------------------------------------------------
\subsection{Example 2}
\label{sss:exam-5bdd}
%-------------------------------------------------------------------------------

What about non-divisible polynomials like this?

\[
(3x^2 + 3x - 4) / (x - 4)
\]

1. We first consider the highest-degree terms from both the dividend and divisor, the
result is the first term of our quotient.

\[
\frac{3x^2}{x} = 3x
\]

\[
\begin{array}{r|ccc}
   x - 4 & 3x^2 & 3x & -4 \\   
\end{array}
\]

2. Then we multiply this by our divisor.

   \[
(3x) \times (x - 4) = 3x^2 - 12x
\]

3. And subtract the result from our dividend.

\[
(3x^2 + 3x - 4) - (3x^2 - 12x) = 15x - 4
\]

\[
\begin{array}{r|ccc}
x - 4 & 3x^2 & 3x & -4 \\\   3x(x - 4) & 3x^2 & -12x & \\   \end{array}
\]

4. Now once again with the highest-degree terms of the remaining polynomial, and we got
the second term of our quotient.

\[
\frac{15x}{x} = 15
\]

5. Multiplying...

\[
(15) \times (x - 4) = 15x - 60
\]

6. Subtracting...

\[
(15x - 4) - (15x - 60) = 56
\]

\[
\begin{array}{r|ccc}
x - 4 & 3x^2 & 3x & -4 \\\   3x(x - 4) & 3x^2 & -12x & \\\   15(x - 4) & 15x & -60 \\   \end{array}
\]

7. We are left with a constant term - our remainder:

\[
Q(x) = 3x + 15, \quad R = 56
\]

So finally:

\[
(3x^2 + 3x - 4) = (3x + 15) \times (x - 4) + 56
\]

%-------------------------------------------------------------------------------
\subsection{Resources}
\label{sss:reso-55b5}
%-------------------------------------------------------------------------------

%-------------------------------------------------------------------------------
\subsection{Dividing Polynomials With Long Division}
\label{sss:divi-5583}
%-------------------------------------------------------------------------------

\begin{itemize}
    \item \href{https://www.khanacademy.org/math/algebra2/x2ec2f6f830c9fb89:poly-div/x2ec2f6f830c9fb89:quad-div-by-linear/v/polynomial-division}{Intro to Polynomial Long Division}, Khan Academy
    \item \href{https://courses.lumenlearning.com/waymakercollegealgebra/chapter/polynomial-long-division/}{Polynomial Long Division}, Lumen Learning
    \item \href{https://www.purplemath.com/modules/polydiv2.htm}{Polynomial Long Division}, Purple Math
    \item \href{https://www.youtube.com/watch?v=_FSXJmESFmQ}{Long Division With Polynomials}, The Organic Chemistry Tutor
    \item \href{https://www.youtube.com/watch?v=l6_ghhd7kwQ}{Long Division of Polynomials}, patrickJMT
\end{itemize}

%-------------------------------------------------------------------------------
\subsection{Dividing Polynomials with Synthetic Division}
\label{sss:divi-32bb}
%-------------------------------------------------------------------------------

\begin{itemize}
    \item \href{https://www.khanacademy.org/math/algebra-home/alg-polynomials/alg-synthetic-division-of-polynomials/v/synthetic-division}{Synthetic Division of Polynomials}, Khan Academy
    \item \href{https://www.purplemath.com/modules/synthdiv.htm}{Synthetic Division}, Purple Math
    \item \href{https://www.wtamu.edu/academic/anns/mps/math/mathlab/col_algebra/col_alg_tut37_syndiv.htm}{Synthetic Division}, WTAMU VirtualMathLab
    \item \href{https://www.youtube.com/watch?v=bZoMz1Cy1T4}{Synthetic Division Example}, patrickJMT
\end{itemize}

%-------------------------------------------------------------------------------
\subsection{Licensing}
\label{sss:lice-7352}
%-------------------------------------------------------------------------------

Content obtained and/or adapted from:

\begin{itemize}
    \item \href{https://en.wikibooks.org/wiki/Calculus/Algebra}{Algebra, Wikibooks: Calculus} under a CC BY-SA license
    \item \href{https://en.wikibooks.org/wiki/High_School_Mathematics_Extensions/Supplementary/Polynomial_Division}{Polynomial Division, Wikibooks: High School Mathematics Extensions} under a CC BY-SA license
\end{itemize}

%-------------------------------------------------------------------------------
\subsection{Learning Objectives}
\label{ssec:lear-ae78}
%-------------------------------------------------------------------------------

\begin{enumerate}
    \item SLO 1: Understand the Concept of Polynomial Long Division
    \item Students will be able to explain the steps involved in polynomial long division, including how to set up the division problem and identify the divisor and dividend. SLO 2: Perform Polynomial Long Division
    \item Students will be able to divide a given polynomial by another polynomial using the long division method, accurately performing each step to find the quotient and remainder. SLO 3: Apply Polynomial Long Division to Various Problems
    \item Students will be able to apply the polynomial long division technique to solve problems involving both divisible and non-divisible polynomials, demonstrating their ability to handle different scenarios.
\end{enumerate}

%-------------------------------------------------------------------------------
\subsection{Assessment Instrument}
\label{ssec:asse-45f7}
%-------------------------------------------------------------------------------

Assess through problems requiring students to divide polynomials using both long division
and synthetic division, expressing results in the form $p(x) = d(x)q(x) + r(x)$. At least
one problem should require students to verify the result using the Remainder Theorem.

\textbf{Example problems.}

\begin{enumerate}
    \item Divide $3x^3 - 2x^2 + x - 5$ by $x - 2$ using synthetic division and verify the remainder using the Remainder Theorem.
    \item Perform polynomial long division: $\dfrac{2x^4 - x^3 + 3x - 7}{x^2 + 1}$.
    \item Given that $(x + 3)$ is a factor of $x^3 + 4x^2 - x - 12$, use synthetic division to find the remaining factor and list all zeros.
\end{enumerate}

%===============================================================================
\section{Zeros of Polynomials}
\label{sec:zero-ebde}
%===============================================================================

Zeros of Polynomials

%-------------------------------------------------------------------------------
\subsection{Assessment Instrument}
\label{ssec:asse-7f8e}
%-------------------------------------------------------------------------------

Assess through problems requiring students to apply the Rational Zero Theorem to list
candidate rational zeros, use synthetic division to test them, and factor the polynomial
completely to find all real and complex zeros.

\textbf{Example problems.}

\begin{enumerate}
    \item List all possible rational zeros of $f(x) = 2x^3 - 3x^2 - 11x + 6$ using the Rational Zero Theorem, then find all actual zeros.
    \item Find all zeros of $g(x) = x^4 - 5x^2 + 4$, including any complex zeros.
    \item Use the Factor Theorem and synthetic division to factor $h(x) = x^3 + 2x^2 - 5x - 6$ completely over $\mathbb{R}$.
\end{enumerate}

%===============================================================================
\section{Power and Polynomial Functions}
\label{sec:powe-3059}
%===============================================================================

Power and Polynomial Functions

%-------------------------------------------------------------------------------
\subsection{Assessment Instrument}
\label{ssec:asse-85d1}
%-------------------------------------------------------------------------------

Assess through problems requiring students to identify end behavior using the leading term
test, classify polynomials by degree and number of terms, and describe the turning points
and zeros of a polynomial from its factored form.

\textbf{Example problems.}

\begin{enumerate}
    \item Describe the end behavior of $p(x) = -3x^5 + 2x^3 - x + 7$ and justify using the leading term test.
    \item Given $f(x) = x^3(x - 2)^2(x + 1)$, identify all zeros, their multiplicities, and the maximum number of turning points.
    \item Determine whether $f(x) = x^4 - 16$ is a power function. If so, identify the coefficient and the power; if not, explain why.
\end{enumerate}

%===============================================================================
\section{Graphs of Polynomials}
\label{sec:grap-918e}
%===============================================================================

%-------------------------------------------------------------------------------
\subsection{Polynomial Functions}
\label{sss:poly-bcf7}
%-------------------------------------------------------------------------------

%-------------------------------------------------------------------------------
\subsection{Polynomial of Degree 0:}
\label{sss:poly-2a4f}
%-------------------------------------------------------------------------------

$
f(x) = 2
$
% [Image: graph]

%-------------------------------------------------------------------------------
\subsection{Polynomial of Degree 1:}
\label{sss:poly-39b4}
%-------------------------------------------------------------------------------

$
f(x) = 2x + 1
$
% [Image: graph]

%-------------------------------------------------------------------------------
\subsection{Polynomial of Degree 2:}
\label{sss:poly-bf2f}
%-------------------------------------------------------------------------------

$
f(x) = x^2 - x - 2
$

$
= (x + 1)(x - 2)
$
% [Image: graph]

%-------------------------------------------------------------------------------
\subsection{Polynomial of Degree 3:}
\label{sss:poly-bd32}
%-------------------------------------------------------------------------------

$
f(x) = \frac{1}{4}x^3 + \frac{3}{4}x^2 - \frac{3}{2}x - 2
$

$
= \frac{1}{4} (x + 4)(x + 1)(x - 2)
$
% [Image: graph]

%-------------------------------------------------------------------------------
\subsection{Polynomial of Degree 4:}
\label{sss:poly-93ac}
%-------------------------------------------------------------------------------

$
f(x) = \frac{1}{14} (x + 4)(x + 1)(x - 1)(x - 3) + 0.5
$
% [Image: graph]

%-------------------------------------------------------------------------------
\subsection{Polynomial of Degree 5:}
\label{sss:poly-a79e}
%-------------------------------------------------------------------------------

$
f(x) = \frac{1}{20} (x + 4)(x + 2)(x + 1)(x - 1)(x - 3) + 2
$
% [Image: graph]

%-------------------------------------------------------------------------------
\subsection{Polynomial of Degree 6:}
\label{sss:poly-2daf}
%-------------------------------------------------------------------------------

$
f(x) = \frac{1}{100} (x^6 - 2x^5 - 26x^4 + 28x^3 + 145x^2 - 26x - 80)
$
% [Image: graph]

%-------------------------------------------------------------------------------
\subsection{Polynomial of Degree 7:}
\label{sss:poly-6100}
%-------------------------------------------------------------------------------

$
f(x) = (x - 3)(x - 2)(x - 1)x(x + 1)(x + 2)(x + 3)
$
% [Image: graph]

%-------------------------------------------------------------------------------
\subsection{Polynomial Functions Definition}
\label{sss:poly-1269}
%-------------------------------------------------------------------------------

A polynomial function is a function that can be defined by evaluating a polynomial. More precisely, a function $f$ of one argument from a given domain is a polynomial function if there exists a polynomial:
\[
a\_n x^n + a\_{n-1} x^{n-1} + \cdots + a\_2 x^2 + a\_1 x + a\_0
\]
that evaluates to $f(x)$ for all $x$ in the domain of $f$ (here, $n$ is a non-negative integer and $a_0, a_1, a_2, \ldots, a_n$ are constant coefficients). Generally, unless otherwise specified, polynomial functions have complex coefficients, arguments, and values. In particular, a polynomial, restricted to have real coefficients, defines a function from the complex numbers to the complex numbers. If the domain of this function is also restricted to the reals, the resulting function is a real function that maps reals to reals.

For example, the function $f$, defined by:
\[
f(x) = x^3 - x
\]
is a polynomial function of one variable. Polynomial functions of several variables are
similarly defined, using polynomials in more than one indeterminate, as in:
\[
f(x, y) = 2x^3 + 4x^2 y + xy^5 + y^2 - 7
\]
According to the definition of polynomial functions, there may be expressions that
obviously are not polynomials but nevertheless define polynomial functions. An example is
the expression:
\[
\left(\sqrt{1 - x^2}\right)^2
\]
which takes the same values as the polynomial:
\[
1 - x^2
\]
on the interval:
\[
[-1, 1]
\]
and thus both expressions define the same polynomial function on this interval.

Every polynomial function is continuous, smooth, and entire.

%-------------------------------------------------------------------------------
\subsection{Graphs}
\label{sss:grap-e870}
%-------------------------------------------------------------------------------

A polynomial function in one real variable can be represented by a graph.

\begin{itemize}
    \item The graph of the zero polynomial $f(x) = 0$ is the x-axis.
    \item The graph of a degree 0 polynomial $f(x) = a_0$, where $a_0 \neq 0$, is a horizontal line with y-intercept $a_0$.
    \item The graph of a degree 1 polynomial (or linear function) $f(x) = a_0 + a_1 x$, where $a_1 \neq 0$, is an oblique line with y-intercept $a_0$ and slope $a_1$.
    \item The graph of a degree 2 polynomial $f(x) = a_0 + a_1 x + a_2 x^2$, where $a_2 \neq 0$, is a parabola.
    \item The graph of a degree 3 polynomial $f(x) = a_0 + a_1 x + a_2 x^2 + a_3 x^3$, where $a_3 \neq 0$, is a cubic curve.
    \item The graph of any polynomial with degree 2 or greater $f(x) = a_0 + a_1 x + a_2 x^2 + \cdots + a_n x^n$, where $a_n \neq 0$ and $n \geq 2$, is a continuous non-linear curve.
\end{itemize}

A non-constant polynomial function tends to infinity when the variable increases indefinitely (in absolute value). If the degree is higher than one, the graph does not have any asymptote. It has two parabolic branches with vertical direction (one branch for positive $x$ and one for negative $x$).

Polynomial graphs are analyzed in calculus using intercepts, slopes, concavity, and end
behavior.

%-------------------------------------------------------------------------------
\subsection{Resources}
\label{sss:reso-55b5}
%-------------------------------------------------------------------------------

\begin{itemize}
    \item \href{https://tutorial.math.lamar.edu/classes/alg/graphingpolynomials.aspx}{Graphing Polynomials, Paul''s Online Notes from Lamar University}
    \item \href{https://courses.lumenlearning.com/collegealgebra2017/chapter/graphs-of-polynomial-functions/}{Graphs of Polynomial Functions, Lumen Learning College Algebra}
    \item \href{https://en.wikipedia.org/wiki/Polynomial}{Polynomial, Wikipedia}
\end{itemize}

%-------------------------------------------------------------------------------
\subsection{Licensing}
\label{sss:lice-7352}
%-------------------------------------------------------------------------------

Content obtained and/or adapted from:

\begin{itemize}
    \item \href{https://en.wikipedia.org/wiki/Polynomial}{Polynomial, Wikipedia under a CC BY-SA license}
\end{itemize}

%-------------------------------------------------------------------------------
\subsection{Learning Objectives}
\label{ssec:lear-8fbd}
%-------------------------------------------------------------------------------

\begin{enumerate}
    \item SRO 1: Identify Polynomial Degrees \textbf{Identify the degree of the polynomial function given its equation and graph.} SRO 2: Analyze Polynomial Graphs \textbf{Analyze the graph of a polynomial function and determine its degree and general shape.} SRO 3: Polynomial Functions Definition \textbf{Define and identify polynomial functions based on their algebraic expressions and graph representations.}
\end{enumerate}

%-------------------------------------------------------------------------------
\subsection{Assessment Instrument}
\label{ssec:asse-9c9d}
%-------------------------------------------------------------------------------

Assess through problems that require students to sketch the graph of a polynomial by
analyzing zeros, multiplicities, end behavior, and the $y$-intercept. At least one problem
should require students to write a possible polynomial equation from a graph description.

\textbf{Example problems.}

\begin{enumerate}
    \item Sketch $f(x) = (x + 2)^2(x - 1)(x - 3)$, labeling all zeros, their multiplicities, the $y$-intercept, and the end behavior.
    \item Describe how the multiplicity of a zero determines whether the graph crosses or touches the $x$-axis at that zero.
    \item Write a polynomial function (up to a constant multiple) whose graph has zeros at $x = -1$ (multiplicity~2), $x = 0$ (multiplicity~1), and $x = 4$ (multiplicity~1).
\end{enumerate}

%===============================================================================
\section{Rational Expressions}
\label{sec:rati-e460}
%===============================================================================

Rational expressions are fractions that have a polynomial in the numerator, denominator,
or both. Although rational expressions can seem complicated because they contain
variables, they can be simplified using techniques similar to those used for simplifying
expressions such as
\[
\frac{4x^3}{12x^2}
\]
combined with techniques for factoring polynomials.

%-------------------------------------------------------------------------------
\subsection{Simplifying Rational Expressions}
\label{sss:simp-4e69}
%-------------------------------------------------------------------------------

To simplify a rational expression, follow these steps:

1. \textbf{Determine the domain.} The excluded values are those values for the variable that result in the expression having a denominator of $0$.

2. \textbf{Factor the numerator and denominator.}

3. \textbf{Find common factors for the numerator and denominator and simplify.}

Consider two polynomials
\[
p(x) = a\_n x^n + a\_{n-1} x^{n-1} + \cdots + a\_1 x + a\_0
\]
and
\[
q(x) = b\_m x^m + b\_{m-1} x^{m-1} + \cdots + b\_1 x + b\_0
\]
When we take the quotient of the two we obtain
\[
\frac{p(x)}{q(x)} = \frac{a\_n x^n + a\_{n-1} x^{n-1} + \cdots + a\_1 x + a\_0}{b\_m x^m + b\_{m-1} x^{m-1} + \cdots + b\_1 x + b\_0}
\]
The ratio of two polynomials is called a rational expression. Many times we would like to
simplify such a beast. For example, say we are given
\[
\frac{x^2 - 1}{x + 1}
\]
We may simplify this in the following way:
\[
\frac{x^2 - 1}{x + 1} = \frac{(x + 1)(x - 1)}{x + 1} = x - 1, \quad x \neq -1
\]
This is nice because we have obtained something we understand quite well, $x - 1$, from something we didn''t.

%-------------------------------------------------------------------------------
\subsection{Resources}
\label{sss:reso-55b5}
%-------------------------------------------------------------------------------

\begin{itemize}
    \item \href{https://courses.lumenlearning.com/beginalgebra/chapter/6-1-1-introduction-to-rational-expressions/}{Identify and Simplify Rational Expressions, Lumen Learning}
    \item \href{https://www.youtube.com/watch?v=uVpsz-xpnPo}{Simplifying Rational Expressions, The Organic Chemistry Tutor}
\end{itemize}

%-------------------------------------------------------------------------------
\subsection{Licensing}
\label{sss:lice-7352}
%-------------------------------------------------------------------------------

Content obtained and/or adapted from:

\begin{itemize}
    \item \href{https://en.wikibooks.org/wiki/Calculus/Algebra}{Algebra, Wikibooks: Calculus} under a CC BY-SA license
\end{itemize}

%-------------------------------------------------------------------------------
\subsection{Learning Objectives}
\label{ssec:lear-6bf9}
%-------------------------------------------------------------------------------

\begin{enumerate}
    \item \textbf{Factor Rational Expressions:} Students will be able to factor polynomials in the numerator and denominator of a rational expression to simplify it.
    \item \textbf{Determine Domain:} Students will determine the domain of a rational expression by identifying values that make the denominator zero.
    \item \textbf{Apply Simplification Techniques:} Students will simplify rational expressions by canceling common factors and explaining the process used.
\end{enumerate}

%-------------------------------------------------------------------------------
\subsection{Assessment Instrument}
\label{ssec:asse-31c0}
%-------------------------------------------------------------------------------

Assess through problems requiring students to simplify, multiply, divide, add, and
subtract rational expressions, and to identify restrictions on the variable. Include at
least one problem requiring simplification of a complex fraction.

\textbf{Example problems.}

\begin{enumerate}
    \item Simplify $\dfrac{x^2 - 9}{x^2 - x - 6}$, stating all values of $x$ for which the expression is undefined.
    \item Perform the operation and simplify: $\dfrac{2}{x + 1} + \dfrac{3}{x - 2}$.
    \item Simplify the complex fraction $\dfrac{\dfrac{1}{x} - \dfrac{1}{y}}{\dfrac{1}{x} + \dfrac{1}{y}}$.
\end{enumerate}

%===============================================================================
\section{Graphs of Rational Functions}
\label{sec:grap-2359}
%===============================================================================

In mathematics, a rational function is any function that can be defined by a rational fraction, which is an algebraic fraction such that both the numerator and the denominator are polynomials. The coefficients of the polynomials need not be rational numbers; they may be taken in any field $K$. In this case, one speaks of a rational function and a rational fraction over $K$. The values of the variables may be taken in any field $L$ containing $K$. Then the domain of the function is the set of the values of the variables for which the denominator is not zero, and the codomain is $L$.

The set of rational functions over a field $K$ is a field, the field of fractions of the ring of the polynomial functions over $K$.

%-------------------------------------------------------------------------------
\subsection{Definitions}
\label{sss:defi-9fdc}
%-------------------------------------------------------------------------------

A function $f(x)$ is called a rational function if and only if it can be written in the form
\[
f(x) = \frac{P(x)}{Q(x)}
\]
where $P$ and $Q$ are polynomial functions of $x$ and $Q$ is not the zero function. The domain of $f$ is the set of all values of $x$ for which the denominator $Q(x)$ is not zero.

However, if $P$ and $Q$ have a non-constant polynomial greatest common divisor $R$, then setting $P = P_1 R$ and $Q = Q_1 R$ produces a rational function
\[
f\_1(x) = \frac{P\_1(x)}{Q\_1(x)},
\]
which may have a larger domain than $f(x)$ and is equal to $f(x)$ on the domain of $f(x)$. It is common usage to identify $f(x)$ and $f_1(x)$, that is to extend ``by continuity'' the domain of $f(x)$ to that of $f_1(x)$. Indeed, one can define a rational fraction as an equivalence class of fractions of polynomials, where two fractions $\frac{A(x)}{B(x)}$ and $\frac{C(x)}{D(x)}$ are considered equivalent if $A(x)D(x) = B(x)C(x)$. In this case, $\frac{P(x)}{Q(x)}$ is equivalent to $\frac{P_1(x)}{Q_1(x)}$.

A proper rational function is a rational function in which the degree of $P(x)$ is no greater than the degree of $Q(x)$ and both are real polynomials.

%-------------------------------------------------------------------------------
\subsection{Degree}
\label{sss:degr-fb81}
%-------------------------------------------------------------------------------

There are several non-equivalent definitions of the degree of a rational function.

Most commonly, the degree of a rational function is the maximum of the degrees of its constituent polynomials $P$ and $Q$, when the fraction is reduced to lowest terms. If the degree of $f$ is $d$, then the equation
\[
f(z) = w
\]
has $d$ distinct solutions in $z$ except for certain values of $w$, called critical values, where two or more solutions coincide or where some solution is rejected at infinity (that is, when the degree of the equation decreases after having cleared the denominator).

In the case of complex coefficients, a rational function with degree one is a Möbius
transformation.

The degree of the graph of a rational function is not the degree as defined above: it is
the maximum of the degree of the numerator and one plus the degree of the denominator.

In some contexts, such as in asymptotic analysis, the degree of a rational function is the
difference between the degrees of the numerator and the denominator.

In network synthesis and network analysis, a rational function of degree two (that is, the
ratio of two polynomials of degree at most two) is often called a biquadratic function.

%-------------------------------------------------------------------------------
\subsection{Examples}
\label{sss:exam-bfeb}
%-------------------------------------------------------------------------------

Rational function of degree 3, with a graph of degree 3:
\[
y = \frac{x^3 - 2x}{2(x^2 - 5)}
\]
% [Image: deg3]

Rational function of degree 2, with a graph of degree 3:
\[
y = \frac{x^2 - 3x - 2}{x^2 - 4}
\]
% [Image: deg2]

The rational function
\[
f(x) = \frac{x^3 - 2x}{2(x^2 - 5)}
\]
is not defined at
\[
x^2 = 5 \Rightarrow x = \pm \sqrt{5}.
\]
It is asymptotic to $\frac{x}{2}$ as $x \to \infty$.

The rational function
\[
f(x) = \frac{x^2 + 2}{x^2 + 1}
\]
is defined for all real numbers, but not for all complex numbers, since if $x$ were a square root of $-1$ (i.e., the imaginary unit or its negative), then formal evaluation would lead to division by zero:
\[
f(i) = \frac{i^2 + 2}{i^2 + 1} = \frac{-1 + 2}{-1 + 1} = \frac{1}{0},
\]
which is undefined.

A constant function such as $f(x) = \pi$ is a rational function since constants are polynomials. The function itself is rational, even though the value of $f(x)$ is irrational for all $x$.

Every polynomial function $f(x) = P(x)$ is a rational function with $Q(x) = 1$. A function that cannot be written in this form, such as $f(x) = \sin(x)$, is not a rational function. However, the adjective ``irrational'' is not generally used for functions.

The rational function
\[
f(x) = \frac{x}{x}
\]
is equal to 1 for all $x$ except 0, where there is a removable singularity. The sum, product, or quotient (excepting division by the zero polynomial) of two rational functions is itself a rational function. However, the process of reduction to standard form may inadvertently result in the removal of such singularities unless care is taken. Using the definition of rational functions as equivalence classes gets around this, since $x/x$ is equivalent to $1/1$.

%-------------------------------------------------------------------------------
\subsection{Sketch a Graph of Rational Functions}
\label{sss:sket-9dfc}
%-------------------------------------------------------------------------------

1. Evaluate the function at $0$ to find the $y$-intercept.

2. Factor the numerator and denominator.

3. For factors in the numerator not common to the denominator, determine where each factor of the numerator is zero to find the $x$-intercepts.

4. Find the multiplicities of the $x$-intercepts to determine the behavior of the graph at those points.

5. For factors in the denominator, note the multiplicities of the zeros to determine the
local behavior. For those factors not common to the numerator, find the vertical
asymptotes by setting those factors equal to zero and then solve.

6. For factors in the denominator common to factors in the numerator, find the removable discontinuities by setting those factors equal to $0$ and then solve.

7. Compare the degrees of the numerator and the denominator to determine the horizontal or
slant asymptotes.

8. Sketch the graph.

%-------------------------------------------------------------------------------
\subsection{References}
\label{sss:refe-44a2}
%-------------------------------------------------------------------------------

\begin{itemize}
    \item Martin J. Corless, Art Frazho, \textit{Linear Systems and Control}, p. 163, CRC Press, 2003. ISBN 0203911377.
    \item Malcolm W. Pownall, \textit{Functions and Graphs: Calculus Preparatory Mathematics}, p. 203, Prentice-Hall, 1983. ISBN 0133323048.
    \item Glisson, Tildon H., \textit{Introduction to Circuit Analysis and Design}, Springer, 2011. ISBN 9048194431.
    \item ``Rational function,'' Encyclopedia of Mathematics, EMS Press, 2001 [1994].
    \item Press, W.H.; Teukolsky, S.A.; Vetterling, W.T.; Flannery, B.P. (2007), ``Section 3.4. Rational Function Interpolation and Extrapolation,'' \textit{Numerical Recipes: The Art of Scientific Computing} (3rd ed.), New York: Cambridge University Press. ISBN 978-0-521-88068-8.
    \item \href{https://courses.lumenlearning.com/ivytech-collegealgebra/chapter/graph-rational-functions/}{Graph Rational Functions}, Lumen Learning.
\end{itemize}

%-------------------------------------------------------------------------------
\subsection{Licensing}
\label{sss:lice-7352}
%-------------------------------------------------------------------------------

Content obtained and/or adapted from:

\begin{itemize}
    \item \href{https://en.wikipedia.org/wiki/Rational_function}{"Rational function," Wikipedia} under a CC BY-SA license.
    \item \href{https://courses.lumenlearning.com/ivytech-collegealgebra/chapter/graph-rational-functions/}{Graph Rational Functions, Lumen Learning College Algebra} under a CC BY license.
\end{itemize}

%-------------------------------------------------------------------------------
\subsection{Learning Objectives}
\label{ssec:lear-eb17}
%-------------------------------------------------------------------------------

\begin{enumerate}
    \item \textbf{Identify the Domain and Range of Rational Functions}: Students will be able to determine the domain and range of rational functions by analyzing the numerator and denominator polynomials.
    \item \textbf{Analyze Asymptotic Behavior}: Students will be able to identify and analyze the vertical, horizontal, and slant asymptotes of rational functions by comparing the degrees of the numerator and denominator.
    \item \textbf{Graph Rational Functions}: Students will be able to sketch graphs of rational functions by finding intercepts, asymptotes, and removable discontinuities, and determining the behavior of the graph near these key points.
\end{enumerate}

%-------------------------------------------------------------------------------
\subsection{Assessment Instrument}
\label{ssec:asse-5038}
%-------------------------------------------------------------------------------

Assess through problems requiring students to identify vertical asymptotes, horizontal or
oblique asymptotes, holes, intercepts, and end behavior for a given rational function, and
to produce a labeled sketch of the graph.

\textbf{Example problems.}

\begin{enumerate}
    \item Find all vertical and horizontal asymptotes of $f(x) = \dfrac{2x^2 - 3x}{x^2 - 4}$ and sketch the graph.
    \item Determine whether $g(x) = \dfrac{x^2 - 1}{x - 1}$ has a hole or a vertical asymptote at $x = 1$. Justify and sketch the function.
    \item Find the oblique asymptote of $h(x) = \dfrac{x^2 + 2x - 3}{x - 1}$ using polynomial division and sketch the graph.
\end{enumerate}

%===============================================================================
\section{Single Transformations of Functions}
\label{sec:sing-9dfc}
%===============================================================================

%-------------------------------------------------------------------------------
\subsection{Introduction}
\label{sss:intr-8800}
%-------------------------------------------------------------------------------

%-------------------------------------------------------------------------------
\subsection{Vertical Shift}
\label{sss:vert-45b1}
%-------------------------------------------------------------------------------

$f(x) = x^2$ (red) and $g(x) = x^2 + 4$ (blue)

% [Image: graph]

%-------------------------------------------------------------------------------
\subsection{Horizontal Shift}
\label{sss:hori-8d28}
%-------------------------------------------------------------------------------

$f(x) = \sin(x)$ (red) and $g(x) = \sin(x + \frac{\pi}{2})$ (blue)

% \href{https://upload.wikimedia.org/wikipedia/commons/thumb/2/2b/Horizontal_shift.png/300px-Horizontal_shift.png}{Image: graph}

%-------------------------------------------------------------------------------
\subsection{Vertical Reflection}
\label{sss:vert-0225}
%-------------------------------------------------------------------------------

$f(x) = \sqrt{x}$ (red) and $g(x) = -\sqrt{x}$ (blue)

% [Image: graph]

%-------------------------------------------------------------------------------
\subsection{Horizontal Reflection}
\label{sss:hori-6e98}
%-------------------------------------------------------------------------------

$f(x) = e^x$ (red) and $g(x) = e^{-x}$ (blue)

% [Image: graph]

%-------------------------------------------------------------------------------
\subsection{Vertical Stretch/Compression}
\label{sss:vert-3d17}
%-------------------------------------------------------------------------------

$f(x) = x^2$ (red), $g(x) = \frac{1}{2}x^2$ (blue, vertical compression), $h(x) = 2x^2$ (green, vertical stretch), and $j(x) = -2x^2$ (black, vertical stretch and reflection)

% [Image: graph]

%-------------------------------------------------------------------------------
\subsection{Horizontal Stretch/Compression}
\label{sss:hori-0121}
%-------------------------------------------------------------------------------

$f(x) = \sqrt{x}$ (red), $g(x) = \sqrt{2x}$ (blue, horizontal compression), $h(x) = \sqrt{\frac{1}{2}x}$ (green, horizontal stretch), and $j(x) = \sqrt{-2x}$ (black, horizontal compression and reflection)

% [Image: graph]

%-------------------------------------------------------------------------------
\subsection{Translations}
\label{sss:tran-2ed1}
%-------------------------------------------------------------------------------

A vertical shift adds or subtracts a constant to the function. For example, $g(x) = x^2 + 4$ is $f(x) = x^2$ shifted up by 4 units, and $g(x) = \sin(x) - 7.7$ is $f(x) = \sin(x)$ shifted down by 7.7 units.

A horizontal shift modifies the function as $g(x) = f(x - h)$. For example, $g(x) = (x - 3)^2$ is $f(x) = x^2$ shifted 3 units to the right, and $g(x) = \sin(x + \frac{\pi}{2})$ is $f(x) = \sin(x)$ shifted $\frac{\pi}{2}$ units to the left.

%-------------------------------------------------------------------------------
\subsection{Reflections}
\label{sss:refl-afab}
%-------------------------------------------------------------------------------

A vertical reflection changes $f(x)$ to $g(x) = -f(x)$, flipping the graph over the x-axis. For example, $g(x) = -\sqrt{x}$ is a vertical reflection of $f(x) = \sqrt{x}$.

A horizontal reflection changes $f(x)$ to $g(x) = f(-x)$, flipping the graph over the y-axis. For example, $g(x) = e^{-x}$ is a horizontal reflection of $f(x) = e^x$.

%-------------------------------------------------------------------------------
\subsection{Even and Odd Functions}
\label{sss:even-a6f0}
%-------------------------------------------------------------------------------

A function $f$ is even if $f(x) = f(-x)$. For example, $f(x) = x^2$ is even. 

A function $f$ is odd if $f(x) = -f(-x)$. For example, $f(x) = x^3$ is odd.

If a function is neither even nor odd, it does not satisfy either condition. For example, $f(x) = x^2 + x$ is neither even nor odd.

%-------------------------------------------------------------------------------
\subsection{Compressions and Stretches}
\label{sss:comp-67a5}
%-------------------------------------------------------------------------------

A vertical stretch/compression is represented by $g(x) = af(x)$. If $a > 1$, the graph stretches. If $0 < a < 1$, it compresses. If $a < 0$, it stretches or compresses with a reflection.

A horizontal stretch/compression is represented by $g(x) = f(bx)$. If $b > 1$, the graph compresses horizontally. If $0 < b < 1$, it stretches. If $b < 0$, it stretches or compresses with a reflection.

%-------------------------------------------------------------------------------
\subsection{Resources}
\label{sss:reso-55b5}
%-------------------------------------------------------------------------------

\begin{itemize}
    \item \href{https://courses.lumenlearning.com/wmopen-collegealgebra/chapter/introduction-transformations-of-functions/}{Intro to Transformations of Functions, Lumen Learning}
\end{itemize}

%-------------------------------------------------------------------------------
\subsection{Licensing}
\label{sss:lice-7352}
%-------------------------------------------------------------------------------

Content obtained and/or adapted from:

\begin{itemize}
    \item \href{https://courses.lumenlearning.com/wmopen-collegealgebra/chapter/introduction-transformations-of-functions/}{Transformations of Functions, Lumen Learning College Algebra under a CC BY-SA license}
\end{itemize}

%-------------------------------------------------------------------------------
\subsection{Learning Objectives}
\label{ssec:lear-8667}
%-------------------------------------------------------------------------------

\begin{enumerate}
    \item \textbf{Vertical and Horizontal Shifts}: Given a function $f(x)$, identify and describe the effect of vertical and horizontal shifts on the function. For example, explain how $g(x) = f(x - h)$ shifts $f(x)$ horizontally.
    \item \textbf{Reflections}: Given a function $f(x)$, determine the result of vertical and horizontal reflections on the function. For example, describe how $g(x) = -f(x)$ reflects $f(x)$ over the x-axis.
    \item \textbf{Stretching and Compression}: Given a function $f(x)$, analyze the impact of vertical and horizontal stretches or compressions. For example, explain the effect of $g(x) = af(x)$ and $g(x) = f(bx)$ on the graph of $f(x)$.
\end{enumerate}

%-------------------------------------------------------------------------------
\subsection{Assessment Instrument}
\label{ssec:asse-37a8}
%-------------------------------------------------------------------------------

Assess through problems requiring students to write the equation of a transformed function
from a verbal description of a single transformation, and conversely to describe the
transformation applied to a parent function. All four transformation types must appear:
vertical and horizontal shifts, reflections, and stretches or compressions.

\textbf{Example problems.}

\begin{enumerate}
    \item Given $f(x) = x^2$, write the equation obtained by: (a)~shifting 3 units left; (b)~reflecting over the $x$-axis; (c)~vertically stretching by a factor of 4.
    \item Describe in words the single transformation that maps $f(x) = \sqrt{x}$ to $g(x) = \sqrt{-x}$, and state whether the domain changes.
    \item Given $h(x) = |2x|$, identify the parent function and describe the transformation applied. Sketch both graphs on the same axes.
\end{enumerate}

%===============================================================================
\section{Multiple Transformations of Functions}
\label{sec:mult-9f5c}
%===============================================================================
See \href{https://mathresearch.utsa.edu/alice/src/Pagina.aspx?pagina=25\&opus=1073\&lng=BRITANNIA}{Single Transformations of Functions} for more information on translating, reflecting, compressing, and stretching functions.

%-------------------------------------------------------------------------------
\subsection{Combining Functions}
\label{sss:comb-7f72}
%-------------------------------------------------------------------------------

\textbf{Vertical and horizontal shift:} $f(x) = x^2 + x$ (red) and $g(x) = (x - 3)^2 + (x - 3) + 5$ (blue; $f(x)$ shifted 3 units right and 5 units up).

If $f(x)$ is some function, then $g(x) = f(x - h) + k$ is the function $f(x)$ shifted $h$ units horizontally (to the right for $h > 0$ and to the left for $h < 0$) and $k$ units vertically (up for $k > 0$ and down for $k < 0$). For example, $g(x) = (x - 3)^2 + (x - 3) + 5$ is the function $f(x) = x^2 + x$ shifted 3 units to the right and 5 units up. $g(x) = \sqrt{x + 2} - 6$ is the function $f(x) = \sqrt{x}$ shifted 2 units to the left and 6 units down.

The order that two transformations are applied can change the resulting function. A
vertical transformation and a horizontal transformation can be applied in either order and
will result in the same function. However, two different vertical transformations or two
different horizontal transformations can result in differing functions depending on the
order they are applied in.

Let $f(x) = x^3 + x^2$ for the following examples:

\textbf{Vertical shift and horizontal reflection:} A 3 unit vertical shift, THEN a horizontal reflection of $f(x)$, will result in the function $g(x) = (-x)^3 + (-x)^2 + 3 = -x^3 + x^2 + 3$. If we apply the horizontal reflection first instead, we will still get the same result.

\textbf{Vertical shift and vertical reflection:} A 3 unit vertical shift, THEN a vertical reflection of $f(x)$, will result in the function $g(x) = -(x^3 + x^2 + 3) = -x^3 - x^2 - 3$. However, if we apply the vertical reflection before the vertical shift, we will get the function $h(x) = -x^3 - x^2 + 3$, which is not equal to our previous result $g(x)$.

\textbf{Horizontal shift and vertical compression:} A 3 unit horizontal shift, then a vertical compression of 2, will result in the function $g(x) = 2((x - 3)^3 + (x - 3)^2)$. If we apply the vertical compression first, we get the same result.

\textbf{Horizontal shift and horizontal compression:} A 3 unit horizontal shift, then a horizontal compression of 2, will result in the function $g(x) = (2x - 3)^3 + (2x - 3)^2$. However, if we apply the horizontal compression first, we get the function $g(x) = 2(x - 3)^3 + 2(x - 3)^2 = (2x - 6)^3 + (2x - 6)^2$, which is not the same as our previous result $g(x)$.

%-------------------------------------------------------------------------------
\subsection{Resources}
\label{sss:reso-55b5}
%-------------------------------------------------------------------------------

\begin{itemize}
    \item \href{https://online.math.uh.edu/Math1330-unpaid/ch1/s13/CombTransf/Combining_Transformations_Math1330_s13.pdf}{Combining Transformations, University of Houston}
    \item \href{https://courses.lumenlearning.com/waymakercollegealgebra/chapter/sequences-of-transformations/}{Sequences of Transformations, Lumen Learning}
\end{itemize}

%-------------------------------------------------------------------------------
\subsection{Learning Objectives}
\label{ssec:lear-74d6}
%-------------------------------------------------------------------------------

\begin{enumerate}
    \item \textbf{Identify and graph multiple transformations of functions:} Students will be able to identify vertical and horizontal shifts, reflections, compressions, and stretches of a given function, and graph the resulting function accurately.
    \item \textbf{Analyze the order of transformations:} Students will be able to explain how the order of applying two transformations, such as reflections and shifts, affects the resulting function, and determine when the order matters.
    \item \textbf{Apply combined transformations to functions:} Students will be able to perform combined transformations on functions, including shifts, reflections, and compressions, and write the equation of the transformed function.
\end{enumerate}

%-------------------------------------------------------------------------------
\subsection{Assessment Instrument}
\label{ssec:asse-e613}
%-------------------------------------------------------------------------------

Assess through problems requiring students to apply a specified sequence of
transformations in order, write the resulting function equation, and analyze how changing
the order of two transformations of the same type affects the outcome. Include a
comparison problem that requires applying two different orderings and comparing the
resulting equations.

\textbf{Example problems.}

\begin{enumerate}
    \item Starting from $f(x) = x^3$, apply the following in order: shift 2 units right, then reflect over the $y$-axis. Write the equation of the final function.
    \item Explain why applying a vertical shift before a vertical reflection yields a different result than applying the reflection first. Illustrate with $f(x) = x^2$.
    \item Apply a horizontal compression by a factor of 3 and then a vertical shift of 4 units up to $f(x) = \sqrt{x}$. Write the equation and graph the result.
\end{enumerate}